% ============================================================
%  Ordinary Differential Equations — Full Course
%  Preamble + Chapters 1–2
% ============================================================

\documentclass[12pt,a4paper,openany]{book}
\usepackage[utf8]{inputenc}
\usepackage[T1]{fontenc}
\usepackage{lmodern}
\usepackage[margin=2.5cm]{geometry}
\usepackage{amsmath,amssymb,amsthm,mathrsfs,mathtools}
\usepackage{tikz}
\usetikzlibrary{arrows.meta,calc,positioning,patterns,decorations.markings,decorations.pathmorphing,cd}
\usepackage{pgfplots}
\pgfplotsset{compat=1.18}
\usepackage[most]{tcolorbox}
\tcbuselibrary{theorems,skins,breakable}
\usepackage[colorlinks=true,linkcolor=blue!60!black,citecolor=green!50!black]{hyperref}
\usepackage{makeidx}
\makeindex
\usepackage{enumitem}
\usepackage{fancyhdr}
\usepackage{caption}
\usepackage{booktabs}
\usepackage{array}
\usepackage{microtype}
\usepackage{xcolor}

\pagestyle{fancy}
\fancyhf{}
\fancyhead[LE]{\slshape\leftmark}
\fancyhead[RO]{\slshape\rightmark}
\fancyfoot[C]{\thepage}

% Theorem environments — using \newtheorem (amsthm)
\theoremstyle{definition}
\newtheorem{definition}{Definition}[chapter]
\newtheorem{example}[definition]{Example}
\newtheorem{exercise}{Exercise}[chapter]

\theoremstyle{plain}
\newtheorem{theorem}[definition]{Theorem}
\newtheorem{proposition}[definition]{Proposition}
\newtheorem{lemma}[definition]{Lemma}
\newtheorem{corollary}[definition]{Corollary}

\theoremstyle{remark}
\newtheorem{remark}[definition]{Remark}
\newtheorem{notation}[definition]{Notation}
\newtheorem{method}[definition]{Method}

% tcolorbox wrapping (visual only — does NOT use \newtcbtheorem)
\tcolorboxenvironment{definition}{enhanced,breakable,colback=blue!5,colframe=blue!70!black,fonttitle=\bfseries,before skip=10pt,after skip=10pt,left=8pt,right=8pt}
\tcolorboxenvironment{theorem}{enhanced,breakable,colback=red!5,colframe=red!70!black,fonttitle=\bfseries,before skip=10pt,after skip=10pt,left=8pt,right=8pt}
\tcolorboxenvironment{proposition}{enhanced,breakable,colback=orange!5,colframe=orange!70!black,fonttitle=\bfseries,before skip=10pt,after skip=10pt,left=8pt,right=8pt}
\tcolorboxenvironment{lemma}{enhanced,breakable,colback=yellow!5,colframe=yellow!50!black,fonttitle=\bfseries,before skip=10pt,after skip=10pt,left=8pt,right=8pt}
\tcolorboxenvironment{corollary}{enhanced,breakable,colback=green!5,colframe=green!50!black,fonttitle=\bfseries,before skip=10pt,after skip=10pt,left=8pt,right=8pt}
\tcolorboxenvironment{remark}{enhanced,breakable,colback=gray!5,colframe=gray!50!black,before skip=8pt,after skip=8pt,left=6pt,right=6pt}
\tcolorboxenvironment{example}{enhanced,breakable,colback=purple!5,colframe=purple!50!black,before skip=8pt,after skip=8pt,left=6pt,right=6pt}

% Custom commands
\newcommand{\R}{\mathbb{R}}
\newcommand{\C}{\mathbb{C}}
\newcommand{\N}{\mathbb{N}}
\newcommand{\Z}{\mathbb{Z}}
\newcommand{\Q}{\mathbb{Q}}
\newcommand{\dd}{\mathrm{d}}
\newcommand{\abs}[1]{\left|#1\right|}
\newcommand{\norm}[1]{\left\|#1\right\|}
\DeclareMathOperator{\tr}{tr}
\DeclareMathOperator{\Sp}{Sp}
\DeclareMathOperator{\diag}{diag}

% ============================================================
\begin{document}

% ============================================================
%  TITLE PAGE
% ============================================================
\begin{titlepage}
\begin{tikzpicture}[remember picture, overlay]
  \fill[blue!10!white] (current page.south west) rectangle (current page.north east);
  \fill[blue!80!black, rotate around={-15:(current page.center)}]
    ([xshift=-4cm,yshift=-1cm]current page.center) rectangle ([xshift=4cm,yshift=1.5cm]current page.center);
  \node[anchor=center, text=white, font=\Huge\bfseries, rotate=-15]
    at ([yshift=0.3cm]current page.center) {Ordinary Differential Equations};
  \node[anchor=center, text=orange!80!yellow, font=\Large, rotate=-15]
    at ([yshift=-1cm]current page.center) {Lecture Notes};
  \node[anchor=north, text=blue!70!black, font=\Large\itshape]
    at ([yshift=-4cm]current page.center) {Ya\'e Ulrich Gaba};
  \node[anchor=north, text=blue!50!black, font=\large]
    at ([yshift=-5.5cm]current page.center) {Licence L3 --- 2025--2026};
  \fill[orange!70!yellow] ([yshift=3cm]current page.south west) rectangle ([yshift=3.3cm]current page.south east);
  \node[anchor=south, text=gray, font=\small] at ([yshift=1cm]current page.south) {\today};
\end{tikzpicture}
\end{titlepage}

% ============================================================
%  PREFACE
% ============================================================
\frontmatter

\chapter*{Preface}
\addcontentsline{toc}{chapter}{Preface}

Ordinary differential equations are among the most fundamental and widely applied tools in all of mathematics. They arise naturally whenever a quantity changes at a rate that depends on its current state: the swing of a pendulum, the growth of a population, the discharge of a capacitor, the trajectory of a planet. Understanding how to formulate, solve, and analyse such equations is essential for students of mathematics, physics, engineering, biology, and economics alike.

This textbook is intended for a one-semester or two-semester course at the advanced undergraduate level. The only prerequisites are a solid grounding in single-variable and multivariable calculus, and a first course in linear algebra. Some familiarity with basic real analysis is helpful but not strictly required; all essential analytic arguments are presented in full.

The book is structured as follows. Chapters~1 and~2 introduce the vocabulary of the subject and develop systematic techniques for solving first-order equations. Chapters~3 and~4 treat higher-order linear equations and the Laplace transform. Chapters~5 and~6 address systems of equations and qualitative theory. Later chapters cover power series methods, boundary value problems, and an introduction to partial differential equations.

Each chapter contains:
\begin{itemize}[leftmargin=2cm]
  \item Precise definitions and clearly stated theorems with complete proofs.
  \item Numerous worked examples illustrating every technique.
  \item Graphical illustrations produced with TikZ and \texttt{pgfplots}.
  \item Exercises graded by difficulty: one star for routine practice, two stars for intermediate problems, and three stars for challenging or theoretical questions.
  \item A chapter summary for efficient review.
\end{itemize}

We have striven to balance rigour with accessibility. Every major existence and uniqueness theorem is proved, but the emphasis throughout is on understanding \emph{why} the methods work, not merely on applying recipes.

\vspace{1cm}
\hfill\textit{The authors, 2026}

\tableofcontents

\mainmatter

% ============================================================
% ============================================================
%  CHAPTER 1: INTRODUCTION AND MODELLING
% ============================================================
% ============================================================
\chapter{Introduction and Modelling --- What Is an ODE?}
\label{ch:introduction}
\index{ordinary differential equation}

\section{Why Differential Equations?}
\label{sec:why-odes}

Many of the fundamental laws of nature are not expressed as direct relationships between quantities, but as relationships between quantities and their \emph{rates of change}. Newton's second law, for instance, does not tell us the position of a particle directly; it tells us how the acceleration (the second derivative of position) relates to the forces acting on the particle. To find the actual motion, we must solve a differential equation.

\index{Newton's second law}

In this introductory chapter we survey several real-world models that lead to ordinary differential equations, establish the basic terminology and classification of such equations, describe what it means to ``solve'' an ODE, and introduce the geometric viewpoint of direction fields. By the end of the chapter the reader should have a clear picture of the landscape we shall explore in the rest of the book.

% --------------------------------------------------------
\section{Real-World Models}
\label{sec:models}

\subsection{The Simple Pendulum}
\label{subsec:pendulum}
\index{pendulum}

Consider a point mass $m$ suspended from a rigid, massless rod of length $\ell$, free to swing in a vertical plane. Let $\theta(t)$ denote the angle the rod makes with the downward vertical at time $t$. Balancing the tangential component of gravity against the tangential acceleration gives
\begin{equation}\label{eq:pendulum}
  m\ell\,\ddot{\theta} = -mg\sin\theta,
\end{equation}
or equivalently,
\begin{equation}\label{eq:pendulum-standard}
  \ddot{\theta} + \frac{g}{\ell}\sin\theta = 0.
\end{equation}
This is a \emph{second-order, nonlinear}\index{nonlinear ODE} ordinary differential equation in the unknown function $\theta(t)$. For small angles one may approximate $\sin\theta\approx\theta$, yielding the linearised equation
\begin{equation}\label{eq:pendulum-linear}
  \ddot{\theta} + \frac{g}{\ell}\,\theta = 0,
\end{equation}
whose solutions are sinusoidal oscillations of period $T = 2\pi\sqrt{\ell/g}$.

\subsection{Population Growth}
\label{subsec:population}
\index{population growth}

Let $P(t)$ be the size of a population at time $t$. The simplest model assumes that the rate of growth is proportional to the current population:
\begin{equation}\label{eq:malthus}
  \frac{\dd P}{\dd t} = rP, \qquad P(0)=P_0,
\end{equation}
where $r>0$ is the intrinsic growth rate. This is the \emph{Malthusian} (exponential) growth model\index{Malthusian growth}, with solution $P(t)=P_0 e^{rt}$.

A more realistic model accounts for limited resources by introducing a carrying capacity $K$:
\begin{equation}\label{eq:logistic}
  \frac{\dd P}{\dd t} = rP\!\left(1-\frac{P}{K}\right).
\end{equation}
This is the \emph{logistic equation}\index{logistic equation}, a first-order, nonlinear ODE. Its solution is
\[
  P(t) = \frac{K}{1 + \left(\frac{K}{P_0}-1\right)e^{-rt}}.
\]

\subsection{Radioactive Decay}
\label{subsec:radioactive}
\index{radioactive decay}

A sample of a radioactive isotope decays at a rate proportional to the amount present:
\begin{equation}\label{eq:decay}
  \frac{\dd N}{\dd t} = -\lambda N, \qquad N(0) = N_0,
\end{equation}
where $\lambda > 0$ is the decay constant\index{decay constant}. The solution $N(t)=N_0 e^{-\lambda t}$ gives the well-known exponential decay law. The \emph{half-life}\index{half-life} is $t_{1/2} = \ln 2 / \lambda$.

\subsection{RC Circuits}
\label{subsec:rc-circuit}
\index{RC circuit}

Consider a resistor $R$ in series with a capacitor $C$ and an applied voltage $E(t)$. By Kirchhoff's voltage law\index{Kirchhoff's law}, the charge $q(t)$ on the capacitor satisfies
\begin{equation}\label{eq:rc-circuit}
  R\,\frac{\dd q}{\dd t} + \frac{q}{C} = E(t).
\end{equation}
This is a \emph{first-order, linear}\index{linear ODE} ODE. If $E(t)=E_0$ is constant and $q(0)=0$, the solution is
\[
  q(t) = CE_0\!\left(1 - e^{-t/(RC)}\right).
\]
The product $\tau = RC$ is the \emph{time constant}\index{time constant} of the circuit.

% --------------------------------------------------------
\section{Basic Definitions}
\label{sec:basic-definitions}

\begin{definition}[Ordinary Differential Equation]
\label{def:ode}
\index{ordinary differential equation!definition}
An \emph{ordinary differential equation} (ODE) is an equation involving an unknown function $y = y(t)$ of a single independent variable $t$ and one or more of its derivatives $y', y'', \ldots, y^{(n)}$. The general form is
\begin{equation}\label{eq:general-ode}
  F\!\left(t, y, y', y'', \ldots, y^{(n)}\right) = 0,
\end{equation}
where $F$ is a given function.
\end{definition}

\begin{definition}[Order of an ODE]
\label{def:order}
\index{order of an ODE}
The \emph{order} of an ODE is the order of the highest derivative that appears in the equation. Equation~\eqref{eq:general-ode} is of order~$n$ if $F$ depends non-trivially on $y^{(n)}$.
\end{definition}

\begin{example}
\label{ex:orders}
The equation $y' + 2y = e^t$ is first-order. The pendulum equation~\eqref{eq:pendulum-standard} is second-order. The equation $y''' - 3y'' + y = 0$ is third-order.
\end{example}

\begin{definition}[Degree of an ODE]
\label{def:degree}
\index{degree of an ODE}
When an ODE can be written as a polynomial in the highest-order derivative, the \emph{degree} of the ODE is the power of that highest derivative. For instance, $(y'')^3 + y' = t$ has order~$2$ and degree~$3$.
\end{definition}

\begin{definition}[Linear ODE]
\label{def:linear-ode}
\index{linear ODE!definition}
An ODE of order $n$ is \emph{linear} if it can be written in the form
\begin{equation}\label{eq:linear-ode}
  a_n(t)\,y^{(n)} + a_{n-1}(t)\,y^{(n-1)} + \cdots + a_1(t)\,y' + a_0(t)\,y = g(t),
\end{equation}
where the coefficients $a_0, a_1, \ldots, a_n$ and the right-hand side $g$ are given functions of $t$ alone. An ODE that is not linear is called \emph{nonlinear}\index{nonlinear ODE}.
\end{definition}

\begin{remark}
\label{rem:linearity-check}
Linearity requires that the unknown $y$ and all its derivatives appear to the first power and are not multiplied together. Thus $y'' + y^2 = 0$ is nonlinear (because of $y^2$), and $yy' = 1$ is nonlinear (because of the product $yy'$).
\end{remark}

\begin{definition}[Autonomous ODE]
\label{def:autonomous}
\index{autonomous ODE}
An ODE is \emph{autonomous} if the independent variable $t$ does not appear explicitly. In first-order form: $y' = f(y)$. The logistic equation~\eqref{eq:logistic} and the pendulum equation~\eqref{eq:pendulum-standard} are both autonomous.
\end{definition}

\begin{definition}[Homogeneous Linear ODE]
\label{def:homogeneous}
\index{homogeneous ODE}
A linear ODE~\eqref{eq:linear-ode} is \emph{homogeneous} if $g(t) \equiv 0$; otherwise it is \emph{non-homogeneous}\index{non-homogeneous ODE}.
\end{definition}

% --------------------------------------------------------
\section{Solution Concepts}
\label{sec:solution-concepts}

\begin{definition}[Solution of an ODE]
\label{def:solution}
\index{solution!of an ODE}
A \emph{solution} of the ODE $F(t,y,y',\ldots,y^{(n)})=0$ on an interval $I\subset\R$ is a function $\varphi\colon I\to\R$ that is $n$ times differentiable and satisfies
\[
  F\!\left(t,\varphi(t),\varphi'(t),\ldots,\varphi^{(n)}(t)\right) = 0
  \quad\text{for all } t\in I.
\]
\end{definition}

\begin{definition}[General Solution]
\label{def:general-solution}
\index{solution!general}
The \emph{general solution} of an $n$-th order ODE is a family of solutions containing $n$ arbitrary constants $c_1, c_2, \ldots, c_n$, such that (under appropriate conditions) every solution can be obtained by choosing particular values of these constants.
\end{definition}

\begin{definition}[Particular Solution]
\label{def:particular-solution}
\index{solution!particular}
A \emph{particular solution} is any single solution obtained from the general solution by assigning specific values to the arbitrary constants.
\end{definition}

\begin{definition}[Singular Solution]
\label{def:singular-solution}
\index{solution!singular}
A \emph{singular solution} is a solution that cannot be obtained from the general solution for any choice of the constants. Singular solutions typically arise as envelopes of the family of curves defined by the general solution.
\end{definition}

\begin{example}
\label{ex:singular}
The ODE $(y')^2 = 4y$ has the general solution $y = (t-c)^2$ for $c\in\R$, representing a family of parabolas. The function $y\equiv 0$ is also a solution, but it is not a member of this family for any finite $c$. It is a singular solution and is, geometrically, the envelope of the family of parabolas.
\end{example}

% --------------------------------------------------------
\section{Initial Value Problems}
\label{sec:ivp}

\begin{definition}[Initial Value Problem]
\label{def:ivp}
\index{initial value problem}
An \emph{initial value problem} (IVP) for an $n$-th order ODE consists of the equation together with $n$ initial conditions:
\begin{equation}\label{eq:ivp}
\begin{cases}
  y^{(n)} = f(t, y, y', \ldots, y^{(n-1)}), \\[4pt]
  y(t_0) = y_0,\quad y'(t_0)=y_1,\quad \ldots,\quad y^{(n-1)}(t_0) = y_{n-1}.
\end{cases}
\end{equation}
The point $(t_0, y_0, y_1, \ldots, y_{n-1})$ specifies the initial state of the system.
\end{definition}

\begin{example}
\label{ex:ivp-decay}
The radioactive decay problem $N' = -\lambda N$, $N(0) = N_0$ is an IVP. Given the initial amount $N_0$, there is exactly one solution: $N(t) = N_0 e^{-\lambda t}$.
\end{example}

We state without proof the central existence and uniqueness theorem; a complete proof will be given in Chapter~5 when we have the tools of the contraction mapping principle at our disposal.

\begin{theorem}[Picard--Lindel\"of]
\label{thm:picard-lindelof-preview}
\index{Picard--Lindel\"of theorem}
Consider the IVP
\[
  y' = f(t,y), \qquad y(t_0)=y_0.
\]
If $f$ is continuous on an open rectangle $R = \{(t,y) : \abs{t-t_0}<a,\;\abs{y-y_0}<b\}$ and satisfies a Lipschitz condition\index{Lipschitz condition} in $y$, i.e., there exists $L>0$ such that
\[
  \abs{f(t,y_1)-f(t,y_2)} \le L\abs{y_1-y_2}
  \quad\text{for all } (t,y_1),(t,y_2)\in R,
\]
then there exists $\delta>0$ and a unique function $\varphi\colon(t_0-\delta,\,t_0+\delta)\to\R$ satisfying the IVP.
\end{theorem}

\begin{remark}
\label{rem:lipschitz-sufficient}
If $\partial f/\partial y$ exists and is continuous on $R$, then $f$ is automatically Lipschitz in $y$ on any compact subset, so the theorem applies. The condition $\abs{\partial f/\partial y}\le L$ on $R$ suffices.
\end{remark}

% --------------------------------------------------------
\section{Direction Fields and Integral Curves}
\label{sec:direction-fields}
\index{direction field}
\index{integral curve}

For a first-order ODE $y' = f(t,y)$, the function $f$ assigns a slope to every point $(t,y)$ in its domain. By drawing short line segments of slope $f(t,y)$ at a grid of points, we obtain a \emph{direction field} (also called a \emph{slope field}\index{slope field}). A solution curve, called an \emph{integral curve}, is a curve $y = \varphi(t)$ that is tangent to the direction field at every point.

\begin{center}
\begin{tikzpicture}
  \begin{axis}[
    title={Direction field for $y' = y - t$},
    width=11cm, height=8cm,
    axis lines=middle,
    xlabel={$t$}, ylabel={$y$},
    xmin=-1.5, xmax=4, ymin=-2, ymax=4,
    xtick={-1,0,...,3}, ytick={-2,-1,...,4},
    grid=both,
    grid style={gray!20},
    every axis x label/.style={at={(ticklabel* cs:1)},anchor=west},
    every axis y label/.style={at={(ticklabel* cs:1)},anchor=south},
  ]
    % Direction field
    \foreach \xt in {-1,-0.5,...,3.5} {
      \foreach \yt in {-1.5,-1,...,3.5} {
        \pgfmathsetmacro{\slope}{\yt - \xt}
        \pgfmathsetmacro{\len}{0.2/sqrt(1+(\slope)^2)}
        \edef\temp{
          \noexpand\draw[-,blue!50!black,thin]
            (axis cs:{\xt-\len},{\yt-\slope*\len})
            -- (axis cs:{\xt+\len},{\yt+\slope*\len});
        }
        \temp
      }
    }
    % A few integral curves
    \addplot[red, very thick, domain=-1:3.5, samples=80, smooth]
      {x + 1 + 2*exp(x)};
    \addplot[red!70!black, thick, domain=-1:3.5, samples=80, smooth]
      {x + 1};
    \addplot[red, very thick, domain=-1:3.5, samples=80, smooth]
      {x + 1 - exp(x)};
  \end{axis}
\end{tikzpicture}
\end{center}

The direction field provides qualitative insight into the behaviour of solutions even when an explicit formula is unavailable. In the figure above, the red curves are integral curves of $y' = y - t$; note how they follow the pattern set by the short blue line segments.

\begin{remark}
\label{rem:direction-field-uniqueness}
By the Picard--Lindel\"of theorem, if $f$ satisfies the Lipschitz condition, integral curves cannot cross. Two distinct solution curves may approach one another but never intersect.
\end{remark}

% --------------------------------------------------------
\section{Classification of ODEs}
\label{sec:classification}
\index{classification of ODEs}

We summarise the main classification criteria in the following table.

\begin{center}
\begin{tabular}{@{}lll@{}}
\toprule
\textbf{Criterion} & \textbf{Types} & \textbf{Example} \\
\midrule
Order & First, second, \ldots, $n$-th & $y'=2y$ (first), $y''+y=0$ (second) \\
Linearity & Linear / Nonlinear & $y''+y=t$ (linear), $y''+y^2=0$ (nonlinear) \\
Autonomy & Autonomous / Non-autonomous & $y'=y^2$ (auto.), $y'=ty$ (non-auto.) \\
Homogeneity & Homogeneous / Non-homogeneous & $y''+y=0$ (hom.), $y''+y=\sin t$ (non-hom.) \\
Coefficients & Constant / Variable & $y''+3y'+2y=0$ (const.), $ty''+y=0$ (var.) \\
\bottomrule
\end{tabular}
\end{center}

% --------------------------------------------------------
\section{Historical Context}
\label{sec:history}
\index{history of ODEs}

The study of differential equations began in the late seventeenth century, immediately after the independent invention of calculus by Newton\index{Newton, Isaac} and Leibniz\index{Leibniz, Gottfried Wilhelm}. Key milestones include:

\begin{itemize}[leftmargin=1.5cm]
  \item \textbf{1690s}: The Bernoulli brothers, Jakob\index{Bernoulli, Jakob} and Johann\index{Bernoulli, Johann}, solved several families of first-order equations, including separable and what are now called Bernoulli equations.
  \item \textbf{1739}: Euler\index{Euler, Leonhard} developed systematic methods for linear equations with constant coefficients and introduced the exponential ansatz.
  \item \textbf{1760s--1780s}: Lagrange\index{Lagrange, Joseph-Louis} introduced variation of parameters and studied singular solutions.
  \item \textbf{1820s--1830s}: Cauchy\index{Cauchy, Augustin-Louis} gave the first rigorous existence proofs using the method of successive approximations.
  \item \textbf{1890}: Picard\index{Picard, \'Emile} refined Cauchy's method into the iteration scheme that bears his name, and Lindel\"of\index{Lindel\"of, Ernst} sharpened the uniqueness conditions.
  \item \textbf{1880s--1910s}: Poincar\'e\index{Poincar\'e, Henri} and Lyapunov\index{Lyapunov, Aleksandr} created the qualitative theory of ODEs, shifting focus from finding explicit formulas to understanding the geometric and stability properties of solutions.
\end{itemize}

% --------------------------------------------------------
\section{Exercises}
\label{sec:exercises-ch1}

\begin{exercise}[$\bigstar$]
\label{exer:ch1-classify}
Classify each of the following ODEs by order, linearity, and whether they are autonomous:
\begin{enumerate}[label=(\alph*)]
  \item $y' = 3y + \sin t$
  \item $y'' + y' + y^3 = 0$
  \item $(y')^2 + y = t$
  \item $y''' - 2y'' + y' = e^t$
  \item $\displaystyle\frac{\dd^2 \theta}{\dd t^2} + \sin\theta = 0$
\end{enumerate}
\end{exercise}

\begin{exercise}[$\bigstar$]
\label{exer:ch1-verify}
Verify that $y(t)=Ce^{-2t}+\frac{1}{3}e^{t}$ is the general solution of $y'+2y=e^{t}$.
\end{exercise}

\begin{exercise}[$\bigstar$]
\label{exer:ch1-ivp-solve}
Solve the IVP $y' = 3y$, $y(0) = 5$, and determine the interval of existence.
\end{exercise}

\begin{exercise}[$\bigstar\bigstar$]
\label{exer:ch1-singular}
Show that $y \equiv 0$ is a singular solution of $(y')^2 = 4y$. Find the general solution and verify that $y = 0$ is the envelope of the family.
\end{exercise}

\begin{exercise}[$\bigstar\bigstar$]
\label{exer:ch1-rc}
An RC circuit has $R = 100\,\Omega$, $C = 10^{-3}\,\text{F}$, and is connected to a constant voltage $E_0 = 12\,\text{V}$ at time $t=0$ with no initial charge. Write the IVP for $q(t)$, solve it, and find the time at which the charge reaches $90\%$ of its final value.
\end{exercise}

\begin{exercise}[$\bigstar\bigstar$]
\label{exer:ch1-logistic}
A population satisfies the logistic equation with $r = 0.5$, $K = 1000$, and $P(0) = 50$. Find $P(t)$ explicitly and determine the time at which $P = 500$.
\end{exercise}

\begin{exercise}[$\bigstar\bigstar$]
\label{exer:ch1-direction-field}
Sketch (by hand or using software) the direction field for $y' = t^2 - y$ on $[-2,2]\times[-2,4]$. Draw at least three integral curves starting from different initial conditions. What long-term behaviour do you observe?
\end{exercise}

\begin{exercise}[$\bigstar\bigstar\bigstar$]
\label{exer:ch1-lipschitz}
Let $f(t,y) = \sqrt{\abs{y}}$ on $\R^2$. Show that $f$ does not satisfy a Lipschitz condition in $y$ near $y=0$. Find two distinct solutions of the IVP $y' = \sqrt{\abs{y}}$, $y(0) = 0$, and explain why this does not contradict the Picard--Lindel\"of theorem.
\end{exercise}

\begin{exercise}[$\bigstar\bigstar\bigstar$]
\label{exer:ch1-pendulum-energy}
For the nonlinear pendulum equation $\ddot{\theta} + \omega^2\sin\theta = 0$ (with $\omega^2 = g/\ell$), define the energy
\[
  E(\theta,\dot{\theta}) = \tfrac{1}{2}\dot{\theta}^2 - \omega^2\cos\theta.
\]
Show that $E$ is constant along solutions. Sketch the level curves of $E$ in the $(\theta,\dot{\theta})$-plane and identify the equilibria and their stability.
\end{exercise}

% --------------------------------------------------------
\section*{Chapter Summary}
\addcontentsline{toc}{section}{Chapter Summary}

\begin{itemize}[leftmargin=1.5cm]
  \item An \textbf{ODE} relates an unknown function of one variable to its derivatives.
  \item ODEs are classified by \textbf{order}, \textbf{linearity}, \textbf{autonomy}, and \textbf{homogeneity}.
  \item The \textbf{general solution} of an $n$-th order ODE contains $n$ arbitrary constants; a \textbf{particular solution} results from fixing these constants; a \textbf{singular solution} lies outside the general family.
  \item An \textbf{initial value problem} specifies $n$ conditions at a single point and, under the hypotheses of the \textbf{Picard--Lindel\"of theorem}, has a unique local solution.
  \item \textbf{Direction fields} give qualitative insight into the behaviour of solutions.
  \item ODEs arise naturally in physics (mechanics, circuits), biology (population dynamics), chemistry (reaction rates), and many other fields.
\end{itemize}


% ============================================================
% ============================================================
%  CHAPTER 2: FIRST-ORDER EQUATIONS
% ============================================================
% ============================================================
\chapter{First-Order Equations}
\label{ch:first-order}
\index{first-order ODE}

In this chapter we develop systematic methods for solving the most important classes of first-order ordinary differential equations. The general first-order ODE has the form
\begin{equation}\label{eq:first-order-general}
  y' = f(t,y),
\end{equation}
or equivalently $M(t,y)\,\dd t + N(t,y)\,\dd y = 0$. While no single technique covers all cases, the methods presented here handle a very large proportion of the equations that arise in applications.

% --------------------------------------------------------
\section{Separable Equations}
\label{sec:separable}
\index{separable equation}

\begin{definition}[Separable Equation]
\label{def:separable}
A first-order ODE is \emph{separable} if it can be written in the form
\begin{equation}\label{eq:separable}
  \frac{\dd y}{\dd t} = g(t)\,h(y),
\end{equation}
where $g$ depends only on $t$ and $h$ depends only on $y$.
\end{definition}

\begin{method}[Separation of Variables]
\label{meth:separation}
\index{separation of variables}
If $h(y)\ne 0$, we may rewrite the equation as
\[
  \frac{\dd y}{h(y)} = g(t)\,\dd t,
\]
and integrate both sides:
\[
  \int \frac{\dd y}{h(y)} = \int g(t)\,\dd t + C.
\]
The resulting implicit or explicit relation defines the solution.
\end{method}

\begin{remark}
\label{rem:separable-singular}
Any constant $y_0$ satisfying $h(y_0)=0$ gives an equilibrium solution $y(t)\equiv y_0$. Such solutions may be singular; they must always be checked separately.
\end{remark}

\begin{example}
\label{ex:separable-1}
Solve $y' = ty^2$.

\textbf{Solution.} Separating variables (for $y\ne 0$):
\[
  \frac{\dd y}{y^2} = t\,\dd t
  \qquad\Longrightarrow\qquad
  -\frac{1}{y} = \frac{t^2}{2} + C
  \qquad\Longrightarrow\qquad
  y = \frac{-2}{t^2 + C_1},
\]
where $C_1 = 2C$. Additionally, $y \equiv 0$ is a constant solution.
\end{example}

\begin{example}
\label{ex:separable-logistic}
Solve the logistic equation $P' = rP(1 - P/K)$, $P(0) = P_0$, with $0 < P_0 < K$.

\textbf{Solution.} Separating variables:
\[
  \frac{\dd P}{P(1-P/K)} = r\,\dd t.
\]
Partial fractions give
\[
  \frac{1}{P(1-P/K)} = \frac{1}{P} + \frac{1/K}{1-P/K} = \frac{1}{P} + \frac{1}{K-P}.
\]
Integrating:
\[
  \ln\abs{P} - \ln\abs{K-P} = rt + C_0.
\]
Applying $P(0)=P_0$ yields $C_0 = \ln\!\left(\frac{P_0}{K-P_0}\right)$. Solving for $P$:
\begin{equation}\label{eq:logistic-solution}
  P(t) = \frac{K}{1 + \left(\frac{K-P_0}{P_0}\right)e^{-rt}}.
\end{equation}
As $t\to\infty$, $P(t)\to K$: the population approaches the carrying capacity.
\end{example}

\begin{center}
\begin{tikzpicture}
  \begin{axis}[
    title={Logistic growth curves ($r=1$, $K=10$)},
    width=11cm, height=7cm,
    axis lines=left,
    xlabel={$t$}, ylabel={$P(t)$},
    xmin=0, xmax=8, ymin=0, ymax=12,
    ytick={0,2,...,12},
    grid=both, grid style={gray!15},
    every axis x label/.style={at={(ticklabel* cs:1)},anchor=west},
    every axis y label/.style={at={(ticklabel* cs:1)},anchor=south},
    legend pos=south east,
  ]
    \addplot[thick, blue, domain=0:8, samples=50] {10 / (1 + ((10-0.5)/0.5)*exp(-x))};
    \addplot[thick, red, domain=0:8, samples=50] {10 / (1 + ((10-2)/2)*exp(-x))};
    \addplot[thick, green!50!black, domain=0:8, samples=50] {10 / (1 + ((10-5)/5)*exp(-x))};
    \addplot[thick, orange, domain=0:8, samples=50] {10 / (1 + ((10-8)/8)*exp(-x))};
    \addplot[thick, purple, domain=0:8, samples=50] {10 / (1 + ((10-12)/12)*exp(-x))};
    \addplot[dashed, black, domain=0:8] {10};
    \node at (axis cs:7,10.5) {$K=10$};
    \legend{$P_0=0.5$,$P_0=2$,$P_0=5$,$P_0=8$,$P_0=12$}
  \end{axis}
\end{tikzpicture}
\end{center}

% --------------------------------------------------------
\section{Linear First-Order Equations}
\label{sec:linear-first-order}
\index{linear first-order ODE}

\begin{definition}[Linear First-Order ODE]
\label{def:linear-first-order}
A \emph{linear first-order ODE} has the standard form
\begin{equation}\label{eq:linear-first-order}
  y' + p(t)\,y = q(t),
\end{equation}
where $p$ and $q$ are continuous functions on an interval $I$.
\end{definition}

\begin{method}[Integrating Factor]
\label{meth:integrating-factor}
\index{integrating factor!for linear equations}
Define the \emph{integrating factor}\index{integrating factor}
\begin{equation}\label{eq:integrating-factor}
  \mu(t) = e^{\int p(t)\,\dd t}.
\end{equation}
Multiplying both sides of~\eqref{eq:linear-first-order} by $\mu(t)$:
\[
  \frac{\dd}{\dd t}\!\bigl[\mu(t)\,y\bigr] = \mu(t)\,q(t).
\]
Integrating:
\begin{equation}\label{eq:linear-solution}
  y(t) = \frac{1}{\mu(t)}\left[\int \mu(t)\,q(t)\,\dd t + C\right].
\end{equation}
\end{method}

\begin{theorem}[Existence and Uniqueness for Linear First-Order ODEs]
\label{thm:linear-existence}
\index{existence and uniqueness!linear first-order}
If $p$ and $q$ are continuous on an open interval $I$ containing $t_0$, then the IVP
\[
  y' + p(t)y = q(t), \qquad y(t_0) = y_0,
\]
has a unique solution defined on the entire interval~$I$.
\end{theorem}

\begin{proof}
The integrating factor $\mu(t) = e^{\int_{t_0}^{t} p(s)\,\dd s}$ is well-defined, positive, and differentiable on~$I$. Multiplying the equation by $\mu$ and integrating from $t_0$ to $t$ gives
\[
  \mu(t)y(t) - y_0 = \int_{t_0}^{t}\mu(s)\,q(s)\,\dd s,
\]
so
\[
  y(t) = \frac{1}{\mu(t)}\left[y_0 + \int_{t_0}^{t}\mu(s)\,q(s)\,\dd s\right].
\]
This formula defines a unique continuous function on all of~$I$, and direct differentiation confirms it satisfies the ODE. Uniqueness follows because if $y_1$ and $y_2$ are two solutions with the same initial condition, then $w = y_1 - y_2$ satisfies $w' + p(t)w = 0$, $w(t_0) = 0$, whence $w(t) = 0\cdot e^{-\int_{t_0}^t p(s)\,\dd s} = 0$.
\end{proof}

\begin{example}
\label{ex:linear-1}
Solve $y' + 2ty = t$, $y(0) = 1$.

\textbf{Solution.} Here $p(t) = 2t$, $q(t) = t$. The integrating factor is
\[
  \mu(t) = e^{\int 2t\,\dd t} = e^{t^2}.
\]
Then
\[
  \frac{\dd}{\dd t}\!\bigl[e^{t^2}y\bigr] = t\,e^{t^2},
\]
so
\[
  e^{t^2}y = \int t\,e^{t^2}\,\dd t = \tfrac{1}{2}e^{t^2} + C.
\]
Hence $y = \frac{1}{2} + Ce^{-t^2}$. Applying $y(0)=1$: $C = \frac{1}{2}$. The solution is
\[
  y(t) = \frac{1}{2}\!\left(1 + e^{-t^2}\right).
\]
\end{example}

\begin{example}
\label{ex:linear-rc}
Solve the RC circuit equation $Rq' + q/C = E_0$ with $q(0)=0$.

\textbf{Solution.} Dividing by $R$: $q' + \frac{1}{RC}\,q = \frac{E_0}{R}$. With $\tau = RC$, the integrating factor is $\mu(t) = e^{t/\tau}$. Then
\[
  e^{t/\tau}q = \frac{E_0}{R}\int e^{t/\tau}\,\dd t = E_0 C\,e^{t/\tau} + K.
\]
From $q(0)=0$: $K = -E_0 C$. Therefore
\[
  q(t) = E_0 C\!\left(1 - e^{-t/\tau}\right).
\]
\end{example}

% --------------------------------------------------------
\section{Bernoulli Equations}
\label{sec:bernoulli}
\index{Bernoulli equation}

\begin{definition}[Bernoulli Equation]
\label{def:bernoulli}
A \emph{Bernoulli equation} has the form
\begin{equation}\label{eq:bernoulli}
  y' + p(t)\,y = q(t)\,y^n,
\end{equation}
where $n\in\R$, $n\ne 0,1$ (the cases $n=0$ and $n=1$ are linear).
\end{definition}

\begin{method}[Bernoulli Substitution]
\label{meth:bernoulli}
\index{Bernoulli substitution}
Setting $v = y^{1-n}$ and differentiating: $v' = (1-n)\,y^{-n}\,y'$. Dividing the Bernoulli equation by $y^n$ and multiplying by $(1-n)$:
\begin{equation}\label{eq:bernoulli-linear}
  v' + (1-n)\,p(t)\,v = (1-n)\,q(t).
\end{equation}
This is a linear first-order ODE in $v$, solvable by the integrating factor method.
\end{method}

\begin{example}
\label{ex:bernoulli-1}
Solve $y' + y = y^3$.

\textbf{Solution.} Here $n=3$, so $v = y^{-2}$, $v' = -2y^{-3}y'$. Dividing by $y^3$: $y^{-3}y' + y^{-2} = 1$, so $-\frac{1}{2}v' + v = 1$, i.e.,
\[
  v' - 2v = -2.
\]
Integrating factor: $\mu = e^{-2t}$. Then
\[
  \frac{\dd}{\dd t}\!\bigl[e^{-2t}v\bigr] = -2e^{-2t},
  \qquad
  e^{-2t}v = e^{-2t} + C,
  \qquad
  v = 1 + Ce^{2t}.
\]
Returning to $y$:
\[
  y^{-2} = 1 + Ce^{2t}
  \qquad\Longrightarrow\qquad
  y = \pm\frac{1}{\sqrt{1+Ce^{2t}}}.
\]
\end{example}

% --------------------------------------------------------
\section{Exact Equations}
\label{sec:exact}
\index{exact equation}

\begin{definition}[Exact Equation]
\label{def:exact}
The equation
\begin{equation}\label{eq:exact-form}
  M(t,y)\,\dd t + N(t,y)\,\dd y = 0
\end{equation}
is \emph{exact} if there exists a function $F(t,y)$ such that
\[
  \frac{\partial F}{\partial t} = M
  \qquad\text{and}\qquad
  \frac{\partial F}{\partial y} = N.
\]
In this case $F(t,y)=C$ defines the solution implicitly.
\end{definition}

\begin{theorem}[Exactness Criterion]
\label{thm:exactness}
\index{exactness criterion}
If $M$, $N$, $\partial M/\partial y$, and $\partial N/\partial t$ are continuous on a simply connected domain~$D$, then equation~\eqref{eq:exact-form} is exact if and only if
\begin{equation}\label{eq:exactness-condition}
  \frac{\partial M}{\partial y} = \frac{\partial N}{\partial t}
  \qquad\text{on } D.
\end{equation}
\end{theorem}

\begin{proof}
(\emph{Necessity}) If $F$ exists with $F_t = M$ and $F_y = N$, then $M_y = F_{ty} = F_{yt} = N_t$ by equality of mixed partials.

(\emph{Sufficiency}) Define
\[
  F(t,y) = \int_{t_0}^{t} M(s,y)\,\dd s + g(y),
\]
where $g$ is to be determined. Then $F_t = M$. Now require $F_y = N$:
\[
  F_y = \int_{t_0}^{t} M_y(s,y)\,\dd s + g'(y) = \int_{t_0}^{t} N_t(s,y)\,\dd s + g'(y) = N(t,y) - N(t_0,y) + g'(y).
\]
Setting $F_y = N$ gives $g'(y) = N(t_0,y)$, which depends only on $y$ and can be integrated.
\end{proof}

\begin{method}[Solving Exact Equations]
\label{meth:exact}
\index{exact equation!solving}
\begin{enumerate}[label=\textbf{Step~\arabic*:}]
  \item Check exactness: verify $M_y = N_t$.
  \item Compute $F(t,y) = \int M(t,y)\,\dd t + g(y)$ (treating $y$ as a constant in the integral).
  \item Determine $g(y)$ by differentiating $F$ with respect to $y$ and setting $F_y = N$.
  \item The solution is $F(t,y)=C$.
\end{enumerate}
\end{method}

\begin{example}
\label{ex:exact-1}
Solve $(2ty + 3)\,\dd t + (t^2 + 4y)\,\dd y = 0$.

\textbf{Solution.} Here $M = 2ty+3$ and $N = t^2+4y$. Check: $M_y = 2t$ and $N_t = 2t$. The equation is exact.

Integrate $M$ with respect to $t$:
\[
  F(t,y) = \int(2ty+3)\,\dd t = t^2 y + 3t + g(y).
\]
Now $F_y = t^2 + g'(y)$. Set equal to $N = t^2 + 4y$: $g'(y) = 4y$, so $g(y) = 2y^2$.

The solution is
\[
  t^2 y + 3t + 2y^2 = C.
\]
\end{example}

\subsection{Integrating Factors for Non-Exact Equations}
\label{subsec:integrating-factors-exact}
\index{integrating factor!for exact equations}

When $M_y \ne N_t$, we seek a function $\mu$ such that $\mu M\,\dd t + \mu N\,\dd y = 0$ is exact. This requires $(\mu M)_y = (\mu N)_t$.

\begin{proposition}
\label{prop:integrating-factor-t}
If $\displaystyle\frac{M_y - N_t}{N}$ depends only on $t$, then
\[
  \mu(t) = \exp\!\left(\int \frac{M_y - N_t}{N}\,\dd t\right)
\]
is an integrating factor depending only on $t$.
\end{proposition}

\begin{proposition}
\label{prop:integrating-factor-y}
If $\displaystyle\frac{N_t - M_y}{M}$ depends only on $y$, then
\[
  \mu(y) = \exp\!\left(\int \frac{N_t - M_y}{M}\,\dd y\right)
\]
is an integrating factor depending only on $y$.
\end{proposition}

\begin{example}
\label{ex:integrating-factor-exact}
Solve $(3y + 2t)\,\dd t + t\,\dd y = 0$.

\textbf{Solution.} Here $M = 3y+2t$, $N = t$. Check: $M_y = 3$, $N_t = 1$, so $M_y \ne N_t$. Compute
\[
  \frac{M_y - N_t}{N} = \frac{3 - 1}{t} = \frac{2}{t},
\]
which depends only on $t$. Thus $\mu(t) = e^{\int 2/t\,\dd t} = t^2$. Multiplying:
\[
  (3t^2 y + 2t^3)\,\dd t + t^3\,\dd y = 0.
\]
Now $M^* = 3t^2 y + 2t^3$, $N^* = t^3$: $M^*_y = 3t^2 = N^*_t$. Exact. Integrate:
\[
  F = \int t^3\,\dd y = t^3 y + h(t).
\]
Then $F_t = 3t^2 y + h'(t) = 3t^2 y + 2t^3$, so $h'(t) = 2t^3$, $h(t) = \frac{t^4}{2}$.

The solution is $t^3 y + \frac{t^4}{2} = C$.
\end{example}

% --------------------------------------------------------
\section{Homogeneous Equations}
\label{sec:homogeneous-first-order}
\index{homogeneous equation (first-order)}

\begin{definition}[Homogeneous Function]
\label{def:homogeneous-function}
\index{homogeneous function}
A function $f(t,y)$ is \emph{homogeneous of degree $k$} if $f(\lambda t, \lambda y) = \lambda^k f(t,y)$ for all $\lambda > 0$.
\end{definition}

\begin{definition}[Homogeneous First-Order ODE]
\label{def:homogeneous-ode}
A first-order ODE $M(t,y)\,\dd t + N(t,y)\,\dd y = 0$ is \emph{homogeneous} (in the sense of homogeneous functions) if $M$ and $N$ are both homogeneous of the same degree. Equivalently, the equation can be written as
\begin{equation}\label{eq:homogeneous-form}
  \frac{\dd y}{\dd t} = \Phi\!\left(\frac{y}{t}\right)
\end{equation}
for some function $\Phi$.
\end{definition}

\begin{method}[Substitution $y = ut$]
\label{meth:homogeneous-sub}
\index{substitution!$y = ut$}
Setting $y = ut$ (so $u = y/t$), we have $y' = u + tu'$. Equation~\eqref{eq:homogeneous-form} becomes
\[
  u + tu' = \Phi(u),
\]
which is separable:
\[
  \frac{\dd u}{\Phi(u)-u} = \frac{\dd t}{t}.
\]
\end{method}

\begin{example}
\label{ex:homogeneous-1}
Solve $t\,y' = y + \sqrt{t^2 + y^2}$, $t > 0$.

\textbf{Solution.} Rewrite as $y' = \frac{y}{t} + \sqrt{1 + (y/t)^2}$, which has the form $y' = \Phi(y/t)$ with $\Phi(u) = u + \sqrt{1+u^2}$. Setting $y = ut$:
\[
  u + tu' = u + \sqrt{1+u^2},
\]
so $tu' = \sqrt{1+u^2}$. Separating:
\[
  \frac{\dd u}{\sqrt{1+u^2}} = \frac{\dd t}{t}.
\]
Integrating: $\sinh^{-1}(u) = \ln t + C_1$, i.e., $u = \sinh(\ln t + C_1)$. Finally,
\[
  y = t\,\sinh(\ln t + C_1).
\]
\end{example}

% --------------------------------------------------------
\section{Riccati Equations}
\label{sec:riccati}
\index{Riccati equation}

\begin{definition}[Riccati Equation]
\label{def:riccati}
A \emph{Riccati equation} has the form
\begin{equation}\label{eq:riccati}
  y' = P(t) + Q(t)\,y + R(t)\,y^2.
\end{equation}
\end{definition}

In general, Riccati equations cannot be solved in closed form. However, if one particular solution $y_1(t)$ is known, the substitution $y = y_1 + 1/v$ reduces the equation to a linear first-order ODE in $v$.

\begin{method}[Riccati Reduction]
\label{meth:riccati}
\index{Riccati reduction}
Suppose $y_1(t)$ is a known particular solution of~\eqref{eq:riccati}. Set $y = y_1 + \frac{1}{v}$. Then
\[
  y' = y_1' - \frac{v'}{v^2}.
\]
Substituting into~\eqref{eq:riccati} and simplifying (using $y_1' = P + Qy_1 + Ry_1^2$):
\begin{equation}\label{eq:riccati-reduced}
  v' + (Q + 2Ry_1)\,v = -R.
\end{equation}
This is a linear ODE in $v$.
\end{method}

\begin{example}
\label{ex:riccati-1}
Solve $y' = 1 + t^2 - 2ty + y^2$, given that $y_1(t) = t$ is a particular solution.

\textbf{Solution.} Check: $y_1' = 1$ and $1 + t^2 - 2t\cdot t + t^2 = 1$. Confirmed.

Here $P(t)=1+t^2$, $Q(t)=-2t$, $R(t)=1$. Setting $y = t + 1/v$:
\[
  v' + (-2t + 2t)\,v = -1 \qquad\Longrightarrow\qquad v' = -1,
\]
so $v = -t + C$. Therefore
\[
  y = t + \frac{1}{C - t}.
\]
\end{example}

% --------------------------------------------------------
\section{Clairaut Equations}
\label{sec:clairaut}
\index{Clairaut equation}

\begin{definition}[Clairaut Equation]
\label{def:clairaut}
A \emph{Clairaut equation} has the form
\begin{equation}\label{eq:clairaut}
  y = ty' + \psi(y'),
\end{equation}
where $\psi$ is a given function.
\end{definition}

\begin{method}[Solving Clairaut Equations]
\label{meth:clairaut}
\index{Clairaut equation!solution}
Differentiate~\eqref{eq:clairaut} with respect to $t$:
\[
  y' = y' + (t + \psi'(y'))\,y''.
\]
Hence either $y'' = 0$ or $t + \psi'(y') = 0$.

\textbf{Case 1:} $y'' = 0$ means $y' = c$ (constant), and the general solution is the family of straight lines
\[
  y = ct + \psi(c).
\]

\textbf{Case 2:} $t = -\psi'(p)$ with $p = y'$, combined with $y = tp + \psi(p)$, gives a parametric representation of the \emph{singular solution}, which is the envelope of the family.
\end{method}

\begin{example}
\label{ex:clairaut-1}
Solve $y = ty' + (y')^2$.

\textbf{Solution.} This is a Clairaut equation with $\psi(p) = p^2$.

\textit{General solution:} $y = ct + c^2$ (a family of lines).

\textit{Singular solution:} From $t + \psi'(p) = 0$: $t + 2p = 0$, so $p = -t/2$. Then
\[
  y = t\!\left(-\frac{t}{2}\right) + \left(-\frac{t}{2}\right)^2 = -\frac{t^2}{2} + \frac{t^2}{4} = -\frac{t^2}{4}.
\]
The singular solution is the parabola $y = -t^2/4$, which is the envelope of the family of lines.
\end{example}

\begin{center}
\begin{tikzpicture}
  \begin{axis}[
    title={Clairaut equation: $y = ty' + (y')^2$},
    width=11cm, height=8cm,
    axis lines=middle,
    xlabel={$t$}, ylabel={$y$},
    xmin=-5, xmax=5, ymin=-5, ymax=5,
    grid=both, grid style={gray!15},
    every axis x label/.style={at={(ticklabel* cs:1)},anchor=west},
    every axis y label/.style={at={(ticklabel* cs:1)},anchor=south},
  ]
    % General solution: y = ct + c^2 for various c
    \foreach \c in {-3,-2.5,...,3} {
      \addplot[blue!60!black, thin, domain=-5:5, samples=2]
        {\c*x + \c*\c};
    }
    % Singular solution: y = -t^2/4
    \addplot[red, very thick, domain=-5:5, samples=100, smooth]
      {-x*x/4};
    \node[red] at (axis cs:3.5,-3.5) {$y=-\tfrac{t^2}{4}$};
  \end{axis}
\end{tikzpicture}
\end{center}

% --------------------------------------------------------
\section{Lagrange Equations}
\label{sec:lagrange}
\index{Lagrange equation}

\begin{definition}[Lagrange Equation]
\label{def:lagrange}
A \emph{Lagrange equation} (also called a d'Alembert equation\index{d'Alembert equation}) has the form
\begin{equation}\label{eq:lagrange}
  y = t\,\varphi(y') + \psi(y'),
\end{equation}
where $\varphi$ and $\psi$ are given functions. A Clairaut equation is the special case $\varphi(p) = p$.
\end{definition}

\begin{method}[Solving Lagrange Equations]
\label{meth:lagrange}
\index{Lagrange equation!solution}
Set $p = y'$ and differentiate~\eqref{eq:lagrange} with respect to $t$:
\[
  p = \varphi(p) + t\,\varphi'(p)\,p' + \psi'(p)\,p'.
\]
If $\varphi(p)\ne p$, we can rearrange:
\[
  \frac{\dd t}{\dd p} = \frac{t\,\varphi'(p) + \psi'(p)}{p - \varphi(p)},
\]
which is a linear first-order ODE in $t$ as a function of $p$. Solving this and substituting back via~\eqref{eq:lagrange} gives the solution parametrically in terms of $p$.
\end{method}

\begin{example}
\label{ex:lagrange-1}
Solve $y = 2ty' - (y')^2$.

\textbf{Solution.} Here $\varphi(p) = 2p$, $\psi(p) = -p^2$. Differentiating:
\[
  p = 2p + 2tp' - 2pp',
\]
so
\[
  -p = (2t - 2p)\,p',
  \qquad\text{i.e.,}\qquad
  \frac{\dd t}{\dd p} = \frac{2t - 2p}{-p} = -\frac{2t}{p} + 2.
\]
This is linear: $t' + \frac{2}{p}\,t = 2$. Integrating factor: $\mu = p^2$. Then
\[
  \frac{\dd}{\dd p}\!\bigl[p^2 t\bigr] = 2p^2,
  \qquad
  p^2 t = \frac{2p^3}{3} + C,
  \qquad
  t = \frac{2p}{3} + \frac{C}{p^2}.
\]
From the original equation: $y = 2tp - p^2 = 2p\!\left(\frac{2p}{3}+\frac{C}{p^2}\right) - p^2 = \frac{p^2}{3} + \frac{2C}{p}$.

The solution is given parametrically:
\[
  \begin{cases}
    t = \dfrac{2p}{3} + \dfrac{C}{p^2}, \\[8pt]
    y = \dfrac{p^2}{3} + \dfrac{2C}{p}.
  \end{cases}
\]
\end{example}

% --------------------------------------------------------
\section{Table of Standard Forms}
\label{sec:standard-forms-table}
\index{standard forms table}

\begin{center}
\renewcommand{\arraystretch}{1.6}
\begin{tabular}{@{}p{3cm}p{4.5cm}p{5.5cm}@{}}
\toprule
\textbf{Type} & \textbf{Standard Form} & \textbf{Method} \\
\midrule
Separable\index{separable equation} &
  $y' = g(t)\,h(y)$ &
  Separate and integrate \\
Linear\index{linear first-order ODE} &
  $y' + p(t)y = q(t)$ &
  Integrating factor $\mu = e^{\int p\,\dd t}$ \\
Bernoulli\index{Bernoulli equation} &
  $y' + p(t)y = q(t)y^n$ &
  $v = y^{1-n}$ reduces to linear \\
Exact\index{exact equation} &
  $M\,\dd t + N\,\dd y = 0$, $M_y = N_t$ &
  Find $F$ with $F_t=M$, $F_y=N$ \\
Homogeneous\index{homogeneous equation (first-order)} &
  $y' = \Phi(y/t)$ &
  $y = ut$ reduces to separable \\
Riccati\index{Riccati equation} &
  $y' = P + Qy + Ry^2$ &
  Know $y_1$; set $y = y_1+1/v$ \\
Clairaut\index{Clairaut equation} &
  $y = ty' + \psi(y')$ &
  General: $y=ct+\psi(c)$; singular: envelope \\
Lagrange\index{Lagrange equation} &
  $y = t\varphi(y')+\psi(y')$ &
  Differentiate; linear ODE in $t(p)$ \\
\bottomrule
\end{tabular}
\end{center}

% --------------------------------------------------------
\section{Additional Worked Examples}
\label{sec:additional-examples}

\begin{example}
\label{ex:extra-separable}
Solve $\displaystyle\frac{\dd y}{\dd t} = \frac{t^2}{1+y^2}$.

\textbf{Solution.} Separating: $(1+y^2)\,\dd y = t^2\,\dd t$. Integrating:
\[
  y + \frac{y^3}{3} = \frac{t^3}{3} + C.
\]
This defines $y$ implicitly as a function of $t$.
\end{example}

\begin{example}
\label{ex:extra-linear}
Solve $y' - \frac{2y}{t} = t^2\cos t$, $t > 0$.

\textbf{Solution.} Here $p(t) = -2/t$, so
\[
  \mu(t) = e^{\int -2/t\,\dd t} = e^{-2\ln t} = t^{-2}.
\]
Then
\[
  \frac{\dd}{\dd t}\!\left[\frac{y}{t^2}\right] = \cos t,
\]
so $\frac{y}{t^2} = \sin t + C$, giving $y = t^2(\sin t + C)$.
\end{example}

\begin{example}
\label{ex:extra-bernoulli}
Solve $y' + \frac{y}{t} = t\sqrt{y}$, $t > 0$, $y > 0$.

\textbf{Solution.} This is a Bernoulli equation with $n = 1/2$. Set $v = y^{1/2}$, so $v' = \frac{1}{2}y^{-1/2}y'$. Dividing by $y^{1/2}$: $y^{-1/2}y' + \frac{y^{1/2}}{t} = t$, so $2v' + \frac{v}{t} = t$, i.e.,
\[
  v' + \frac{v}{2t} = \frac{t}{2}.
\]
Integrating factor: $\mu = e^{\int 1/(2t)\,\dd t} = t^{1/2}$. Then
\[
  \frac{\dd}{\dd t}\!\bigl[t^{1/2}\,v\bigr] = \frac{t^{3/2}}{2},
  \qquad
  t^{1/2}\,v = \frac{t^{5/2}}{5} + C.
\]
So $v = \frac{t^2}{5} + Ct^{-1/2}$, and $y = v^2 = \left(\frac{t^2}{5} + Ct^{-1/2}\right)^2$.
\end{example}

\begin{example}
\label{ex:extra-exact}
Solve $(e^y + 2ty)\,\dd t + (te^y + t^2 + 1)\,\dd y = 0$.

\textbf{Solution.} Here $M = e^y + 2ty$, $N = te^y + t^2 + 1$. Check:
\[
  M_y = e^y + 2t, \qquad N_t = e^y + 2t.
\]
Exact. Integrate $M$ with respect to $t$:
\[
  F = te^y + t^2 y + g(y).
\]
Then $F_y = te^y + t^2 + g'(y) = te^y + t^2 + 1$, so $g'(y) = 1$, $g(y)=y$. The solution is
\[
  te^y + t^2 y + y = C.
\]
\end{example}

\begin{example}
\label{ex:extra-homogeneous}
Solve $(t^2 + y^2)\,\dd t - 2ty\,\dd y = 0$.

\textbf{Solution.} Both $M = t^2+y^2$ and $N = -2ty$ are homogeneous of degree~$2$. Set $y = ut$:
\[
  y' = \frac{t^2+u^2t^2}{2t\cdot ut} = \frac{1+u^2}{2u}.
\]
Since $y' = u + tu'$:
\[
  tu' = \frac{1+u^2}{2u} - u = \frac{1-u^2}{2u}.
\]
Separating:
\[
  \frac{2u\,\dd u}{1-u^2} = \frac{\dd t}{t}.
\]
Integrating: $-\ln\abs{1-u^2} = \ln\abs{t} + C_0$, so $\frac{1}{\abs{1-u^2}} = A\abs{t}$ where $A = e^{C_0}$. Returning to $y$: $u = y/t$, so
\[
  \frac{t^2}{t^2 - y^2} = At,
  \qquad\text{i.e.,}\qquad
  t = A(t^2 - y^2).
\]
Equivalently, $t^2 - y^2 = Ct$ for some constant $C$.
\end{example}

\begin{center}
\begin{tikzpicture}
  \begin{axis}[
    title={Solution curves for $(t^2+y^2)\,\dd t = 2ty\,\dd y$},
    width=10cm, height=8cm,
    axis lines=middle,
    xlabel={$t$}, ylabel={$y$},
    xmin=-1, xmax=6, ymin=-4, ymax=4,
    grid=both, grid style={gray!15},
    every axis x label/.style={at={(ticklabel* cs:1)},anchor=west},
    every axis y label/.style={at={(ticklabel* cs:1)},anchor=south},
    samples=200,
  ]
    % t^2 - y^2 = Ct  =>  y^2 = t^2 - Ct  =>  y = +/- sqrt(t^2 - Ct)
    \foreach \cc in {-4,-3,-2,-1,0,1,2,3,4,5} {
      \addplot[blue!70!black, thick, domain=0.01:5.5, smooth]
        ({x},{sqrt(max(x*x - \cc*x, 0.0001))});
      \addplot[blue!70!black, thick, domain=0.01:5.5, smooth]
        ({x},{-sqrt(max(x*x - \cc*x, 0.0001))});
    }
  \end{axis}
\end{tikzpicture}
\end{center}

% --------------------------------------------------------
\section{Exercises}
\label{sec:exercises-ch2}

\subsection*{Separable Equations}

\begin{exercise}[$\bigstar$]
\label{exer:ch2-sep1}
Solve $y' = \dfrac{t}{y}$, $y(1) = 2$.
\end{exercise}

\begin{exercise}[$\bigstar$]
\label{exer:ch2-sep2}
Solve $y' = y\sin t$, $y(0) = 1$.
\end{exercise}

\begin{exercise}[$\bigstar\bigstar$]
\label{exer:ch2-sep3}
Solve $\dfrac{\dd y}{\dd t} = \dfrac{t^2 y - y}{t+1}$ and determine all equilibrium solutions.
\end{exercise}

\subsection*{Linear Equations}

\begin{exercise}[$\bigstar$]
\label{exer:ch2-lin1}
Solve $y' + 3y = 6$, $y(0)=1$.
\end{exercise}

\begin{exercise}[$\bigstar$]
\label{exer:ch2-lin2}
Solve $ty' + 2y = t^3$, $t>0$, $y(1) = 0$.
\end{exercise}

\begin{exercise}[$\bigstar\bigstar$]
\label{exer:ch2-lin3}
Solve $y' + y\tan t = \sec t$, $-\pi/2 < t < \pi/2$, $y(0) = 0$.
\end{exercise}

\subsection*{Bernoulli Equations}

\begin{exercise}[$\bigstar\bigstar$]
\label{exer:ch2-ber1}
Solve $y' - y = -y^2 e^t$.
\end{exercise}

\begin{exercise}[$\bigstar\bigstar$]
\label{exer:ch2-ber2}
Solve $y' + \dfrac{y}{t} = \dfrac{\ln t}{y}$, $t > 0$, $y > 0$.
\end{exercise}

\subsection*{Exact Equations}

\begin{exercise}[$\bigstar$]
\label{exer:ch2-exact1}
Solve $(2t + 3y)\,\dd t + (3t + 2y)\,\dd y = 0$.
\end{exercise}

\begin{exercise}[$\bigstar\bigstar$]
\label{exer:ch2-exact2}
Show that $(y^2 + t)\,\dd t + 2ty\,\dd y = 0$ is not exact, find an integrating factor depending on $t$, and solve.
\end{exercise}

\begin{exercise}[$\bigstar\bigstar$]
\label{exer:ch2-exact3}
Solve $\left(\dfrac{1}{t} + \dfrac{y}{t^2+y^2}\right)\dd t + \left(\dfrac{t}{t^2+y^2}\right)\dd y = 0$, $t > 0$.
\end{exercise}

\subsection*{Homogeneous Equations}

\begin{exercise}[$\bigstar$]
\label{exer:ch2-hom1}
Solve $y' = \dfrac{y+t}{t}$.
\end{exercise}

\begin{exercise}[$\bigstar\bigstar$]
\label{exer:ch2-hom2}
Solve $(t^2 + 3y^2)\,\dd t - 2ty\,\dd y = 0$.
\end{exercise}

\subsection*{Riccati, Clairaut, Lagrange}

\begin{exercise}[$\bigstar\bigstar$]
\label{exer:ch2-ric1}
Solve $y' = -y^2 + \dfrac{2}{t^2}$, given that $y_1 = \dfrac{1}{t}$ is a particular solution.
\end{exercise}

\begin{exercise}[$\bigstar\bigstar$]
\label{exer:ch2-clair1}
Solve the Clairaut equation $y = ty' - (y')^3$. Find both the general solution and the singular solution.
\end{exercise}

\begin{exercise}[$\bigstar\bigstar\bigstar$]
\label{exer:ch2-lag1}
Solve $y = t(y')^2 - (y')^3$ parametrically.
\end{exercise}

\subsection*{Mixed and Challenging Problems}

\begin{exercise}[$\bigstar\bigstar$]
\label{exer:ch2-mixed1}
For each equation below, identify the type and solve:
\begin{enumerate}[label=(\alph*)]
  \item $(3t^2 y + y^3)\,\dd t + (t^3 + 3ty^2)\,\dd y = 0$
  \item $y' = \dfrac{y^2 - t^2}{2ty}$
  \item $y' + y\cot t = 2\cos t$
\end{enumerate}
\end{exercise}

\begin{exercise}[$\bigstar\bigstar\bigstar$]
\label{exer:ch2-existence}
Consider the IVP $y' = 3y^{2/3}$, $y(0) = 0$.
\begin{enumerate}[label=(\alph*)]
  \item Show that $y \equiv 0$ is a solution.
  \item Show that for any $a \ge 0$, the function
    \[
      y_a(t) = \begin{cases} 0 & \text{if } t \le a, \\ (t-a)^3 & \text{if } t > a, \end{cases}
    \]
    is also a solution.
  \item Explain why this does not contradict the Picard--Lindel\"of theorem.
\end{enumerate}
\end{exercise}

\begin{exercise}[$\bigstar\bigstar\bigstar$]
\label{exer:ch2-orthogonal}
\index{orthogonal trajectories}
The family of curves $y = Ct^2$ fills the upper half-plane. Find the \emph{orthogonal trajectories}\index{orthogonal trajectories} of this family, i.e., the curves that intersect each member of the family at right angles. Sketch both families.
\end{exercise}

% --------------------------------------------------------
\section*{Chapter Summary}
\addcontentsline{toc}{section}{Chapter Summary}

\begin{itemize}[leftmargin=1.5cm]
  \item A \textbf{separable} equation $y' = g(t)h(y)$ is solved by separating variables and integrating; always check for equilibrium (possibly singular) solutions $h(y_0)=0$.
  \item A \textbf{linear} equation $y' + py = q$ is solved using the integrating factor $\mu = e^{\int p\,\dd t}$. The solution exists and is unique on any interval where $p$ and $q$ are continuous.
  \item A \textbf{Bernoulli} equation $y' + py = qy^n$ is reduced to linear form by $v = y^{1-n}$.
  \item An equation $M\,\dd t + N\,\dd y = 0$ is \textbf{exact} when $M_y = N_t$; the solution is $F(t,y)=C$ where $F_t = M$, $F_y = N$. Non-exact equations may admit integrating factors $\mu(t)$ or $\mu(y)$.
  \item A \textbf{homogeneous} equation (equal-degree numerator and denominator) is reduced to separable form by $y = ut$.
  \item A \textbf{Riccati} equation requires one known solution; then a substitution reduces it to a linear equation.
  \item \textbf{Clairaut} and \textbf{Lagrange} equations are solved by differentiation: Clairaut yields a family of lines and a singular envelope; Lagrange leads to a linear ODE in $t(p)$.
  \item The table of standard forms (Section~\ref{sec:standard-forms-table}) is a reliable guide for identifying the type of a given first-order ODE.
\end{itemize}

% === END OF CHUNK preamble+ch1-2 ===
% === START OF CHUNK ch3-4 ===

% ══════════════════════════════════════════════════════════════
%  CHAPTER 3 — EXISTENCE AND UNIQUENESS THEOREMS
% ══════════════════════════════════════════════════════════════
\chapter{Existence and Uniqueness Theorems}\label{ch:existence-uniqueness}
\index{existence and uniqueness}

The previous chapters showed us how to solve certain classes of ODEs explicitly.
But explicit formulae are the exception rather than the rule: most differential
equations cannot be solved in closed form. This raises fundamental questions.
Given an initial value problem
\[
  y' = f(t,y), \qquad y(t_0) = y_0,
\]
does a solution exist at all? If it does, is it unique? How far can it be
extended? And does the solution depend continuously on the initial data?
This chapter provides rigorous answers via the Picard--Lindel\"of theorem,
Peano's theorem, and Gr\"onwall's lemma.\index{initial value problem}

% ────────────────────────────────────────────────
\section{Motivating Examples}\label{sec:exist-motivation}
\index{non-uniqueness}

Before stating any theorems, let us see that the questions above are not
vacuous.

\begin{example}[Non-uniqueness]\label{ex:non-unique-y23}
\index{non-uniqueness!example $y'=y^{2/3}$}
Consider the initial value problem
\[
  y' = y^{2/3}, \qquad y(0) = 0.
\]
One solution is $y(t) = 0$ for all $t$. Another is
\[
  y(t) = \begin{cases} 0 & t \le c, \\[4pt]
  \displaystyle\left(\frac{t-c}{3}\right)^{3} & t > c,
  \end{cases}
\]
for any $c \ge 0$. Thus uniqueness fails: there are infinitely many solutions
through the origin. The problem is that $f(t,y) = y^{2/3}$ is not Lipschitz
in $y$ near $y = 0$.
\end{example}

\begin{example}[Finite-time blow-up]\label{ex:blowup-y2}
\index{blow-up}
Consider
\[
  y' = y^2, \qquad y(0) = 1.
\]
Separation of variables gives $y(t) = \dfrac{1}{1-t}$, which blows up at
$t = 1$. The solution exists on $(-\infty, 1)$ but cannot be extended to all
of $\R$. Existence is only \emph{local} in general.
\end{example}

\begin{center}
\begin{tikzpicture}
  \begin{axis}[
    width=11cm, height=7cm,
    xlabel={$t$}, ylabel={$y$},
    xmin=-0.5, xmax=2.5, ymin=-0.5, ymax=3.5,
    axis lines=middle,
    title={Non-unique solutions of $y' = y^{2/3}$, $y(0)=0$},
    legend pos=north west,
    samples=200, thick,
  ]
    % Trivial solution
    \addplot[blue, domain=-0.5:2.5] {0};
    \addlegendentry{$y \equiv 0$}
    % Non-trivial with c=0
    \addplot[red, domain=0:2.5] {(x/3)^3};
    \addlegendentry{$y = (t/3)^3$}
    % Non-trivial with c=1
    \addplot[orange, domain=0:2.5] {ifthenelse(x<=1, 0, ((x-1)/3)^3)};
    \addlegendentry{$y = ((t-1)/3)^3$ for $t>1$}
    % Non-trivial with c=0.5
    \addplot[green!60!black, domain=0:2.5] {ifthenelse(x<=0.5, 0, ((x-0.5)/3)^3)};
    \addlegendentry{$y = ((t-0.5)/3)^3$ for $t>0.5$}
  \end{axis}
\end{tikzpicture}
\end{center}

% ────────────────────────────────────────────────
\section{The Lipschitz Condition}\label{sec:lipschitz}
\index{Lipschitz condition}

The key hypothesis that guarantees uniqueness is a Lipschitz condition on
the right-hand side.

\begin{definition}[Lipschitz condition]\label{def:lipschitz}
\index{Lipschitz condition!definition}
Let $D \subset \R \times \R^n$ be open. A function $f\colon D \to \R^n$ is
\textbf{Lipschitz in $y$} (uniformly in $t$) on $D$ if there exists a constant
$L \ge 0$ such that
\[
  \norm{f(t,y_1) - f(t,y_2)} \le L\,\norm{y_1 - y_2}
\]
for all $(t,y_1), (t,y_2) \in D$. The constant $L$ is called a
\textbf{Lipschitz constant}\index{Lipschitz constant}.
\end{definition}

\begin{definition}[Locally Lipschitz]\label{def:locally-lipschitz}
\index{Lipschitz condition!locally Lipschitz}
The function $f$ is \textbf{locally Lipschitz in $y$} on $D$ if for each
compact $K \subset D$, the restriction $f|_K$ satisfies a Lipschitz condition
(with a constant $L_K$ that may depend on $K$).
\end{definition}

\begin{proposition}[Sufficient condition for Lipschitz]\label{prop:C1-lipschitz}
\index{Lipschitz condition!sufficient condition}
If $f\colon D \to \R^n$ is continuous and $\dfrac{\partial f}{\partial y}$
exists and is continuous on the open set $D$, then $f$ is locally Lipschitz
in $y$ on $D$.
\end{proposition}

\begin{proof}
Let $K \subset D$ be a compact convex set. Since $\partial f/\partial y$ is
continuous on $K$, it is bounded: there exists $L \ge 0$ with
$\norm{\partial f/\partial y(t,y)} \le L$ for all $(t,y) \in K$.
For $(t,y_1), (t,y_2) \in K$, the mean value inequality gives
\[
  \norm{f(t,y_1) - f(t,y_2)}
  \le \sup_{s \in [0,1]} \norm*{\frac{\partial f}{\partial y}\bigl(t, y_2 + s(y_1 - y_2)\bigr)}
       \cdot \norm{y_1 - y_2}
  \le L\,\norm{y_1 - y_2}. \qedhere
\]
\end{proof}

\begin{example}[Checking the Lipschitz condition]\label{ex:lipschitz-check}
\begin{enumerate}[label=(\alph*)]
\item $f(t,y) = ty + \sin y$. We have $\partial f/\partial y = t + \cos y$,
  which is continuous everywhere. Hence $f$ is locally Lipschitz.
\item $f(t,y) = y^{2/3}$. Here $\partial f/\partial y = \frac{2}{3}\,y^{-1/3}$,
  which is unbounded near $y=0$. Indeed, $f$ is \emph{not} Lipschitz near $y=0$.
\item $f(t,y) = y^2$. Then $\partial f/\partial y = 2y$, continuous everywhere.
  So $f$ is locally Lipschitz but not globally Lipschitz (the derivative is
  unbounded as $\abs{y} \to \infty$).
\end{enumerate}
\end{example}

% ────────────────────────────────────────────────
\section{The Picard--Lindel\"of Theorem}\label{sec:picard-lindelof}
\index{Picard--Lindel\"of theorem}
\index{Cauchy--Lipschitz theorem|see{Picard--Lindel\"of theorem}}

\subsection{Integral formulation}

The starting point is to reformulate the initial value problem as an integral
equation.

\begin{proposition}[Integral equation equivalence]\label{prop:integral-eq}
\index{integral equation!equivalence}
Let $f$ be continuous on an open set $D \subset \R \times \R^n$ containing
$(t_0, y_0)$. A continuous function $y\colon I \to \R^n$ (where $I$ is an
interval containing $t_0$) is a solution of
\[
  y' = f(t,y), \qquad y(t_0) = y_0,
\]
if and only if it satisfies the integral equation
\begin{equation}\label{eq:volterra}
  y(t) = y_0 + \int_{t_0}^{t} f\bigl(s, y(s)\bigr)\,\dd s
  \qquad \text{for all } t \in I.
\end{equation}
\end{proposition}

\begin{proof}
If $y$ solves the IVP, integrate both sides of $y'(s) = f(s,y(s))$ from $t_0$
to $t$ and use $y(t_0) = y_0$. Conversely, if \eqref{eq:volterra} holds,
differentiate with respect to $t$ (using continuity of the integrand) to
recover $y'(t) = f(t,y(t))$, and setting $t = t_0$ gives $y(t_0) = y_0$.
\end{proof}

\subsection{Picard iteration}\label{subsec:picard-iteration}
\index{Picard iteration}

The idea is to define a sequence of successive approximations:
\begin{equation}\label{eq:picard-iterates}
\begin{aligned}
  y_0(t) &= y_0, \\
  y_{n+1}(t) &= y_0 + \int_{t_0}^{t} f\bigl(s, y_n(s)\bigr)\,\dd s,
  \qquad n = 0, 1, 2, \dotsc
\end{aligned}
\end{equation}
If this sequence converges (uniformly), its limit $y(t) = \lim_{n\to\infty} y_n(t)$
will satisfy the integral equation~\eqref{eq:volterra} and hence solve the IVP.

\begin{example}[Picard iterates for $y' = y$, $y(0)=1$]\label{ex:picard-exp}
\index{Picard iteration!example}
Here $f(t,y) = y$, $t_0 = 0$, $y_0 = 1$. We compute:
\begin{align*}
  y_0(t) &= 1, \\
  y_1(t) &= 1 + \int_0^t 1\,\dd s = 1 + t, \\
  y_2(t) &= 1 + \int_0^t (1+s)\,\dd s = 1 + t + \frac{t^2}{2}, \\
  y_3(t) &= 1 + \int_0^t \Bigl(1 + s + \frac{s^2}{2}\Bigr)\dd s
           = 1 + t + \frac{t^2}{2} + \frac{t^3}{6}.
\end{align*}
By induction, $y_n(t) = \displaystyle\sum_{k=0}^{n} \frac{t^k}{k!}$, which
converges to $e^t$ as $n \to \infty$.
\end{example}

\begin{center}
\begin{tikzpicture}
  \begin{axis}[
    width=11cm, height=7cm,
    xlabel={$t$}, ylabel={$y$},
    xmin=-0.3, xmax=3, ymin=0, ymax=8,
    axis lines=middle,
    title={Picard iterates converging to $e^t$},
    legend pos=north west,
    samples=100, thick,
  ]
    \addplot[gray!50, very thick, domain=0:2.8] {exp(x)};
    \addlegendentry{$e^t$ (exact)}
    \addplot[blue, dashed, domain=0:2.8] {1};
    \addlegendentry{$y_0 = 1$}
    \addplot[red, dashed, domain=0:2.8] {1+x};
    \addlegendentry{$y_1 = 1+t$}
    \addplot[orange, domain=0:2.8] {1+x+x^2/2};
    \addlegendentry{$y_2$}
    \addplot[green!60!black, domain=0:2.8] {1+x+x^2/2+x^3/6};
    \addlegendentry{$y_3$}
    \addplot[purple, domain=0:2.8] {1+x+x^2/2+x^3/6+x^4/24};
    \addlegendentry{$y_4$}
  \end{axis}
\end{tikzpicture}
\end{center}

\begin{example}[Picard iterates for $y' = t + y^2$, $y(0) = 0$]\label{ex:picard-nonlinear}
We compute the first few iterates:
\begin{align*}
  y_0(t) &= 0, \\
  y_1(t) &= \int_0^t (s + 0)\,\dd s = \frac{t^2}{2}, \\
  y_2(t) &= \int_0^t \biggl(s + \frac{s^4}{4}\biggr)\dd s
           = \frac{t^2}{2} + \frac{t^5}{20}, \\
  y_3(t) &= \int_0^t \biggl(s + \Bigl(\frac{s^2}{2} + \frac{s^5}{20}\Bigr)^{\!2}\biggr)\dd s
           = \frac{t^2}{2} + \frac{t^5}{20} + \frac{t^8}{160} + \frac{t^{11}}{4400}.
\end{align*}
Even without a closed form, the iterates give systematic polynomial approximations
to the solution.
\end{example}

\subsection{Statement and proof}

\begin{theorem}[Picard--Lindel\"of]\label{thm:picard-lindelof}
\index{Picard--Lindel\"of theorem!statement and proof}
Let $f\colon D \to \R^n$ be continuous on the open set
$D \subset \R \times \R^n$, and suppose $f$ is Lipschitz in $y$ with constant
$L$ on the closed rectangle
\[
  R = \bigl\{(t,y) \in \R \times \R^n : \abs{t - t_0} \le a,\;
  \norm{y - y_0} \le b\bigr\} \subset D.
\]
Let $M = \max_{(t,y) \in R} \norm{f(t,y)}$ and
$\alpha = \min\bigl(a, \, b/M\bigr)$. Then the initial value problem
\[
  y' = f(t,y), \qquad y(t_0) = y_0,
\]
has a \textbf{unique} solution $y \in C^1\bigl([t_0 - \alpha,\, t_0 + \alpha];\, \R^n\bigr)$.
\end{theorem}

\begin{proof}
We use the Banach fixed point theorem\index{Banach fixed point theorem}
applied to the Picard operator on an appropriate complete metric space.

\medskip
\textbf{Step 1: Setup.}
Let $I = [t_0 - \alpha, \, t_0 + \alpha]$ and define the space
\[
  X = \bigl\{\varphi \in C(I, \R^n) : \norm{\varphi(t) - y_0} \le b
  \text{ for all } t \in I\bigr\},
\]
equipped with the weighted supremum norm
\[
  \norm{\varphi}_\lambda = \sup_{t \in I}
  e^{-\lambda \abs{t - t_0}} \norm{\varphi(t) - y_0},
\]
where $\lambda > L$ will be chosen later. The space $(X, \norm{\cdot}_\lambda)$
is a closed subset of $C(I, \R^n)$ with the $\norm{\cdot}_\lambda$ norm, hence
it is a complete metric space.

\medskip
\textbf{Step 2: The Picard operator maps $X$ to $X$.}
Define $T\colon X \to C(I, \R^n)$ by
\[
  (T\varphi)(t) = y_0 + \int_{t_0}^{t} f\bigl(s, \varphi(s)\bigr)\,\dd s.
\]
For $\varphi \in X$ and $t \in I$, since $\bigl(s, \varphi(s)\bigr) \in R$ we have
$\norm{f(s,\varphi(s))} \le M$, so
\[
  \norm{(T\varphi)(t) - y_0}
  \le \int_{t_0}^{t} \norm{f(s, \varphi(s))}\,\dd s
  \le M \abs{t - t_0}
  \le M\alpha \le b.
\]
Thus $T\varphi \in X$.

\medskip
\textbf{Step 3: $T$ is a contraction.}
For $\varphi, \psi \in X$,
\begin{align*}
  \norm{(T\varphi)(t) - (T\psi)(t)}
  &\le \biggl|\int_{t_0}^{t}
       \norm{f\bigl(s,\varphi(s)\bigr) - f\bigl(s,\psi(s)\bigr)}\,\dd s\biggr| \\
  &\le L \biggl|\int_{t_0}^{t} \norm{\varphi(s) - \psi(s)}\,\dd s\biggr| \\
  &= L \biggl|\int_{t_0}^{t} e^{\lambda\abs{s-t_0}}
     e^{-\lambda\abs{s-t_0}} \norm{\varphi(s) - \psi(s)}\,\dd s\biggr| \\
  &\le L\,\norm{\varphi - \psi}_\lambda
     \biggl|\int_{t_0}^{t} e^{\lambda\abs{s-t_0}}\,\dd s\biggr| \\
  &\le L\,\norm{\varphi - \psi}_\lambda \cdot
     \frac{e^{\lambda\abs{t-t_0}}}{\lambda}.
\end{align*}
Multiplying both sides by $e^{-\lambda\abs{t-t_0}}$ and taking the supremum
over $t \in I$,
\[
  \norm{T\varphi - T\psi}_\lambda
  \le \frac{L}{\lambda}\,\norm{\varphi - \psi}_\lambda.
\]
Choosing $\lambda > L$ makes $q = L/\lambda < 1$, so $T$ is a contraction on
$(X, \norm{\cdot}_\lambda)$.

\medskip
\textbf{Step 4: Conclusion.}
By the Banach fixed point theorem, $T$ has a unique fixed point $y \in X$, i.e.,
\[
  y(t) = y_0 + \int_{t_0}^t f\bigl(s, y(s)\bigr)\,\dd s.
\]
By Proposition~\ref{prop:integral-eq}, $y$ is the unique solution of the IVP on
$I$. Since $f$ is continuous and $y$ satisfies the integral equation, $y$ is in
fact $C^1$.
\end{proof}

\begin{remark}\label{rem:standard-proof}
Many textbooks prove the Picard--Lindel\"of theorem using the standard supremum
norm $\norm{\cdot}_\infty$ and the factor $\alpha^n L^n / n!$ to control the
$n$-th iterate directly. The weighted-norm approach above avoids induction and
gives the contraction property in one step.
\end{remark}

% ────────────────────────────────────────────────
\section{Peano's Existence Theorem}\label{sec:peano}
\index{Peano existence theorem}

If we drop the Lipschitz condition and only assume continuity, we lose
uniqueness but retain existence.

\begin{theorem}[Peano]\label{thm:peano}
\index{Peano existence theorem!statement}
Let $f\colon D \to \R^n$ be continuous on an open set
$D \subset \R \times \R^n$ with $(t_0, y_0) \in D$. Then the IVP
$y' = f(t,y)$, $y(t_0) = y_0$ has at least one solution on some interval
$[t_0 - \alpha, \, t_0 + \alpha]$ with $\alpha > 0$.
\end{theorem}

\begin{remark}\label{rem:peano-no-proof}
The proof uses the Arzel\`a--Ascoli theorem and is beyond our scope here.
The crucial point is that \emph{continuity alone guarantees existence but not
uniqueness}, as Example~\ref{ex:non-unique-y23} demonstrates.
\end{remark}

% ────────────────────────────────────────────────
\section{Maximal Solutions and Blow-Up}\label{sec:maximal-solutions}
\index{maximal solution}
\index{blow-up}

\begin{definition}[Maximal solution]\label{def:maximal-solution}
A solution $y\colon (\omega_-, \omega_+) \to \R^n$ is called \textbf{maximal}
(or \textbf{non-extendable}) if it cannot be extended to a solution on any
strictly larger interval. The interval $(\omega_-, \omega_+)$ is the
\textbf{maximal interval of existence}.
\end{definition}

\begin{theorem}[Blow-up alternative]\label{thm:blowup-alternative}
\index{blow-up!alternative}
Let $f$ be continuous and locally Lipschitz in $y$ on an open set
$D \subset \R \times \R^n$. Let $y\colon (\omega_-, \omega_+) \to \R^n$ be
the unique maximal solution of $y' = f(t,y)$, $y(t_0) = y_0$. Then either
$\omega_+ = +\infty$, or $\bigl(t, y(t)\bigr)$ leaves every compact subset
of $D$ as $t \to \omega_+^-$ (and similarly for $\omega_-$).
\end{theorem}

\begin{remark}\label{rem:blowup-cases}
In practice, ``leaving every compact set'' means either:
\begin{enumerate}[label=(\roman*)]
  \item $\norm{y(t)} \to \infty$ as $t \to \omega_+^-$ (blow-up in finite time), or
  \item $\bigl(t, y(t)\bigr)$ approaches the boundary of $D$.
\end{enumerate}
For $D = \R \times \R^n$, only case (i) is possible for finite $\omega_+$.
\end{remark}

\begin{example}[Blow-up revisited]\label{ex:blowup-revisited}
For $y' = y^2$, $y(0) = y_0 > 0$: the solution $y(t) = (y_0^{-1} - t)^{-1}$
has $\omega_+ = 1/y_0$. As $t \to (1/y_0)^-$, we get $y(t) \to +\infty$.
The larger the initial condition, the sooner the blow-up.
\end{example}

% ────────────────────────────────────────────────
\section{Gr\"onwall's Lemma}\label{sec:gronwall}
\index{Gr\"onwall's lemma}

Gr\"onwall's lemma is one of the most useful tools in the theory of ODEs.
It turns a differential or integral inequality into a pointwise bound.

\begin{lemma}[Gr\"onwall, integral form]\label{lem:gronwall}
\index{Gr\"onwall's lemma!integral form}
Let $u, \beta\colon [a,b] \to \R$ be continuous and non-negative, and let
$\alpha \ge 0$ be a constant. If
\begin{equation}\label{eq:gronwall-hyp}
  u(t) \le \alpha + \int_a^t \beta(s)\,u(s)\,\dd s
  \qquad \text{for all } t \in [a,b],
\end{equation}
then
\begin{equation}\label{eq:gronwall-concl}
  u(t) \le \alpha \exp\!\biggl(\int_a^t \beta(s)\,\dd s\biggr)
  \qquad \text{for all } t \in [a,b].
\end{equation}
\end{lemma}

\begin{proof}
Define $U(t) = \alpha + \displaystyle\int_a^t \beta(s)\,u(s)\,\dd s$, so that
$U(a) = \alpha$ and $u(t) \le U(t)$. Then
\[
  U'(t) = \beta(t)\,u(t) \le \beta(t)\,U(t).
\]
If $\alpha > 0$, then $U(t) > 0$ on $[a,b]$, and we may write
\[
  \frac{U'(t)}{U(t)} \le \beta(t).
\]
Integrating from $a$ to $t$:
\[
  \ln U(t) - \ln U(a) \le \int_a^t \beta(s)\,\dd s,
\]
so $U(t) \le U(a) \exp\!\bigl(\int_a^t \beta(s)\,\dd s\bigr)
= \alpha \exp\!\bigl(\int_a^t \beta(s)\,\dd s\bigr)$.
Since $u(t) \le U(t)$, the result follows.

For $\alpha = 0$, apply the result with $\alpha$ replaced by $\varepsilon > 0$
(the hypothesis still holds since $u(t) \le \varepsilon + \int_a^t \beta(s)\,u(s)\,\dd s$), obtaining
$u(t) \le \varepsilon\exp\!\bigl(\int_a^t \beta(s)\,\dd s\bigr)$. Letting
$\varepsilon \to 0$ yields $u(t) \le 0$, hence $u(t) = 0$.
\end{proof}

\begin{corollary}[Gr\"onwall, constant coefficient]\label{cor:gronwall-const}
\index{Gr\"onwall's lemma!constant coefficient}
If $u(t) \le \alpha + L\displaystyle\int_a^t u(s)\,\dd s$ for all $t \in [a,b]$
with $\alpha \ge 0$ and $L \ge 0$, then
\[
  u(t) \le \alpha\,e^{L(t-a)}.
\]
In particular, if $\alpha = 0$, then $u \equiv 0$.
\end{corollary}

% ────────────────────────────────────────────────
\section{Continuous Dependence on Initial Data}\label{sec:continuous-dependence}
\index{continuous dependence on initial data}

Solutions of an ODE are not just unique under a Lipschitz condition; they also
depend continuously on the initial values.

\begin{theorem}[Continuous dependence]\label{thm:continuous-dependence}
Let $f$ be continuous and Lipschitz in $y$ with constant $L$ on
$D \subset \R \times \R^n$. Let $y$ and $\tilde{y}$ be solutions of
\[
  y' = f(t,y),\quad y(t_0)=y_0
  \qquad\text{and}\qquad
  \tilde{y}' = f(t,\tilde{y}),\quad \tilde{y}(t_0) = \tilde{y}_0,
\]
on a common interval $[t_0, t_0 + T]$. Then
\begin{equation}\label{eq:cont-dep}
  \norm{y(t) - \tilde{y}(t)} \le \norm{y_0 - \tilde{y}_0}\,e^{L(t-t_0)}
  \qquad \text{for all } t \in [t_0, t_0 + T].
\end{equation}
\end{theorem}

\begin{proof}
Let $w(t) = \norm{y(t) - \tilde{y}(t)}$. From the integral formulation,
\[
  w(t) \le \norm{y_0 - \tilde{y}_0}
  + \int_{t_0}^t \norm{f(s,y(s)) - f(s,\tilde{y}(s))}\,\dd s
  \le \norm{y_0 - \tilde{y}_0} + L \int_{t_0}^t w(s)\,\dd s.
\]
By Gr\"onwall's lemma (Corollary~\ref{cor:gronwall-const}) with
$\alpha = \norm{y_0 - \tilde{y}_0}$, we get the result.
\end{proof}

\begin{remark}\label{rem:exponential-sensitivity}
\index{sensitivity!exponential}
The bound~\eqref{eq:cont-dep} shows that perturbations in initial data grow
at most exponentially. This exponential sensitivity is sharp in general (think
of $y' = Ly$). For chaotic systems the effective Lipschitz constant is large,
explaining rapid divergence of nearby trajectories.
\end{remark}

% ────────────────────────────────────────────────
\section{Further Counterexamples}\label{sec:counterexamples}
\index{counterexamples!existence and uniqueness}

\begin{example}[No global existence for $y'=1+y^2$]\label{ex:no-global}
\index{blow-up!$y'=1+y^2$}
The IVP $y' = 1 + y^2$, $y(0) = 0$ has the solution $y(t) = \tan t$, which
blows up at $t = \pi/2$. Here $f(t,y) = 1+y^2$ is locally Lipschitz but
not globally Lipschitz, and no global solution exists.
\end{example}

\begin{example}[Global existence with linear growth]\label{ex:global-linear}
\index{global existence}
If $\norm{f(t,y)} \le A + B\norm{y}$ for constants $A, B \ge 0$ and all
$(t,y) \in \R \times \R^n$, then every maximal solution exists for all time.
Indeed, Gr\"onwall's lemma shows that $\norm{y(t)}$ grows at most exponentially,
preventing finite-time blow-up.
\end{example}

\begin{exercise}\label{exer:linear-growth-global}
Prove the claim of Example~\ref{ex:global-linear} in detail.
\emph{Hint:} Apply Gr\"onwall's lemma to $u(t) = \norm{y(t)}$ and use the
blow-up alternative (Theorem~\ref{thm:blowup-alternative}).
\end{exercise}

% ────────────────────────────────────────────────
\section{Exercises}\label{sec:ch3-exercises}

\begin{exercise}\label{exer:picard-sinusoidal}
\index{Picard iteration!exercise}
Consider $y' = \cos y$, $y(0) = 0$. Compute the first three Picard iterates
$y_1, y_2, y_3$.
\end{exercise}

\begin{exercise}\label{exer:picard-2d}
For the system $x' = y$, $y' = -x$ with $x(0) = 1$, $y(0) = 0$, compute the
Picard iterates $(\mathbf{y}_0, \mathbf{y}_1, \mathbf{y}_2, \mathbf{y}_3)$
and identify the pattern. What is the exact solution?
\end{exercise}

\begin{exercise}\label{exer:lipschitz-verify}
\index{Lipschitz condition!exercises}
Determine whether the following functions are Lipschitz, locally Lipschitz, or
neither in $y$ on the given domain:
\begin{enumerate}[label=(\alph*)]
  \item $f(t,y) = e^{-t}\sin(y)$ on $\R \times \R$.
  \item $f(t,y) = \sqrt{\abs{y}}$ on $\R \times \R$.
  \item $f(t,y) = t^2 y^3$ on $\R \times \R$.
  \item $f(t,y) = \dfrac{y}{1+y^2}$ on $\R \times \R$.
\end{enumerate}
\end{exercise}

\begin{exercise}\label{exer:gronwall-application}
Let $y$ solve $y' = f(t,y)$ on $[0,T]$ with $\norm{f(t,y)} \le 3\norm{y} + 2$.
Show that $\norm{y(t)} \le (\norm{y(0)} + \tfrac{2}{3})\,e^{3t} - \tfrac{2}{3}$.
\end{exercise}

\begin{exercise}\label{exer:uniqueness-failure}
Show that $y' = 3y^{2/3}$, $y(0) = 0$ has a two-parameter family of solutions.
Find them explicitly.
\end{exercise}

\begin{exercise}\label{exer:blow-up-time}
Find the blow-up time for $y' = y^p$ with $y(0) = y_0 > 0$ and $p > 1$.
Express it in terms of $p$ and $y_0$.
\end{exercise}

\begin{exercise}\label{exer:gronwall-sharpen}
\index{Gr\"onwall's lemma!exercise}
Let $u \ge 0$ satisfy $u(t) \le 3 + 2\displaystyle\int_0^t s\,u(s)\,\dd s$
for $t \in [0, T]$. Show that $u(t) \le 3 e^{t^2}$.
\end{exercise}

\begin{exercise}\label{exer:cont-dep-quantitative}
Two solutions of $y' = -y + \sin t$ start at $y_0 = 1$ and
$\tilde{y}_0 = 1.01$. Give an upper bound on $\abs{y(t) - \tilde{y}(t)}$
for $t \ge 0$ using continuous dependence. Can you improve the bound by
exploiting the sign of the Lipschitz constant?
\end{exercise}

% ────────────────────────────────────────────────
\section{Chapter Summary}\label{sec:ch3-summary}

\begin{itemize}
  \item The \textbf{Picard--Lindel\"of theorem} guarantees local existence and
    uniqueness of solutions when $f$ is continuous and Lipschitz in $y$.
  \item The proof constructs the solution as the limit of \textbf{Picard
    iterates} and uses the Banach fixed point theorem.
  \item \textbf{Peano's theorem}: continuity alone suffices for existence, but
    uniqueness may fail (e.g., $y' = y^{2/3}$).
  \item \textbf{Maximal solutions} either exist for all time or leave every
    compact subset of the domain (blow-up alternative).
  \item \textbf{Gr\"onwall's lemma} converts integral inequalities into
    pointwise exponential bounds, and is the key tool for proving uniqueness,
    continuous dependence, and stability estimates.
  \item Solutions depend \textbf{continuously} on initial data, with at most
    exponential amplification of perturbations.
\end{itemize}


% ══════════════════════════════════════════════════════════════
%  CHAPTER 4 — SECOND-ORDER LINEAR ODEs WITH CONSTANT COEFFICIENTS
% ══════════════════════════════════════════════════════════════
\chapter{Second-Order Linear ODEs with Constant Coefficients}\label{ch:second-order}
\index{second-order ODE}
\index{linear ODE!second-order}

Second-order linear ODEs with constant coefficients form one of the most
important classes of differential equations in both pure and applied
mathematics. They model mechanical vibrations, electrical circuits, and
countless other phenomena. In this chapter we develop the complete theory:
characteristic equation, fundamental systems, the Wronskian, particular
solutions by undetermined coefficients, resonance, and the Euler--Cauchy
equation.

% ────────────────────────────────────────────────
\section{Physical Motivation}\label{sec:physical-motivation}
\index{spring-mass system}
\index{RLC circuit}

\subsection{Spring-mass system}

Consider a mass $m$ attached to a spring with stiffness $k > 0$, subject to
a damping force proportional to velocity (coefficient $c \ge 0$) and an
external driving force $F(t)$. Newton's second law gives
\begin{equation}\label{eq:spring-mass}
  m\,y'' + c\,y' + k\,y = F(t),
\end{equation}
where $y(t)$ denotes the displacement from equilibrium.\index{damping}

\subsection{RLC circuit}

An RLC series circuit with inductance $L$, resistance $R$, and capacitance $C$
driven by a voltage source $V(t)$ satisfies
\begin{equation}\label{eq:rlc}
  L\,q'' + R\,q' + \frac{1}{C}\,q = V(t),
\end{equation}
where $q(t)$ is the charge on the capacitor. The analogy with~\eqref{eq:spring-mass}
is exact: $L \leftrightarrow m$, $R \leftrightarrow c$,
$1/C \leftrightarrow k$.

% ────────────────────────────────────────────────
\section{The Homogeneous Equation}\label{sec:homogeneous-2nd}
\index{homogeneous equation!second-order}

We study
\begin{equation}\label{eq:homo-2nd}
  ay'' + by' + cy = 0,
\end{equation}
where $a, b, c \in \R$ with $a \ne 0$.

\subsection{The characteristic equation}\label{subsec:char-eq}
\index{characteristic equation}

Seeking solutions of the form $y(t) = e^{rt}$ and substituting into
\eqref{eq:homo-2nd} gives
\[
  a r^2 e^{rt} + b r e^{rt} + c e^{rt} = 0
  \quad\Longrightarrow\quad
  ar^2 + br + c = 0.
\]

\begin{definition}[Characteristic equation]\label{def:char-eq}
\index{characteristic equation!definition}
The \textbf{characteristic equation} (or \textbf{auxiliary equation}) of
$ay'' + by' + cy = 0$ is the quadratic
\begin{equation}\label{eq:char-poly}
  ar^2 + br + c = 0,
\end{equation}
with discriminant $\Delta = b^2 - 4ac$.\index{discriminant}
\end{definition}

The nature of the roots determines the form of the general solution. We treat
the three cases separately.

\subsection{Case 1: Two distinct real roots ($\Delta > 0$)}\label{subsec:case1}
\index{characteristic equation!distinct real roots}

\begin{theorem}[Distinct real roots]\label{thm:distinct-real}
If $\Delta > 0$, the characteristic equation has two distinct real roots
$r_1 \ne r_2$, and the general solution of \eqref{eq:homo-2nd} is
\begin{equation}\label{eq:gen-sol-case1}
  y(t) = C_1 e^{r_1 t} + C_2 e^{r_2 t},
\end{equation}
where $C_1, C_2 \in \R$ are arbitrary constants.
\end{theorem}

\begin{proof}
Both $y_1(t) = e^{r_1 t}$ and $y_2(t) = e^{r_2 t}$ are solutions by
construction. Their Wronskian (see Section~\ref{sec:wronskian}) is
\[
  W(t) = \begin{vmatrix} e^{r_1 t} & e^{r_2 t} \\
  r_1 e^{r_1 t} & r_2 e^{r_2 t} \end{vmatrix}
  = (r_2 - r_1)\,e^{(r_1+r_2)t} \ne 0,
\]
since $r_1 \ne r_2$. Hence $\{y_1, y_2\}$ is a fundamental system, and
every solution is a linear combination of $y_1$ and $y_2$.
\end{proof}

\begin{example}[Overdamped spring]\label{ex:overdamped}
\index{overdamped system}
Solve $y'' + 5y' + 6y = 0$, with $y(0) = 2$, $y'(0) = -3$.

The characteristic equation $r^2 + 5r + 6 = (r+2)(r+3) = 0$ gives
$r_1 = -2$, $r_2 = -3$. General solution:
$y(t) = C_1 e^{-2t} + C_2 e^{-3t}$.

Applying initial conditions:
\[
  C_1 + C_2 = 2, \qquad -2C_1 - 3C_2 = -3.
\]
Solving: $C_1 = 3$, $C_2 = -1$, so $y(t) = 3e^{-2t} - e^{-3t}$.
\end{example}

\subsection{Case 2: Repeated real root ($\Delta = 0$)}\label{subsec:case2}
\index{characteristic equation!repeated root}

\begin{theorem}[Repeated root]\label{thm:repeated-root}
If $\Delta = 0$, the characteristic equation has a double root
$r = -b/(2a)$, and the general solution is
\begin{equation}\label{eq:gen-sol-case2}
  y(t) = (C_1 + C_2\,t)\,e^{rt}.
\end{equation}
\end{theorem}

\begin{proof}
We already know $y_1(t) = e^{rt}$ is a solution. We seek a second,
linearly independent solution. Try $y_2(t) = t\,e^{rt}$. Then
\[
  y_2' = e^{rt}(1 + rt), \qquad
  y_2'' = e^{rt}(2r + r^2 t).
\]
Substituting into $ay'' + by' + cy$:
\begin{align*}
  a(2r + r^2 t)e^{rt} + b(1+rt)e^{rt} + c\,t\,e^{rt}
  &= e^{rt}\bigl[(ar^2+br+c)t + 2ar + b\bigr] \\
  &= e^{rt}[0 \cdot t + 2ar + b].
\end{align*}
Since $r = -b/(2a)$, we have $2ar + b = 0$. Thus $y_2 = te^{rt}$ is indeed
a solution. Its Wronskian with $y_1$ is
$W(t) = e^{2rt} \ne 0$,
so $\{e^{rt}, te^{rt}\}$ is a fundamental system.
\end{proof}

\begin{example}[Critically damped system]\label{ex:critically-damped}
\index{critically damped system}
Solve $y'' + 4y' + 4y = 0$, with $y(0) = 1$, $y'(0) = 0$.

Characteristic equation: $r^2 + 4r + 4 = (r+2)^2 = 0$, giving $r = -2$.
General solution: $y(t) = (C_1 + C_2 t)e^{-2t}$.

From $y(0) = C_1 = 1$ and $y'(0) = C_2 - 2C_1 = 0$ we get $C_2 = 2$.
Thus $y(t) = (1 + 2t)\,e^{-2t}$.
\end{example}

\subsection{Case 3: Complex conjugate roots ($\Delta < 0$)}\label{subsec:case3}
\index{characteristic equation!complex roots}

\begin{theorem}[Complex conjugate roots]\label{thm:complex-roots}
If $\Delta < 0$, the characteristic equation has conjugate complex roots
$r = \alpha \pm i\beta$ where $\alpha = -b/(2a)$ and
$\beta = \sqrt{4ac - b^2}/(2a) > 0$. The general solution (real-valued) is
\begin{equation}\label{eq:gen-sol-case3}
  y(t) = e^{\alpha t}\bigl(C_1 \cos\beta t + C_2 \sin\beta t\bigr).
\end{equation}
\end{theorem}

\begin{proof}
The complex-valued functions $e^{(\alpha+i\beta)t}$ and
$e^{(\alpha-i\beta)t}$ are solutions. Taking real and imaginary parts of
$e^{(\alpha+i\beta)t} = e^{\alpha t}(\cos\beta t + i\sin\beta t)$ gives two
real solutions:
\[
  y_1(t) = e^{\alpha t}\cos\beta t, \qquad
  y_2(t) = e^{\alpha t}\sin\beta t.
\]
Their Wronskian is
\[
  W(t) = e^{2\alpha t}
  \begin{vmatrix}
    \cos\beta t & \sin\beta t \\
    \alpha\cos\beta t - \beta\sin\beta t &
    \alpha\sin\beta t + \beta\cos\beta t
  \end{vmatrix}
  = \beta\,e^{2\alpha t} \ne 0.
\]
Hence $\{y_1, y_2\}$ is a fundamental system.
\end{proof}

\begin{example}[Underdamped oscillation]\label{ex:underdamped}
\index{underdamped system}
Solve $y'' + 2y' + 5y = 0$, with $y(0) = 1$, $y'(0) = -1$.

Characteristic equation: $r^2 + 2r + 5 = 0$, giving
$r = -1 \pm 2i$. So $\alpha = -1$, $\beta = 2$.
General solution: $y(t) = e^{-t}(C_1\cos 2t + C_2\sin 2t)$.

From $y(0) = C_1 = 1$ and $y'(0) = -C_1 + 2C_2 = -1$, we get $C_2 = 0$.
Thus $y(t) = e^{-t}\cos 2t$.
\end{example}

\begin{center}
\begin{tikzpicture}
  \begin{axis}[
    width=12cm, height=7cm,
    xlabel={$t$}, ylabel={$y(t)$},
    xmin=0, xmax=6, ymin=-1.2, ymax=2.5,
    axis lines=middle,
    title={Three damping regimes},
    legend pos=north east,
    samples=200, thick,
  ]
    % Overdamped
    \addplot[blue, domain=0:6] {3*exp(-2*x) - exp(-3*x)};
    \addlegendentry{Overdamped: $3e^{-2t}-e^{-3t}$}
    % Critically damped
    \addplot[red, domain=0:6] {(1+2*x)*exp(-2*x)};
    \addlegendentry{Critically damped: $(1+2t)e^{-2t}$}
    % Underdamped
    \addplot[green!60!black, domain=0:6] {exp(-x)*cos(deg(2*x))};
    \addlegendentry{Underdamped: $e^{-t}\cos 2t$}
    % Envelope for underdamped
    \addplot[green!60!black, dashed, thin, domain=0:6] {exp(-x)};
    \addplot[green!60!black, dashed, thin, domain=0:6] {-exp(-x)};
  \end{axis}
\end{tikzpicture}
\end{center}

% ────────────────────────────────────────────────
\section{Fundamental Solutions and the Wronskian}\label{sec:wronskian}
\index{Wronskian}
\index{fundamental system}

\begin{definition}[Wronskian]\label{def:wronskian}
Given two differentiable functions $y_1, y_2\colon I \to \R$, their
\textbf{Wronskian} is
\begin{equation}\label{eq:wronskian-def}
  W(y_1, y_2)(t) = \begin{vmatrix} y_1(t) & y_2(t) \\
  y_1'(t) & y_2'(t) \end{vmatrix}
  = y_1(t)\,y_2'(t) - y_2(t)\,y_1'(t).
\end{equation}
\end{definition}

\begin{theorem}[Abel's identity]\label{thm:abel}
\index{Abel's identity}
If $y_1$ and $y_2$ are solutions of $ay'' + by' + cy = 0$ on an interval $I$,
then
\begin{equation}\label{eq:abel}
  W(t) = W(t_0)\,\exp\!\biggl(-\frac{b}{a}(t - t_0)\biggr)
\end{equation}
for any $t_0 \in I$. In particular, $W$ is either identically zero or never
zero on $I$.
\end{theorem}

\begin{proof}
Differentiate the Wronskian:
\begin{align*}
  W' &= y_1 y_2'' - y_2 y_1''.
\end{align*}
From the ODE, $y_i'' = -(b/a)y_i' - (c/a)y_i$ for $i = 1, 2$. Thus
\begin{align*}
  W' &= y_1\bigl[-(b/a)y_2' - (c/a)y_2\bigr]
       - y_2\bigl[-(b/a)y_1' - (c/a)y_1\bigr] \\
     &= -(b/a)(y_1 y_2' - y_2 y_1') = -(b/a)\,W.
\end{align*}
This first-order linear ODE in $W$ has the solution
$W(t) = W(t_0)\,e^{-(b/a)(t-t_0)}$.
\end{proof}

\begin{proposition}[Fundamental system criterion]\label{prop:fund-system}
\index{fundamental system!criterion}
Two solutions $y_1, y_2$ of $ay'' + by' + cy = 0$ form a fundamental system
(i.e., every solution is a linear combination of $y_1$ and $y_2$) if and only
if $W(y_1, y_2)(t_0) \ne 0$ for some (equivalently, every) $t_0 \in I$.
\end{proposition}

% ────────────────────────────────────────────────
\section{Non-Homogeneous Equations: Undetermined Coefficients}\label{sec:undetermined-coeff}
\index{undetermined coefficients}
\index{non-homogeneous equation}

We now consider
\begin{equation}\label{eq:nonhomo-2nd}
  ay'' + by' + cy = g(t),
\end{equation}
where $g(t)$ is a given forcing function.

\begin{theorem}[Structure of general solution]\label{thm:general-nonhomo}
\index{non-homogeneous equation!general solution}
The general solution of~\eqref{eq:nonhomo-2nd} is
\[
  y(t) = y_h(t) + y_p(t),
\]
where $y_h$ is the general solution of the associated homogeneous equation
and $y_p$ is any particular solution of~\eqref{eq:nonhomo-2nd}.
\end{theorem}

\begin{proof}
If $y$ solves \eqref{eq:nonhomo-2nd}, then $y - y_p$ solves the homogeneous
equation: $a(y-y_p)'' + b(y-y_p)' + c(y-y_p) = g - g = 0$.
Hence $y - y_p = y_h$.
\end{proof}

\subsection{Method of undetermined coefficients}\label{subsec:method-undetermined}
\index{undetermined coefficients!method}

\begin{method}[Undetermined coefficients]\label{meth:undetermined}
When $g(t)$ is a polynomial, exponential, sine/cosine, or a product thereof,
we guess the form of $y_p$ and determine the unknown coefficients by
substitution.

The key rule: if the trial form overlaps with a solution of the homogeneous
equation, multiply by $t$ (or $t^2$) until the overlap is removed.
\end{method}

\begin{center}
\renewcommand{\arraystretch}{1.4}
\begin{tabular}{lll}
\toprule
\textbf{Forcing term $g(t)$} & \textbf{Trial $y_p(t)$} & \textbf{Modification} \\
\midrule
$P_n(t)$ (polynomial of degree $n$) & $A_n t^n + \dotsb + A_1 t + A_0$ & $\times\, t^s$ \\
$e^{\gamma t}$ & $A e^{\gamma t}$ & $\times\, t^s$ \\
$\cos\omega t$ or $\sin\omega t$ & $A\cos\omega t + B\sin\omega t$ & $\times\, t^s$ \\
$e^{\gamma t}\cos\omega t$ & $e^{\gamma t}(A\cos\omega t + B\sin\omega t)$ & $\times\, t^s$ \\
$t^n e^{\gamma t}$ & $(A_n t^n + \dotsb + A_0)\,e^{\gamma t}$ & $\times\, t^s$ \\
$t^n e^{\gamma t}\cos\omega t$ &
  $e^{\gamma t}(A_n t^n+\dotsb+A_0)\cos\omega t$ & $\times\, t^s$ \\
& $\quad{}+ e^{\gamma t}(B_n t^n+\dotsb+B_0)\sin\omega t$ & \\
\bottomrule
\end{tabular}
\captionof{table}{Trial particular solutions. Here $s$ is the multiplicity
of $\gamma$ (or $\gamma \pm i\omega$) as a root of the characteristic
equation ($s = 0, 1$, or $2$).}\label{tab:undetermined}
\index{undetermined coefficients!table}
\end{center}

\begin{example}[Polynomial forcing]\label{ex:poly-forcing}
Solve $y'' - 3y' + 2y = 4t^2$.

Characteristic equation: $r^2 - 3r + 2 = (r-1)(r-2) = 0$, so
$y_h = C_1 e^t + C_2 e^{2t}$.

Since $g(t) = 4t^2$ is a degree-2 polynomial and $0$ is not a root of the
characteristic equation ($s=0$), we try $y_p = At^2 + Bt + C$.
Then $y_p' = 2At + B$, $y_p'' = 2A$, so
\[
  2A - 3(2At+B) + 2(At^2+Bt+C) = 2At^2 + (2B-6A)t + (2A-3B+2C).
\]
Matching coefficients with $4t^2 + 0t + 0$:
\[
  2A = 4, \quad 2B - 6A = 0, \quad 2A - 3B + 2C = 0.
\]
So $A = 2$, $B = 6$, $C = 7$, and $y_p = 2t^2 + 6t + 7$.

General solution: $y = C_1 e^t + C_2 e^{2t} + 2t^2 + 6t + 7$.
\end{example}

\begin{example}[Exponential forcing with overlap]\label{ex:exp-overlap}
\index{undetermined coefficients!overlap}
Solve $y'' - 2y' + y = e^t$.

Characteristic equation: $(r-1)^2 = 0$, so $r = 1$ is a double root and
$y_h = (C_1 + C_2 t)e^t$.

Since $e^t$ corresponds to $\gamma = 1$, which is a root of multiplicity
$s = 2$, we try $y_p = At^2 e^t$. Then
\[
  y_p = At^2 e^t, \quad
  y_p' = A(2t + t^2)e^t, \quad
  y_p'' = A(2 + 4t + t^2)e^t.
\]
Substituting:
\[
  A(2+4t+t^2)e^t - 2A(2t+t^2)e^t + At^2 e^t = 2Ae^t.
\]
Setting $2A = 1$ gives $A = 1/2$. Thus $y_p = \dfrac{t^2}{2}\,e^t$.

General solution: $y = (C_1 + C_2 t)e^t + \dfrac{t^2}{2}\,e^t
= \bigl(C_1 + C_2 t + \tfrac{t^2}{2}\bigr)e^t$.
\end{example}

\begin{example}[Sinusoidal forcing]\label{ex:sin-forcing}
Solve $y'' + 4y = 3\sin t$.

Characteristic equation: $r^2 + 4 = 0$, so $r = \pm 2i$ and
$y_h = C_1\cos 2t + C_2\sin 2t$.

Since $\omega = 1 \ne 2$, there is no overlap ($s = 0$).
Try $y_p = A\cos t + B\sin t$.
Then $y_p'' = -A\cos t - B\sin t$, and
\[
  (-A\cos t - B\sin t) + 4(A\cos t + B\sin t) = 3A\cos t + 3B\sin t.
\]
Matching: $3A = 0$, $3B = 3$, so $A = 0$, $B = 1$ and $y_p = \sin t$.

General solution: $y = C_1\cos 2t + C_2\sin 2t + \sin t$.
\end{example}

% ────────────────────────────────────────────────
\section{Resonance}\label{sec:resonance}
\index{resonance}

Resonance occurs when the forcing frequency matches a natural frequency of
the system, causing the response to grow without bound (in the undamped case)
or to reach a maximum amplitude (with damping).

\begin{example}[Pure resonance]\label{ex:resonance}
\index{resonance!pure}
Consider $y'' + \omega_0^2 y = \cos\omega_0 t$ (no damping, forcing at
natural frequency $\omega_0$).

Characteristic roots: $r = \pm i\omega_0$, so
$y_h = C_1\cos\omega_0 t + C_2\sin\omega_0 t$.

Since $\cos\omega_0 t$ appears in $y_h$, we need $s = 1$: try
$y_p = t(A\cos\omega_0 t + B\sin\omega_0 t)$.
Computing:
\begin{align*}
  y_p'  &= A\cos\omega_0 t + B\sin\omega_0 t
           + t(-A\omega_0\sin\omega_0 t + B\omega_0\cos\omega_0 t), \\
  y_p'' &= -2A\omega_0\sin\omega_0 t + 2B\omega_0\cos\omega_0 t
           - t\omega_0^2(A\cos\omega_0 t + B\sin\omega_0 t).
\end{align*}
Substituting into $y'' + \omega_0^2 y$:
\[
  -2A\omega_0\sin\omega_0 t + 2B\omega_0\cos\omega_0 t = \cos\omega_0 t.
\]
So $A = 0$, $B = \dfrac{1}{2\omega_0}$, and
\[
  y_p(t) = \frac{t}{2\omega_0}\sin\omega_0 t.
\]
The factor of $t$ means the amplitude grows linearly---this is resonance.
\end{example}

\begin{center}
\begin{tikzpicture}
  \begin{axis}[
    width=12cm, height=6.5cm,
    xlabel={$t$}, ylabel={$y(t)$},
    xmin=0, xmax=30, ymin=-10, ymax=10,
    axis lines=middle,
    title={Resonance: $y'' + y = \cos t$, \ $y(0)=y'(0)=0$},
    samples=400, thick,
  ]
    % y_p = (t/2) sin(t)
    \addplot[blue, domain=0:30] {(x/2)*sin(deg(x))};
    % Envelope
    \addplot[red, dashed, domain=0:30] {x/2};
    \addplot[red, dashed, domain=0:30] {-x/2};
    \node[red] at (axis cs:28,8) {$\pm t/2$};
  \end{axis}
\end{tikzpicture}
\end{center}

\begin{example}[Near-resonance and beats]\label{ex:beats}
\index{beats}
Consider $y'' + \omega_0^2 y = \cos\omega t$ with $\omega \ne \omega_0$ but
$\omega \approx \omega_0$, and $y(0) = y'(0) = 0$. The solution is
\[
  y(t) = \frac{1}{\omega_0^2 - \omega^2}(\cos\omega t - \cos\omega_0 t)
  = \frac{2}{\omega_0^2 - \omega^2}
    \sin\!\biggl(\frac{\omega_0 - \omega}{2}\,t\biggr)
    \sin\!\biggl(\frac{\omega_0 + \omega}{2}\,t\biggr).
\]
When $\omega_0 \approx \omega$, the first sine factor produces a slowly varying
envelope (the ``beat'' phenomenon), while the second gives rapid oscillations.
\end{example}

% ────────────────────────────────────────────────
\section{The Euler--Cauchy Equation}\label{sec:euler-cauchy}
\index{Euler--Cauchy equation}

\begin{definition}[Euler--Cauchy equation]\label{def:euler-cauchy}
An \textbf{Euler--Cauchy equation} (or \textbf{equidimensional equation}) is
an ODE of the form
\begin{equation}\label{eq:euler-cauchy}
  at^2 y'' + bt y' + cy = 0, \qquad t > 0,
\end{equation}
where $a, b, c$ are real constants with $a \ne 0$.
\end{definition}

\begin{method}[Solving Euler--Cauchy equations]\label{meth:euler-cauchy}
\index{Euler--Cauchy equation!method}
Seek solutions of the form $y = t^r$ (for $t > 0$). Then $y' = rt^{r-1}$,
$y'' = r(r-1)t^{r-2}$, and substitution gives
\[
  a r(r-1) + br + c = 0 \qquad\text{(indicial equation).}
\]
\index{indicial equation}
\end{method}

The three cases mirror those of constant-coefficient equations:

\begin{theorem}[Euler--Cauchy: general solution]\label{thm:euler-cauchy-sol}
\index{Euler--Cauchy equation!general solution}
Let $r_1, r_2$ be the roots of $ar(r-1) + br + c = 0$.
\begin{enumerate}[label=(\roman*)]
  \item \textbf{Distinct real roots} ($r_1 \ne r_2$ real):
    $y(t) = C_1 t^{r_1} + C_2 t^{r_2}$.
  \item \textbf{Repeated root} ($r_1 = r_2 = r$):
    $y(t) = (C_1 + C_2\ln t)\,t^r$.
  \item \textbf{Complex roots} ($r = \alpha \pm i\beta$, $\beta > 0$):
    $y(t) = t^\alpha\bigl(C_1\cos(\beta\ln t) + C_2\sin(\beta\ln t)\bigr)$.
\end{enumerate}
\end{theorem}

\begin{proof}
Case (i) is immediate. For (ii), with $r$ a double root, $y_1 = t^r$ is one
solution; a second is found by reduction of order
(Section~\ref{sec:reduction-of-order}) or by the substitution $\tau = \ln t$,
giving $y_2 = t^r \ln t$. For (iii), write $t^r = t^{\alpha + i\beta}
= t^\alpha e^{i\beta\ln t} = t^\alpha(\cos(\beta\ln t) + i\sin(\beta\ln t))$,
and take real and imaginary parts.
\end{proof}

\begin{example}\label{ex:euler-cauchy-1}
Solve $t^2 y'' - 2y = 0$ for $t > 0$.

Indicial equation: $r(r-1) - 2 = r^2 - r - 2 = (r-2)(r+1) = 0$, so
$r_1 = 2$, $r_2 = -1$. General solution: $y(t) = C_1 t^2 + C_2 t^{-1}$.
\end{example}

\begin{example}\label{ex:euler-cauchy-2}
Solve $t^2 y'' + ty' + y = 0$ for $t > 0$.

Indicial equation: $r(r-1) + r + 1 = r^2 + 1 = 0$, so $r = \pm i$.
Thus $\alpha = 0$, $\beta = 1$, and
$y(t) = C_1\cos(\ln t) + C_2\sin(\ln t)$.
\end{example}

\begin{remark}\label{rem:euler-cauchy-substitution}
\index{Euler--Cauchy equation!substitution}
The substitution $\tau = \ln t$ (so $t = e^\tau$) transforms the Euler--Cauchy
equation~\eqref{eq:euler-cauchy} into a constant-coefficient equation in $\tau$.
This is often the quickest approach.
\end{remark}

% ────────────────────────────────────────────────
\section{Reduction of Order}\label{sec:reduction-of-order}
\index{reduction of order}

If one solution $y_1$ of a second-order linear ODE is known, a second
linearly independent solution can be found by a standard technique.

\begin{method}[Reduction of order]\label{meth:reduction-of-order}
Given $y'' + p(t)y' + q(t)y = 0$ with a known solution $y_1(t) \ne 0$,
set $y = v(t)\,y_1(t)$. Then $v$ satisfies a first-order ODE in $w = v'$:
\begin{equation}\label{eq:reduction-of-order}
  w' + \biggl(\frac{2y_1'}{y_1} + p(t)\biggr)w = 0,
\end{equation}
which is separable. Solving for $w$ and integrating gives $v$, hence
$y_2 = v\,y_1$.
\end{method}

\begin{proof}[Derivation]
Set $y = vy_1$. Then $y' = v'y_1 + vy_1'$,
$y'' = v''y_1 + 2v'y_1' + vy_1''$. Substituting:
\begin{align*}
  (v''y_1 + 2v'y_1' + vy_1'') + p(v'y_1 + vy_1') + q(vy_1) &= 0 \\
  v''y_1 + v'(2y_1' + py_1) + v\underbrace{(y_1'' + py_1' + qy_1)}_{=\,0} &= 0.
\end{align*}
Setting $w = v'$ yields \eqref{eq:reduction-of-order}.
\end{proof}

\begin{example}[Reduction of order]\label{ex:reduction-of-order}
Given that $y_1(t) = e^t$ solves $y'' - 2y' + y = 0$, find a second solution.

Here $p(t) = -2$. Set $w = v'$:
\[
  w' + \biggl(\frac{2e^t}{e^t} - 2\biggr)w = w' + (2 - 2)w = w' = 0.
\]
So $w = C$ (constant), $v = Ct$, and $y_2 = te^t$, confirming the result from
the repeated-root case.
\end{example}

\begin{example}[Reduction of order for Euler--Cauchy]\label{ex:reduction-euler}
Given that $y_1 = t$ solves $t^2 y'' - 2ty' + 2y = 0$ for $t > 0$, find $y_2$.

Rewrite in standard form: $y'' - \frac{2}{t}y' + \frac{2}{t^2}y = 0$,
so $p(t) = -2/t$. With $y_1 = t$, $y_1' = 1$:
\[
  w' + \biggl(\frac{2}{t} - \frac{2}{t}\biggr)w = 0
  \quad\Longrightarrow\quad w = C.
\]
Thus $v = Ct$, and $y_2 = t \cdot t = t^2$.
\end{example}

% ────────────────────────────────────────────────
\section{Worked Examples}\label{sec:ch4-worked-examples}

\begin{example}[Complete IVP]\label{ex:complete-ivp}
Solve $y'' + y = t\cos t$, with $y(0) = 0$, $y'(0) = 0$.

\textbf{Homogeneous solution.} $r^2 + 1 = 0$, $r = \pm i$, so
$y_h = C_1\cos t + C_2\sin t$.

\textbf{Particular solution.} The forcing $g(t) = t\cos t$ has the form
$t\,e^{0\cdot t}\cos t$, corresponding to $\gamma = 0$, $\omega = 1$,
$n = 1$. Since $0 + i\cdot 1 = i$ is a root of the characteristic equation
with multiplicity $s = 1$, we multiply by $t$:
\[
  y_p = t\bigl[(At + B)\cos t + (Ct + D)\sin t\bigr].
\]
After computing $y_p''$ and substituting (a lengthy but routine calculation),
one finds
\[
  y_p = \frac{t^2}{4}\sin t - \frac{t}{4}\cos t.
\]

\textbf{Initial conditions.}
$y(0) = C_1 + 0 = 0$, so $C_1 = 0$.
$y'(0) = C_2 + (-1/4) = 0$, so $C_2 = 1/4$.

\textbf{Solution:}
$y(t) = \dfrac{1}{4}\sin t + \dfrac{t^2}{4}\sin t - \dfrac{t}{4}\cos t
= \dfrac{(1+t^2)\sin t - t\cos t}{4}$.
\end{example}

\begin{example}[Exponential forcing with complex roots]\label{ex:exp-complex}
Solve $y'' + 2y' + 5y = 4e^{-t}\cos 2t$.

Characteristic equation: $r^2 + 2r + 5 = 0$, giving $r = -1 \pm 2i$.
So $y_h = e^{-t}(C_1\cos 2t + C_2\sin 2t)$.

The forcing is $g(t) = 4e^{-t}\cos 2t$, corresponding to $\gamma + i\omega = -1+2i$,
which is a root with multiplicity $s = 1$. Try
$y_p = te^{-t}(A\cos 2t + B\sin 2t)$.

After substitution and simplification:
$-4A\sin 2t + 4B\cos 2t = 4\cos 2t$,
giving $A = 0$, $B = 1$. So $y_p = te^{-t}\sin 2t$.

General solution: $y = e^{-t}(C_1\cos 2t + C_2\sin 2t) + te^{-t}\sin 2t$.
\end{example}

\begin{example}[Superposition]\label{ex:superposition}
\index{superposition principle}
Solve $y'' + 4y = 3\sin t + 5e^{2t}$.

By superposition, find particular solutions for each term separately.

For $g_1 = 3\sin t$: from Example~\ref{ex:sin-forcing}, $y_{p_1} = \sin t$.

For $g_2 = 5e^{2t}$: try $y_{p_2} = Ae^{2t}$.
Then $4Ae^{2t} + 4Ae^{2t} = 8Ae^{2t} = 5e^{2t}$, so $A = 5/8$.

General solution: $y = C_1\cos 2t + C_2\sin 2t + \sin t + \frac{5}{8}e^{2t}$.
\end{example}

% ────────────────────────────────────────────────
\section{Summary of Solution Methods}\label{sec:ch4-methods-table}

\begin{center}
\renewcommand{\arraystretch}{1.5}
\begin{tabular}{lll}
\toprule
\textbf{Roots of char.\ eq.} & \textbf{Form} & \textbf{General solution} \\
\midrule
$r_1, r_2 \in \R$, $r_1 \ne r_2$ & Distinct real &
  $C_1 e^{r_1 t} + C_2 e^{r_2 t}$ \\
$r_1 = r_2 = r \in \R$ & Repeated &
  $(C_1 + C_2 t)e^{rt}$ \\
$r = \alpha \pm i\beta$ & Complex &
  $e^{\alpha t}(C_1\cos\beta t + C_2\sin\beta t)$ \\
\bottomrule
\end{tabular}
\captionof{table}{Homogeneous solutions for $ay''+by'+cy=0$.}\label{tab:homo-solutions}
\end{center}

\begin{center}
\renewcommand{\arraystretch}{1.5}
\begin{tabular}{lcc}
\toprule
\textbf{Damping} & \textbf{Condition} & \textbf{Behavior} \\
\midrule
Overdamped & $b^2 > 4ac$ & Exponential decay, no oscillation \\
Critically damped & $b^2 = 4ac$ & Fastest non-oscillatory decay \\
Underdamped & $b^2 < 4ac$ & Oscillatory decay \\
Undamped & $b = 0$ & Perpetual oscillation \\
\bottomrule
\end{tabular}
\captionof{table}{Damping regimes for $ay''+by'+cy=0$ with $a,c>0$, $b \ge 0$.}\label{tab:damping}
\index{damping!regimes}
\end{center}

% ────────────────────────────────────────────────
\section{Exercises}\label{sec:ch4-exercises}

\begin{exercise}\label{exer:homo-solve}
Solve the following homogeneous equations. Find the general solution and classify
the damping regime.
\begin{enumerate}[label=(\alph*)]
  \item $y'' + 6y' + 8y = 0$
  \item $y'' + 6y' + 9y = 0$
  \item $y'' + 2y' + 10y = 0$
  \item $4y'' + y = 0$
\end{enumerate}
\end{exercise}

\begin{exercise}\label{exer:ivp-homo}
Solve the IVPs:
\begin{enumerate}[label=(\alph*)]
  \item $y'' - y = 0$, $y(0) = 3$, $y'(0) = -1$.
  \item $y'' + 4y' + 4y = 0$, $y(0) = 2$, $y'(0) = 1$.
  \item $y'' + 9y = 0$, $y(\pi/3) = 0$, $y'(\pi/3) = 3$.
\end{enumerate}
\end{exercise}

\begin{exercise}\label{exer:undetermined}
\index{undetermined coefficients!exercises}
Find the general solution using undetermined coefficients:
\begin{enumerate}[label=(\alph*)]
  \item $y'' + y' - 2y = 4e^{3t}$
  \item $y'' + 4y = 8\cos 2t$
  \item $y'' - 4y' + 4y = 6e^{2t}$
  \item $y'' + y = t^2 + 1$
  \item $y'' + 2y' + y = 3te^{-t}$
\end{enumerate}
\end{exercise}

\begin{exercise}\label{exer:resonance-exercise}
\index{resonance!exercise}
An undamped oscillator $y'' + 9y = A\cos 3t$ is driven at its natural frequency.
\begin{enumerate}[label=(\alph*)]
  \item Find the general solution with $y(0) = 0$, $y'(0) = 0$.
  \item Determine the amplitude of oscillation at time $t = 10$.
  \item Sketch the solution for $A = 1$ on $[0, 20]$.
\end{enumerate}
\end{exercise}

\begin{exercise}\label{exer:euler-cauchy-exercise}
\index{Euler--Cauchy equation!exercises}
Solve the Euler--Cauchy equations for $t > 0$:
\begin{enumerate}[label=(\alph*)]
  \item $t^2 y'' + 5ty' + 4y = 0$
  \item $t^2 y'' - ty' + y = 0$
  \item $t^2 y'' + ty' + 4y = 0$
\end{enumerate}
\end{exercise}

\begin{exercise}\label{exer:reduction-practice}
\index{reduction of order!exercise}
Use reduction of order to find a second solution, given $y_1$:
\begin{enumerate}[label=(\alph*)]
  \item $y'' - 4y' + 4y = 0$, $y_1 = e^{2t}$.
  \item $t^2 y'' - 3ty' + 4y = 0$ ($t > 0$), $y_1 = t^2$.
  \item $(1-t)y'' + ty' - y = 0$ ($0 < t < 1$), $y_1 = e^t$.
\end{enumerate}
\end{exercise}

\begin{exercise}\label{exer:wronskian-compute}
\index{Wronskian!exercise}
Compute the Wronskian and verify Abel's identity for the solutions
$y_1 = e^{-t}\cos 2t$ and $y_2 = e^{-t}\sin 2t$ of $y'' + 2y' + 5y = 0$.
\end{exercise}

\begin{exercise}\label{exer:application-spring}
\index{spring-mass system!exercise}
A $2\,\mathrm{kg}$ mass on a spring ($k = 50\,\mathrm{N/m}$) with damping
($c = 4\,\mathrm{N\cdot s/m}$) is released from rest at a displacement of
$0.1\,\mathrm{m}$.
\begin{enumerate}[label=(\alph*)]
  \item Write and solve the IVP.
  \item Is the system underdamped, critically damped, or overdamped?
  \item Find the quasi-frequency and quasi-period.
  \item Determine the time at which the displacement first returns to zero.
\end{enumerate}
\end{exercise}

\begin{exercise}\label{exer:beats-exercise}
\index{beats!exercise}
Solve $y'' + 25y = \cos 4t$ with $y(0) = 0$, $y'(0) = 0$.
Identify the beat frequency and the carrier frequency. Sketch the solution
for $t \in [0, 10\pi]$.
\end{exercise}

% ────────────────────────────────────────────────
\section{Chapter Summary}\label{sec:ch4-summary}

\begin{itemize}
  \item The general solution of $ay'' + by' + cy = 0$ is determined by the
    \textbf{characteristic equation} $ar^2 + br + c = 0$ and its three cases
    (distinct real, repeated, complex conjugate roots).
  \item The \textbf{Wronskian} determines linear independence of solutions.
    Abel's identity shows it is either always zero or never zero.
  \item For \textbf{non-homogeneous} equations $ay'' + by' + cy = g(t)$:
    the general solution is $y_h + y_p$, where $y_p$ can be found by
    undetermined coefficients when $g$ involves polynomials, exponentials,
    and trigonometric functions.
  \item \textbf{Resonance} occurs when the forcing frequency matches a natural
    frequency, causing unbounded growth (undamped) or maximal amplitude (damped).
  \item The \textbf{Euler--Cauchy equation} $at^2y'' + bty' + cy = 0$ is solved
    by the ansatz $y = t^r$, leading to an indicial equation.
  \item \textbf{Reduction of order}: knowing one solution $y_1$, a second
    independent solution can always be constructed.
\end{itemize}

% === END OF CHUNK ch3-4 ===
% === START OF CHUNK ch5-6 ===

%%%%%%%%%%%%%%%%%%%%%%%%%%%%%%%%%%%%%%%%%%%%%%%%%%%%%%%%%%%%%
%  CHAPTER 5 — LINEAR DIFFERENTIAL SYSTEMS
%%%%%%%%%%%%%%%%%%%%%%%%%%%%%%%%%%%%%%%%%%%%%%%%%%%%%%%%%%%%%
\chapter{Linear Differential Systems}\label{ch:linear-systems}
\index{linear differential system}

\section{Motivation: Coupled Dynamical Models}\label{sec:sys-motivation}

Many phenomena in physics, biology, and engineering are naturally described
not by a single differential equation but by a \emph{system} of coupled
equations.  We present two classical examples.

\begin{example}[Coupled oscillators]\label{ex:coupled-oscillators}
\index{coupled oscillators}
Consider two masses $m_1,m_2$ connected by springs with constants
$k_1,k_2,k_3$ arranged in series along a line.  Letting $x_1(t),x_2(t)$
denote the displacements from equilibrium, Newton's second law gives
\[
  m_1 x_1'' = -k_1 x_1 + k_2(x_2-x_1), \qquad
  m_2 x_2'' = -k_3 x_2 - k_2(x_2-x_1).
\]
Setting $X = (x_1, x_1', x_2, x_2')^T$ this becomes a first-order system
$X' = AX$ with a $4\times 4$ constant matrix $A$.
\end{example}

\begin{example}[Lotka--Volterra predator-prey model]\label{ex:lotka-volterra}
\index{Lotka--Volterra model}\index{predator-prey model}
Let $x(t)$ be the prey population and $y(t)$ the predator population.
The classical model reads
\[
  x' = \alpha x - \beta xy, \qquad
  y' = \delta xy - \gamma y,
\]
where $\alpha,\beta,\gamma,\delta>0$.  This is a \emph{nonlinear} system;
however, linearizing around the coexistence equilibrium
$(\gamma/\delta,\,\alpha/\beta)$ produces a linear system that governs the
local dynamics.
\end{example}

\medskip

These examples illustrate that first-order linear systems arise both directly
and as linearizations of nonlinear models.  The general theory developed in
this chapter applies to both settings.

%--------------------------------------------------------------------
\section{General Form and Basic Properties}\label{sec:sys-general-form}
%--------------------------------------------------------------------

\begin{definition}[Linear differential system]\label{def:linear-system}
\index{linear differential system!general form}
A \textbf{linear differential system} on an interval $I\subset\R$ is an
equation of the form
\begin{equation}\label{eq:sys-nonhom}
  X'(t) = A(t)\,X(t) + B(t),
\end{equation}
where $A:I\to\mathcal{M}_n(\R)$ and $B:I\to\R^n$ are continuous,
$X:I\to\R^n$ is the unknown.  When $B\equiv 0$ the system is
\textbf{homogeneous}\index{homogeneous system}:
\begin{equation}\label{eq:sys-hom}
  X'(t) = A(t)\,X(t).
\end{equation}
\end{definition}

\begin{theorem}[Existence and uniqueness for linear systems]\label{thm:sys-exist}
\index{existence and uniqueness!linear system}
Let $A$ and $B$ be continuous on an open interval $I$ and let
$t_0\in I$, $X_0\in\R^n$.  Then the initial value problem
\[
  X' = A(t)X + B(t), \qquad X(t_0) = X_0
\]
has a \emph{unique} solution defined on the \emph{entire} interval $I$.
\end{theorem}

\begin{proof}
The right-hand side $f(t,X)=A(t)X+B(t)$ is globally Lipschitz in $X$
(with Lipschitz constant $\norm{A(t)}$), so the Picard--Lindel\"of theorem
yields a unique local solution.  Since the Lipschitz constant is locally
bounded and the solution cannot blow up in finite time (Gr\"onwall's
inequality gives $\norm{X(t)}\le C e^{\int \norm{A(s)}\dd s}$), the solution
extends to all of $I$.
\end{proof}

\begin{proposition}[Superposition principle]\label{prop:superposition-sys}
\index{superposition principle!systems}
The set of solutions of the homogeneous system \eqref{eq:sys-hom} is a
vector space of dimension $n$.
\end{proposition}

\begin{proof}
Linearity of the equation implies that the solution set is a subspace of
$\mathcal{C}^1(I,\R^n)$.  To see that its dimension is $n$, consider the
linear map $\Phi:X\mapsto X(t_0)$ from the solution space to $\R^n$.
By Theorem~\ref{thm:sys-exist}, $\Phi$ is a bijection.
\end{proof}

%--------------------------------------------------------------------
\section{Homogeneous Systems with Constant Coefficients}\label{sec:sys-const}
\index{constant coefficient system}
%--------------------------------------------------------------------

Throughout this section, $A\in\mathcal{M}_n(\R)$ is a constant matrix
and we study
\begin{equation}\label{eq:sys-const-hom}
  X' = AX.
\end{equation}

\subsection{The Matrix Exponential}\label{subsec:matrix-exp}
\index{matrix exponential}

\begin{definition}[Matrix exponential]\label{def:matrix-exp}
For $A\in\mathcal{M}_n(\R)$ the \textbf{matrix exponential} is
\begin{equation}\label{eq:matrix-exp-def}
  e^{A} = \sum_{k=0}^{\infty} \frac{A^k}{k!}.
\end{equation}
The series converges absolutely in any matrix norm.
\end{definition}

\begin{theorem}[Properties of the matrix exponential]\label{thm:exp-properties}
\index{matrix exponential!properties}
Let $A,B\in\mathcal{M}_n(\R)$ and $t,s\in\R$.
\begin{enumerate}[label=(\roman*)]
  \item $e^{0} = I_n$.
  \item $e^{(t+s)A} = e^{tA}\,e^{sA}$.
  \item If $AB = BA$, then $e^{A+B} = e^{A}\,e^{B}$.
  \item $e^{A}$ is invertible with $(e^{A})^{-1} = e^{-A}$.
  \item $\dfrac{\dd}{\dd t}\,e^{tA} = A\,e^{tA} = e^{tA}\,A$.
  \item If $P$ is invertible, then $e^{PAP^{-1}} = P\,e^{A}\,P^{-1}$.
\end{enumerate}
\end{theorem}

\begin{proof}
Properties (i)--(iv) follow from the series definition and absolute
convergence.  For (v), differentiating term-by-term:
\[
  \frac{\dd}{\dd t}\sum_{k=0}^{\infty}\frac{(tA)^k}{k!}
  = \sum_{k=1}^{\infty}\frac{k\,t^{k-1}A^k}{k!}
  = A\sum_{k=1}^{\infty}\frac{(tA)^{k-1}}{(k-1)!}
  = A\,e^{tA}.
\]
For (vi), note $(PAP^{-1})^k = PA^kP^{-1}$ and sum.
\end{proof}

\begin{theorem}[Solution via matrix exponential]\label{thm:sol-matrix-exp}
\index{matrix exponential!solution of system}
The unique solution of $X'=AX$, $X(0)=X_0$ is
\begin{equation}\label{eq:sol-exp}
  X(t) = e^{tA}\,X_0.
\end{equation}
\end{theorem}

\begin{proof}
By property (v) of Theorem~\ref{thm:exp-properties},
$\frac{\dd}{\dd t}(e^{tA}X_0) = Ae^{tA}X_0 = A\,X(t)$.
At $t=0$, $X(0) = e^{0}X_0 = X_0$.
\end{proof}

\begin{remark}\label{rem:computing-exp}
Computing $e^{tA}$ directly from the series is rarely practical.
The standard strategies are:
\begin{itemize}
  \item Diagonalization: if $A=PDP^{-1}$ with
    $D=\diag(\lambda_1,\dots,\lambda_n)$, then
    $e^{tA}=P\,\diag(e^{t\lambda_1},\dots,e^{t\lambda_n})\,P^{-1}$.
  \item Jordan form: $A=PJP^{-1}$ gives $e^{tA}=Pe^{tJ}P^{-1}$.
  \item Cayley--Hamilton: express $e^{tA}$ as a polynomial of degree
    $\le n-1$ in $A$.
\end{itemize}
\end{remark}

\subsection{Distinct Real Eigenvalues}\label{subsec:distinct-real}
\index{distinct real eigenvalues}

\begin{method}[Distinct real eigenvalues]\label{meth:distinct-real}
Suppose $A\in\mathcal{M}_n(\R)$ has $n$ distinct real eigenvalues
$\lambda_1,\dots,\lambda_n$ with eigenvectors $v_1,\dots,v_n$.
The general solution of $X'=AX$ is
\begin{equation}\label{eq:sol-distinct-real}
  X(t) = c_1 e^{\lambda_1 t}v_1 + c_2 e^{\lambda_2 t}v_2
         + \cdots + c_n e^{\lambda_n t}v_n,
  \qquad c_i\in\R.
\end{equation}
\end{method}

\begin{example}[A $2\times 2$ system with distinct eigenvalues]\label{ex:2x2-distinct}
Solve $X' = \begin{pmatrix} 1 & 2 \\ 3 & 2 \end{pmatrix}X$.

The characteristic polynomial is $\lambda^2 - 3\lambda - 4 = (\lambda-4)(\lambda+1)$.
\begin{itemize}
  \item $\lambda_1 = 4$: eigenvector $v_1 = \begin{pmatrix}2\\3\end{pmatrix}$.
  \item $\lambda_2 = -1$: eigenvector $v_2 = \begin{pmatrix}1\\-1\end{pmatrix}$.
\end{itemize}
General solution:
\[
  X(t) = c_1 e^{4t}\begin{pmatrix}2\\3\end{pmatrix}
        + c_2 e^{-t}\begin{pmatrix}1\\-1\end{pmatrix}.
\]
\end{example}

\subsection{Complex Eigenvalues}\label{subsec:complex-eigenvalues}
\index{complex eigenvalues}

When $A\in\mathcal{M}_n(\R)$ has a complex eigenvalue
$\lambda=\alpha+i\beta$ (with $\beta\neq 0$) and eigenvector
$w = u + iv$ ($u,v\in\R^n$), the conjugate $\bar\lambda = \alpha-i\beta$
is also an eigenvalue with eigenvector $\bar w = u-iv$.

\begin{method}[Complex eigenvalues -- real form]\label{meth:complex-eigen}
\index{complex eigenvalues!real form}
The pair of conjugate eigenvalues $\alpha\pm i\beta$ contributes
two real-valued solutions:
\begin{equation}\label{eq:sol-complex-real}
\begin{aligned}
  X_1(t) &= e^{\alpha t}\bigl(\cos(\beta t)\,u - \sin(\beta t)\,v\bigr),\\
  X_2(t) &= e^{\alpha t}\bigl(\sin(\beta t)\,u + \cos(\beta t)\,v\bigr).
\end{aligned}
\end{equation}
\end{method}

\begin{example}[Complex eigenvalues]\label{ex:complex-eigenvalues}
Solve $X' = \begin{pmatrix} 0 & -2 \\ 2 & 0 \end{pmatrix}X$.

Eigenvalues: $\lambda = \pm 2i$.
For $\lambda = 2i$: $w = \begin{pmatrix}1\\-i\end{pmatrix}$, so
$u = \begin{pmatrix}1\\0\end{pmatrix}$,
$v = \begin{pmatrix}0\\-1\end{pmatrix}$.
General solution:
\[
  X(t) = c_1\begin{pmatrix}\cos 2t\\\sin 2t\end{pmatrix}
        + c_2\begin{pmatrix}-\sin 2t\\\cos 2t\end{pmatrix}.
\]
The trajectories are circles centered at the origin (a \emph{center}).
\end{example}

\subsection{Repeated Eigenvalues and Jordan Form}\label{subsec:jordan}
\index{Jordan form}\index{generalized eigenvector}\index{repeated eigenvalue}

When $A$ does not have $n$ linearly independent eigenvectors, we need
generalized eigenvectors.

\begin{definition}[Generalized eigenvector]\label{def:gen-eigenvector}
\index{generalized eigenvector}
A vector $w\in\R^n$ is a \textbf{generalized eigenvector} of $A$ of
rank $k$ for eigenvalue $\lambda$ if
\[
  (A-\lambda I)^k w = 0 \quad\text{and}\quad (A-\lambda I)^{k-1}w \neq 0.
\]
\end{definition}

\begin{method}[Jordan block contribution]\label{meth:jordan-block}
\index{Jordan block}
For a Jordan block of size $m$ associated to eigenvalue $\lambda$,
with generalized eigenvectors $w_1,\dots,w_m$ satisfying
$(A-\lambda I)w_1=0$, $(A-\lambda I)w_j=w_{j-1}$ for $j\ge 2$,
the corresponding $m$ linearly independent solutions are
\begin{equation}\label{eq:jordan-solutions}
  X_j(t) = e^{\lambda t}\sum_{i=0}^{j-1}\frac{t^i}{i!}\,w_{j-i},
  \qquad j=1,\dots,m.
\end{equation}
\end{method}

\begin{example}[Repeated eigenvalue, $2\times 2$]\label{ex:repeated-2x2}
Solve $X' = \begin{pmatrix}3 & 1 \\ 0 & 3\end{pmatrix}X$.

The eigenvalue $\lambda=3$ has algebraic multiplicity $2$ but geometric
multiplicity $1$.  Eigenvector: $v_1=\begin{pmatrix}1\\0\end{pmatrix}$.
Generalized eigenvector from $(A-3I)v_2 = v_1$:
$v_2 = \begin{pmatrix}0\\1\end{pmatrix}$.
General solution:
\[
  X(t) = c_1 e^{3t}\begin{pmatrix}1\\0\end{pmatrix}
        + c_2 e^{3t}\left[\begin{pmatrix}0\\1\end{pmatrix}
        + t\begin{pmatrix}1\\0\end{pmatrix}\right].
\]
\end{example}

\begin{remark}\label{rem:exp-jordan-block}
For a Jordan block $J = \lambda I + N$ where $N$ is the nilpotent part,
\[
  e^{tJ} = e^{\lambda t}\,e^{tN}
  = e^{\lambda t}\begin{pmatrix}
    1 & t & t^2/2 & \cdots & t^{m-1}/(m-1)!\\
    0 & 1 & t & \cdots & t^{m-2}/(m-2)!\\
    \vdots & & \ddots & & \vdots \\
    0 & 0 & \cdots & 1 & t \\
    0 & 0 & \cdots & 0 & 1
  \end{pmatrix}.
\]
\end{remark}

%--------------------------------------------------------------------
\subsection{Complete Classification of $2\times 2$ Linear Systems}
\label{subsec:2x2-classification}
\index{phase portrait!classification}\index{classification of 2D systems}
%--------------------------------------------------------------------

Let $A\in\mathcal{M}_2(\R)$ with trace $\tau = \tr(A)$ and determinant
$\Delta = \det(A)$.  The eigenvalues satisfy
$\lambda^2 - \tau\lambda + \Delta = 0$
with discriminant $D = \tau^2 - 4\Delta$.
\index{trace-determinant plane}

\begin{theorem}[Classification of $2\times 2$ linear systems]\label{thm:2x2-class}
Assuming $\Delta\neq 0$ (so the origin is the only equilibrium):
\begin{enumerate}[label=(\roman*)]
  \item \textbf{$D>0$, $\Delta>0$, $\tau<0$:} stable node
    \index{stable node} (two negative real eigenvalues).
  \item \textbf{$D>0$, $\Delta>0$, $\tau>0$:} unstable node
    \index{unstable node} (two positive real eigenvalues).
  \item \textbf{$\Delta<0$:} saddle point
    \index{saddle point} (one positive, one negative eigenvalue).
  \item \textbf{$D<0$, $\tau<0$:} stable spiral (focus)
    \index{stable spiral} (complex eigenvalues with negative real part).
  \item \textbf{$D<0$, $\tau>0$:} unstable spiral (focus)
    \index{unstable spiral} (complex eigenvalues with positive real part).
  \item \textbf{$D<0$, $\tau=0$:} center
    \index{center} (purely imaginary eigenvalues).
  \item \textbf{$D=0$, $\tau\neq 0$:} degenerate (star) node
    \index{star node}\index{degenerate node} or defective node depending on
    geometric multiplicity.
\end{enumerate}
\end{theorem}

%--------------------------------------------------------------------
\section{Phase Portraits for 2D Linear Systems}\label{sec:phase-portraits}
\index{phase portrait}
%--------------------------------------------------------------------

We illustrate the main types with phase portraits.

\begin{figure}[htbp]
\centering
\begin{tikzpicture}[>=Stealth, scale=0.85]
  % Stable node: eigenvalues -1, -3
  \draw[->] (-2.5,0)--(2.5,0) node[right]{$x_1$};
  \draw[->] (0,-2.5)--(0,2.5) node[above]{$x_2$};
  \node at (0,-3){Stable node: $\lambda_1=-3,\;\lambda_2=-1$};
  \foreach \a in {0,30,...,330}{
    \draw[->,blue!70!black,thick]
      ({2*cos(\a)},{2*sin(\a)}) --
      ({2*cos(\a)*0.4},{2*sin(\a)*0.65});
  }
  \foreach \a in {0,45,...,315}{
    \draw[->,blue!70!black]
      ({1.2*cos(\a)},{1.2*sin(\a)}) --
      ({1.2*cos(\a)*0.3},{1.2*sin(\a)*0.55});
  }
  \filldraw (0,0) circle(1.5pt);
\end{tikzpicture}
\hfill
\begin{tikzpicture}[>=Stealth, scale=0.85]
  % Unstable node: eigenvalues 1, 2
  \draw[->] (-2.5,0)--(2.5,0) node[right]{$x_1$};
  \draw[->] (0,-2.5)--(0,2.5) node[above]{$x_2$};
  \node at (0,-3){Unstable node: $\lambda_1=1,\;\lambda_2=2$};
  \foreach \a in {0,30,...,330}{
    \draw[->,red!70!black,thick]
      ({0.5*cos(\a)},{0.5*sin(\a)}) --
      ({0.5*cos(\a)*2.8},{0.5*sin(\a)*2.1});
  }
  \filldraw (0,0) circle(1.5pt);
\end{tikzpicture}
\caption{Stable node (left) and unstable node (right).}
\label{fig:nodes}
\end{figure}

\begin{figure}[htbp]
\centering
\begin{tikzpicture}[>=Stealth, scale=0.85]
  % Saddle: eigenvalues -2, 1
  \draw[->] (-2.5,0)--(2.5,0) node[right]{$x_1$};
  \draw[->] (0,-2.5)--(0,2.5) node[above]{$x_2$};
  \node at (0,-3){Saddle: $\lambda_1=-2,\;\lambda_2=1$};
  % Stable manifold along x-axis
  \draw[->,green!50!black,very thick] (2.2,0)--(0.4,0);
  \draw[->,green!50!black,very thick] (-2.2,0)--(-0.4,0);
  % Unstable manifold along y-axis
  \draw[->,red!60!black,very thick] (0,0.4)--(0,2.2);
  \draw[->,red!60!black,very thick] (0,-0.4)--(0,-2.2);
  % Hyperbolic trajectories
  \foreach \s in {1,-1}{
    \foreach \c in {0.4,0.8,1.3}{
      \draw[blue!60!black,thick,->] plot[domain=0.25:1.5,samples=20]
        ({\s*\c*exp(-2*\x)},{\c*0.5*exp(\x)});
      \draw[blue!60!black,thick,->] plot[domain=0.25:1.5,samples=20]
        ({\s*\c*exp(-2*\x)},{-\c*0.5*exp(\x)});
    }
  }
  \filldraw (0,0) circle(1.5pt);
\end{tikzpicture}
\hfill
\begin{tikzpicture}[>=Stealth, scale=0.85]
  % Center: eigenvalues ±2i
  \draw[->] (-2.5,0)--(2.5,0) node[right]{$x_1$};
  \draw[->] (0,-2.5)--(0,2.5) node[above]{$x_2$};
  \node at (0,-3){Center: $\lambda=\pm 2i$};
  \foreach \r in {0.5,1.0,1.5,2.0}{
    \draw[blue!60!black,thick,
      decoration={markings,mark=at position 0.15 with {\arrow{>}}},
      postaction={decorate}]
      (0,0) circle(\r);
  }
  \filldraw (0,0) circle(1.5pt);
\end{tikzpicture}
\caption{Saddle point (left) and center (right).}
\label{fig:saddle-center}
\end{figure}

\begin{figure}[htbp]
\centering
\begin{tikzpicture}[>=Stealth, scale=0.85]
  % Stable spiral: eigenvalues -0.3 ± 2i
  \draw[->] (-2.5,0)--(2.5,0) node[right]{$x_1$};
  \draw[->] (0,-2.5)--(0,2.5) node[above]{$x_2$};
  \node at (0,-3){Stable spiral: $\lambda=-0.3\pm 2i$};
  \draw[blue!70!black,thick,
    decoration={markings,mark=at position 0.85 with {\arrow{>}}},
    postaction={decorate}]
    plot[domain=0:10,samples=50]
      ({2*exp(-0.3*\x)*cos(2*\x r)},{2*exp(-0.3*\x)*sin(2*\x r)});
  \draw[blue!70!black,thick,
    decoration={markings,mark=at position 0.85 with {\arrow{>}}},
    postaction={decorate}]
    plot[domain=0:10,samples=50]
      ({-2*exp(-0.3*\x)*cos(2*\x r)},{-2*exp(-0.3*\x)*sin(2*\x r)});
  \filldraw (0,0) circle(1.5pt);
\end{tikzpicture}
\hfill
\begin{tikzpicture}[>=Stealth, scale=0.85]
  % Unstable spiral: eigenvalues 0.3 ± 2i
  \draw[->] (-2.5,0)--(2.5,0) node[right]{$x_1$};
  \draw[->] (0,-2.5)--(0,2.5) node[above]{$x_2$};
  \node at (0,-3){Unstable spiral: $\lambda=0.3\pm 2i$};
  \draw[red!70!black,thick,
    decoration={markings,mark=at position 0.15 with {\arrow{>}}},
    postaction={decorate}]
    plot[domain=0:10,samples=50]
      ({2*exp(-0.3*\x)*cos(2*\x r)},{2*exp(-0.3*\x)*sin(2*\x r)});
  \draw[red!70!black,thick,
    decoration={markings,mark=at position 0.15 with {\arrow{>}}},
    postaction={decorate}]
    plot[domain=0:10,samples=50]
      ({-2*exp(-0.3*\x)*cos(2*\x r)},{-2*exp(-0.3*\x)*sin(2*\x r)});
  \filldraw (0,0) circle(1.5pt);
\end{tikzpicture}
\caption{Stable spiral (left) and unstable spiral (right).}
\label{fig:spirals}
\end{figure}

\begin{figure}[htbp]
\centering
\begin{tikzpicture}[>=Stealth, scale=0.85]
  % Trace-determinant plane
  \draw[->] (-3.2,0)--(3.2,0) node[right]{$\tau$};
  \draw[->] (0,-1.5)--(0,3.5) node[above]{$\Delta$};
  % Parabola D=0: Delta = tau^2/4
  \draw[thick,orange,domain=-3:3,samples=40]
    plot(\x,{\x*\x/4}) node[right]{$D=0$};
  % Regions
  \node at (-1.8,2.8) {\footnotesize stable spiral};
  \node at (1.8,2.8) {\footnotesize unstable spiral};
  \node at (-2.2,0.8) {\footnotesize stable node};
  \node at (2.2,0.8) {\footnotesize unstable node};
  \node at (0,-0.8) {\footnotesize saddle};
  \node[blue] at (0,1.8) {\footnotesize centers};
  % Delta axis marking
  \draw[blue,very thick] (0,0.05)--(0,3.3);
\end{tikzpicture}
\caption{Classification in the trace-determinant plane.}
\label{fig:trace-det}
\end{figure}

%--------------------------------------------------------------------
\section{Fundamental Matrix and Wronskian}\label{sec:fundamental-matrix}
%--------------------------------------------------------------------

\begin{definition}[Fundamental matrix]\label{def:fundamental-matrix}
\index{fundamental matrix}
A \textbf{fundamental matrix} for $X'=A(t)X$ is an $n\times n$ matrix
$\Phi(t)$ whose columns form a basis of the solution space.  Equivalently,
$\Phi'(t)=A(t)\Phi(t)$ and $\det\Phi(t)\neq 0$ for all $t\in I$.
\end{definition}

\begin{definition}[Wronskian of a system]\label{def:wronskian-sys}
\index{Wronskian!of a system}
If $X_1,\dots,X_n$ are $n$ solutions of $X'=A(t)X$, their
\textbf{Wronskian} is $W(t) = \det\bigl(X_1(t)\mid\cdots\mid X_n(t)\bigr)$.
\end{definition}

\begin{theorem}[Abel--Liouville formula for systems]\label{thm:abel-liouville-sys}
\index{Abel--Liouville formula}\index{Liouville formula}
If $\Phi(t)$ is a fundamental matrix for $X'=A(t)X$, then
\begin{equation}\label{eq:abel-liouville-sys}
  W(t) = \det\Phi(t) = W(t_0)\,\exp\!\left(\int_{t_0}^{t}\tr A(s)\,\dd s\right).
\end{equation}
\end{theorem}

\begin{proof}
Write $W(t)=\det\Phi(t)$.  Using the formula
$\frac{\dd}{\dd t}\det\Phi = \det\Phi\cdot\tr(\Phi^{-1}\Phi')$
and noting $\Phi^{-1}\Phi' = \Phi^{-1}A\Phi$, so
$\tr(\Phi^{-1}\Phi') = \tr(\Phi^{-1}A\Phi)=\tr(A)$, we get
\[
  W'(t) = \tr A(t)\cdot W(t).
\]
This is a scalar linear ODE whose solution is \eqref{eq:abel-liouville-sys}.
\end{proof}

\begin{corollary}\label{cor:wronskian-zero}
\index{Wronskian!zero or nonzero}
The Wronskian $W(t)$ is either identically zero or never zero on $I$.
\end{corollary}

\begin{remark}\label{rem:principal-matrix}
The \textbf{principal fundamental matrix}\index{principal fundamental matrix}
at $t_0$ is the unique fundamental matrix $\Phi$ satisfying
$\Phi(t_0)=I_n$.  For constant coefficients, $\Phi(t)=e^{(t-t_0)A}$.
\end{remark}

%--------------------------------------------------------------------
\section{Non-Homogeneous Systems: Variation of Parameters}
\label{sec:sys-var-params}
\index{variation of parameters!systems}
%--------------------------------------------------------------------

\begin{theorem}[Variation of parameters for systems]\label{thm:var-params-sys}
Let $\Phi(t)$ be a fundamental matrix of $X'=A(t)X$.  The general
solution of $X'=A(t)X+B(t)$ is
\begin{equation}\label{eq:var-params-sys}
  X(t) = \Phi(t)\,C + \Phi(t)\int_{t_0}^{t}\Phi^{-1}(s)\,B(s)\,\dd s,
\end{equation}
where $C\in\R^n$ is an arbitrary constant vector.
\end{theorem}

\begin{proof}
We seek a particular solution of the form $X_p(t) = \Phi(t)\,U(t)$
with $U:I\to\R^n$ to be determined.  Differentiating:
\[
  X_p' = \Phi' U + \Phi U' = A\Phi U + \Phi U'.
\]
We need $X_p' = AX_p + B = A\Phi U + B$, hence $\Phi U' = B$,
i.e., $U'(t) = \Phi^{-1}(t)B(t)$.  Integrating gives the result.
\end{proof}

\begin{example}[Non-homogeneous $2\times 2$ system]\label{ex:nonhom-sys}
Solve $X' = \begin{pmatrix}2&1\\0&2\end{pmatrix}X
+ \begin{pmatrix}0\\e^{2t}\end{pmatrix}$,\quad $X(0)=\begin{pmatrix}0\\0\end{pmatrix}$.

The homogeneous system has a repeated eigenvalue $\lambda=2$.
A fundamental matrix is
\[
  \Phi(t) = e^{2t}\begin{pmatrix}1&t\\0&1\end{pmatrix},
  \qquad
  \Phi^{-1}(t) = e^{-2t}\begin{pmatrix}1&-t\\0&1\end{pmatrix}.
\]
Then $\Phi^{-1}(t)B(t) = e^{-2t}\begin{pmatrix}1&-t\\0&1\end{pmatrix}
\begin{pmatrix}0\\e^{2t}\end{pmatrix}
= \begin{pmatrix}-t\\1\end{pmatrix}$.

Integrating from $0$ to $t$:
$\int_0^t \begin{pmatrix}-s\\1\end{pmatrix}\dd s
= \begin{pmatrix}-t^2/2\\t\end{pmatrix}$.

So $X_p(t) = e^{2t}\begin{pmatrix}1&t\\0&1\end{pmatrix}
\begin{pmatrix}-t^2/2\\t\end{pmatrix}
= e^{2t}\begin{pmatrix}-t^2/2+t^2\\t\end{pmatrix}
= e^{2t}\begin{pmatrix}t^2/2\\t\end{pmatrix}$.

With $X(0)=0$, the solution is
$X(t) = e^{2t}\begin{pmatrix}t^2/2\\t\end{pmatrix}$.
\end{example}

%--------------------------------------------------------------------
\section{Exercises}\label{sec:sys-exercises}
%--------------------------------------------------------------------

\begin{exercise}\label{exo:sys-1}
Solve the system $X'=\begin{pmatrix}-1&1\\0&-2\end{pmatrix}X$
with $X(0)=\begin{pmatrix}3\\1\end{pmatrix}$.
\end{exercise}

\begin{exercise}\label{exo:sys-2}
Compute $e^{tA}$ for $A=\begin{pmatrix}0&1\\-1&0\end{pmatrix}$.
\end{exercise}

\begin{exercise}\label{exo:sys-3}
Show that if $A$ is antisymmetric ($A^T=-A$) then $e^{tA}$ is orthogonal
for every $t\in\R$.
\end{exercise}

\begin{exercise}\label{exo:sys-4}
Classify the equilibrium at the origin and sketch the phase portrait for:
\begin{enumerate}[label=(\alph*)]
  \item $X'=\begin{pmatrix}1&-5\\1&-3\end{pmatrix}X$.
  \item $X'=\begin{pmatrix}3&-2\\4&-1\end{pmatrix}X$.
  \item $X'=\begin{pmatrix}1&1\\4&1\end{pmatrix}X$.
\end{enumerate}
\end{exercise}

\begin{exercise}\label{exo:sys-5}
Solve $X'=\begin{pmatrix}1&0\\0&2\end{pmatrix}X
+\begin{pmatrix}e^t\\0\end{pmatrix}$ with $X(0)=\begin{pmatrix}1\\1\end{pmatrix}$.
\end{exercise}

\begin{exercise}\label{exo:sys-6}
Let $A=\begin{pmatrix}2&1&0\\0&2&1\\0&0&2\end{pmatrix}$.
Compute $e^{tA}$ and solve $X'=AX$ with $X(0)=(1,0,0)^T$.
\end{exercise}

\begin{exercise}\label{exo:sys-7}
Verify the Abel--Liouville formula for the system
$X'=\begin{pmatrix}0&1\\-1&0\end{pmatrix}X$
by computing $W(t)$ directly and via \eqref{eq:abel-liouville-sys}.
\end{exercise}

\begin{exercise}\label{exo:sys-8}
Consider the $3\times 3$ system $X'=AX$ with
$A=\begin{pmatrix}-1&0&0\\0&0&-1\\0&1&0\end{pmatrix}$.
Find the general solution and describe the long-term behavior.
\end{exercise}

%--------------------------------------------------------------------
\section{Chapter Summary}\label{sec:sys-summary}
%--------------------------------------------------------------------

\begin{itemize}
  \item A linear system $X'=A(t)X+B(t)$ has a unique solution on the
    entire interval where $A,B$ are continuous.
  \item The solution space of the homogeneous system is an $n$-dimensional
    vector space.
  \item For constant coefficients, the solution is $X(t)=e^{tA}X_0$, and
    $e^{tA}$ can be computed via diagonalization, Jordan form, or
    Cayley--Hamilton.
  \item For $2\times 2$ systems, the type of equilibrium (node, saddle,
    spiral, center) is determined by the trace and determinant of $A$.
  \item The Wronskian satisfies $W'=\tr(A)\,W$ (Abel--Liouville), so it
    is either always zero or never zero.
  \item Variation of parameters:
    $X_p(t)=\Phi(t)\int_{t_0}^t\Phi^{-1}(s)B(s)\,\dd s$.
\end{itemize}


%%%%%%%%%%%%%%%%%%%%%%%%%%%%%%%%%%%%%%%%%%%%%%%%%%%%%%%%%%%%%
%  CHAPTER 6 — EXPLICIT SOLUTION METHODS
%%%%%%%%%%%%%%%%%%%%%%%%%%%%%%%%%%%%%%%%%%%%%%%%%%%%%%%%%%%%%
\chapter{Explicit Solution Methods}\label{ch:explicit-methods}
\index{explicit solution methods}

This chapter develops powerful analytic techniques for solving linear
ODEs and systems: the Wronskian, variation of parameters for higher-order
equations, Green's functions, power series methods, and reduction of order.

%--------------------------------------------------------------------
\section{The Wronskian: Definition, Properties, and Abel's Theorem}
\label{sec:wronskian}
\index{Wronskian}
%--------------------------------------------------------------------

\begin{definition}[Wronskian of $n$ functions]\label{def:wronskian}
\index{Wronskian!definition}
Given $n$ functions $y_1,\dots,y_n\in\mathcal{C}^{n-1}(I)$, their
\textbf{Wronskian} is
\begin{equation}\label{eq:wronskian-def}
  W(y_1,\dots,y_n)(t) = \det\begin{pmatrix}
    y_1 & y_2 & \cdots & y_n \\
    y_1' & y_2' & \cdots & y_n' \\
    \vdots & \vdots & & \vdots \\
    y_1^{(n-1)} & y_2^{(n-1)} & \cdots & y_n^{(n-1)}
  \end{pmatrix}(t).
\end{equation}
\end{definition}

\begin{proposition}[Wronskian and linear independence]\label{prop:wronskian-indep}
\index{Wronskian!linear independence}
Let $y_1,\dots,y_n$ be solutions of the $n$th-order linear ODE
\begin{equation}\label{eq:nth-order-hom}
  y^{(n)} + p_{n-1}(t)y^{(n-1)} + \cdots + p_1(t)y' + p_0(t)y = 0
\end{equation}
with continuous coefficients on $I$.  Then $y_1,\dots,y_n$ are linearly
independent if and only if $W(y_1,\dots,y_n)(t)\neq 0$ for some
(equivalently, all) $t\in I$.
\end{proposition}

\begin{theorem}[Abel's theorem]\label{thm:abel}
\index{Abel's theorem}
Let $y_1,\dots,y_n$ be solutions of \eqref{eq:nth-order-hom}.
Then
\begin{equation}\label{eq:abel}
  W(t) = W(t_0)\,\exp\!\left(-\int_{t_0}^{t}p_{n-1}(s)\,\dd s\right).
\end{equation}
\end{theorem}

\begin{proof}
We give the full proof for general $n$.  Write the equation in system
form by setting $X=(y,y',\dots,y^{(n-1)})^T$.  Then $X'=A(t)X$ where
\[
  A(t) = \begin{pmatrix}
    0 & 1 & 0 & \cdots & 0 \\
    0 & 0 & 1 & \cdots & 0 \\
    \vdots & & & \ddots & \vdots \\
    0 & 0 & 0 & \cdots & 1 \\
    -p_0 & -p_1 & -p_2 & \cdots & -p_{n-1}
  \end{pmatrix}.
\]
The trace of this companion matrix\index{companion matrix} is
$\tr A(t) = -p_{n-1}(t)$.  By the Abel--Liouville formula
(Theorem~\ref{thm:abel-liouville-sys}),
\[
  W(t) = W(t_0)\exp\!\left(\int_{t_0}^t \tr A(s)\,\dd s\right)
  = W(t_0)\exp\!\left(-\int_{t_0}^t p_{n-1}(s)\,\dd s\right).
  \qedhere
\]
\end{proof}

\begin{example}[Abel's theorem for second-order equations]\label{ex:abel-2nd}
For $y''+p(t)y'+q(t)y=0$ with solutions $y_1,y_2$:
\[
  W(y_1,y_2)(t) = W(t_0)\,e^{-\int_{t_0}^t p(s)\,\dd s}.
\]
In particular, for $y''-y=0$ ($p\equiv 0$), the solutions
$\cosh t, \sinh t$ have constant Wronskian $W=1$.
\end{example}

\begin{proposition}[Properties of the Wronskian]\label{prop:wronskian-props}
\index{Wronskian!properties}
Let $y_1,\dots,y_n$ be solutions of \eqref{eq:nth-order-hom}.
\begin{enumerate}[label=(\roman*)]
  \item $W$ is either identically zero or never zero on $I$.
  \item $W\neq 0$ if and only if $\{y_1,\dots,y_n\}$ is a fundamental
    system of solutions.
  \item $W$ satisfies the first-order ODE $W' + p_{n-1}(t)W = 0$.
  \item If the equation has constant coefficients ($p_i$ constant),
    then $W(t)=W(0)\,e^{-p_{n-1}t}$.
\end{enumerate}
\end{proposition}

%--------------------------------------------------------------------
\section{Variation of Parameters}\label{sec:var-params}
\index{variation of parameters}
%--------------------------------------------------------------------

\subsection{For $n$th-Order Linear ODEs}\label{subsec:var-params-nth}
\index{variation of parameters!$n$th-order}

Consider the non-homogeneous equation
\begin{equation}\label{eq:nth-order-nonhom}
  y^{(n)} + p_{n-1}(t)y^{(n-1)} + \cdots + p_0(t)y = g(t).
\end{equation}

\begin{theorem}[Variation of parameters -- $n$th-order]\label{thm:var-params-nth}
Let $y_1,\dots,y_n$ be a fundamental system of solutions of the
homogeneous equation \eqref{eq:nth-order-hom} with Wronskian $W$.
A particular solution of \eqref{eq:nth-order-nonhom} is
\begin{equation}\label{eq:var-params-nth}
  y_p(t) = \sum_{k=1}^{n} y_k(t)\int_{t_0}^{t}
    \frac{W_k(s)}{W(s)}\,g(s)\,\dd s,
\end{equation}
where $W_k$ is obtained from $W$ by replacing the $k$th column with
$(0,\dots,0,1)^T$.
\end{theorem}

\begin{proof}
We seek $y_p = \sum_{k=1}^n u_k(t)\,y_k(t)$ with the $n-1$ constraints
\begin{equation}\label{eq:vp-constraints}
  \sum_{k=1}^n u_k'(t)\,y_k^{(j)}(t) = 0, \qquad j=0,1,\dots,n-2.
\end{equation}
Under these constraints, $y_p^{(j)} = \sum_k u_k y_k^{(j)}$ for
$j=0,\dots,n-1$, and
\[
  y_p^{(n)} = \sum_{k=1}^n u_k y_k^{(n)} + \sum_{k=1}^n u_k' y_k^{(n-1)}.
\]
Substituting into \eqref{eq:nth-order-nonhom} and using the fact that each
$y_k$ solves the homogeneous equation:
\[
  \sum_{k=1}^n u_k' y_k^{(n-1)} = g(t).
\]
Together with \eqref{eq:vp-constraints}, we have the system
\begin{equation}\label{eq:vp-system}
  \begin{pmatrix}
    y_1 & \cdots & y_n \\
    y_1' & \cdots & y_n' \\
    \vdots & & \vdots \\
    y_1^{(n-1)} & \cdots & y_n^{(n-1)}
  \end{pmatrix}
  \begin{pmatrix} u_1' \\ u_2' \\ \vdots \\ u_n' \end{pmatrix}
  = \begin{pmatrix} 0 \\ 0 \\ \vdots \\ g(t) \end{pmatrix}.
\end{equation}
By Cramer's rule, $u_k'(t) = W_k(t)\,g(t)/W(t)$.
Integrating gives \eqref{eq:var-params-nth}.
\end{proof}

\begin{example}[Second-order variation of parameters]\label{ex:var-params-2}
\index{variation of parameters!second order}
Solve $y'' - y = \dfrac{1}{\cosh t}$.

Fundamental system: $y_1=e^t$, $y_2=e^{-t}$, with $W=-2$.
\begin{align*}
  W_1 &= \det\begin{pmatrix}0 & e^{-t}\\ 1/\cosh t & -e^{-t}\end{pmatrix}
       = -\frac{e^{-t}}{\cosh t}, \\
  W_2 &= \det\begin{pmatrix}e^t & 0\\ e^t & 1/\cosh t\end{pmatrix}
       = \frac{e^t}{\cosh t}.
\end{align*}
Thus
\begin{align*}
  u_1'(t) &= \frac{W_1}{W}\cdot\frac{1}{1} = \frac{e^{-t}}{2\cosh t}
           = \frac{1}{1+e^{2t}}, \\
  u_2'(t) &= \frac{W_2}{W}\cdot\frac{1}{1} = -\frac{e^t}{2\cosh t}
           = -\frac{1}{1+e^{-2t}}.
\end{align*}
Integrating: $u_1(t) = \arctan(e^t)$ and
$u_2(t) = \arctan(e^{-t})$.
A particular solution is
\[
  y_p(t) = e^t\arctan(e^t) + e^{-t}\arctan(e^{-t}).
\]
\end{example}

\begin{example}[Third-order variation of parameters]\label{ex:var-params-3}
Solve $y'''+y'=\sec t$ on $(-\pi/2,\pi/2)$.

The characteristic equation $r^3+r=r(r^2+1)=0$ gives $r=0,\pm i$.
Fundamental system: $y_1=1$, $y_2=\cos t$, $y_3=\sin t$.
The Wronskian is
\[
  W = \det\begin{pmatrix}
    1 & \cos t & \sin t \\
    0 & -\sin t & \cos t \\
    0 & -\cos t & -\sin t
  \end{pmatrix} = 1.
\]
By Cramer's rule:
\begin{align*}
  u_1' &= \sec t\cdot\det\begin{pmatrix}
    0 & \cos t \\ 0 & -\sin t
  \end{pmatrix}\Big/1 \;(\text{with appropriate expansion})\\
  &= \sec t.
\end{align*}
Working through all the cofactors yields
$u_1'=\sec t$, $u_2'=0$, $u_3'=-1$, giving
$u_1=\ln\abs{\sec t+\tan t}$, $u_2=0$, $u_3=-t$.

Particular solution: $y_p = \ln\abs{\sec t+\tan t} - t\sin t$.
\end{example}

\subsection{For Systems}\label{subsec:var-params-systems}
\index{variation of parameters!systems}

The variation of parameters formula for systems was already established in
Theorem~\ref{thm:var-params-sys}.  We restate it here with the initial
condition formulation.

\begin{proposition}[Duhamel formula]\label{prop:duhamel}
\index{Duhamel formula}
The solution of $X'=A(t)X+B(t)$, $X(t_0)=X_0$ is
\begin{equation}\label{eq:duhamel}
  X(t) = \Phi(t,t_0)\,X_0 + \int_{t_0}^{t}\Phi(t,s)\,B(s)\,\dd s,
\end{equation}
where $\Phi(t,s)$ is the state transition matrix\index{state transition matrix},
i.e., $\Phi(t,s) = \Phi(t)\Phi^{-1}(s)$ with $\Phi$ any fundamental matrix.
For constant $A$, $\Phi(t,s)=e^{(t-s)A}$.
\end{proposition}

%--------------------------------------------------------------------
\section{Green's Function}\label{sec:green}
\index{Green's function}
%--------------------------------------------------------------------

The variation of parameters formula can be recast using an integral kernel
known as the \emph{Green's function}.

\begin{definition}[Green's function for an IVP]\label{def:green-ivp}
\index{Green's function!initial value problem}
For the $n$th-order equation \eqref{eq:nth-order-nonhom} with initial
conditions $y(t_0)=y'(t_0)=\cdots=y^{(n-1)}(t_0)=0$, the
\textbf{Green's function} $G(t,s)$ is defined for $t_0\le s\le t$ by
\begin{equation}\label{eq:green-ivp}
  G(t,s) = \sum_{k=1}^{n}\frac{W_k(s)}{W(s)}\,y_k(t),
\end{equation}
where $y_1,\dots,y_n$ and $W_k$ are as in Theorem~\ref{thm:var-params-nth}.
\end{definition}

\begin{theorem}[Solution via Green's function]\label{thm:green-solution}
The solution of \eqref{eq:nth-order-nonhom} with zero initial conditions
at $t_0$ is
\begin{equation}\label{eq:green-sol}
  y(t) = \int_{t_0}^{t}G(t,s)\,g(s)\,\dd s.
\end{equation}
\end{theorem}

\begin{proposition}[Properties of the Green's function]\label{prop:green-props}
\index{Green's function!properties}
The Green's function $G(t,s)$ satisfies:
\begin{enumerate}[label=(\roman*)]
  \item As a function of $t$ (with $s$ fixed), $G(\cdot,s)$ solves the
    homogeneous equation for $t>s$.
  \item $G(s,s)=G'(s,s)=\cdots=G^{(n-2)}(s,s)=0$ and $G^{(n-1)}(s,s)=1$,
    where derivatives are with respect to $t$.
  \item $G(t,s)=0$ for $t<s$ (causality).
\end{enumerate}
\end{proposition}

\begin{example}[Green's function for $y''+y=g(t)$]\label{ex:green-2nd}
\index{Green's function!example}
Fundamental system: $y_1=\cos t$, $y_2=\sin t$, $W=1$.
\[
  G(t,s) = \frac{-\sin s\cdot\cos t + \cos s\cdot\sin t}{1}
         = \sin(t-s).
\]
The solution with $y(0)=y'(0)=0$ is
$y(t)=\int_0^t\sin(t-s)\,g(s)\,\dd s$.
This is a \emph{convolution}\index{convolution} when $g$ is defined on $[0,\infty)$.
\end{example}

\begin{example}[Green's function for a boundary value problem]\label{ex:green-bvp}
\index{Green's function!boundary value problem}
Consider $y''+y=g(t)$ on $[0,\pi/2]$ with $y(0)=0$, $y(\pi/2)=0$.

The homogeneous solutions satisfying one boundary condition each are:
$\phi_1(t)=\sin t$ (satisfies $\phi_1(0)=0$) and
$\phi_2(t)=\cos t$ (satisfies $\phi_2(\pi/2)=0$).
The Wronskian of $\phi_1,\phi_2$ is $W=-1$.

The Green's function is
\[
  G(t,s) = \begin{cases}
    -\sin(s)\cos(t) & \text{if } t\ge s, \\
    -\sin(t)\cos(s) & \text{if } t\le s.
  \end{cases}
\]
The solution is $y(t)=\int_0^{\pi/2}G(t,s)\,g(s)\,\dd s$.
\end{example}

%--------------------------------------------------------------------
\section{Power Series Solutions}\label{sec:power-series}
\index{power series solutions}
%--------------------------------------------------------------------

\subsection{Ordinary Points}\label{subsec:ordinary-points}
\index{ordinary point}

Consider the second-order equation
\begin{equation}\label{eq:2nd-order-standard}
  y'' + P(t)y' + Q(t)y = 0.
\end{equation}

\begin{definition}[Ordinary and singular points]\label{def:ordinary-singular}
\index{singular point!ODE}
A point $t_0$ is an \textbf{ordinary point} of \eqref{eq:2nd-order-standard}
if $P$ and $Q$ are analytic at $t_0$.  Otherwise $t_0$ is a
\textbf{singular point}.
\end{definition}

\begin{theorem}[Power series at an ordinary point]\label{thm:power-series-ordinary}
\index{power series solutions!ordinary point}
If $t_0$ is an ordinary point of \eqref{eq:2nd-order-standard}, then every
solution can be represented as a power series
\begin{equation}\label{eq:ps-solution}
  y(t) = \sum_{n=0}^{\infty}a_n(t-t_0)^n
\end{equation}
convergent at least on $\abs{t-t_0}<R$, where $R$ is the distance from
$t_0$ to the nearest singular point (in $\C$).
\end{theorem}

\begin{method}[Finding power series solutions]\label{meth:power-series}
\index{power series solutions!method}
\begin{enumerate}[label=\textbf{Step \arabic*:}]
  \item Substitute $y=\sum_{n=0}^{\infty}a_n(t-t_0)^n$ into the ODE.
  \item Compute $y'=\sum_{n=1}^{\infty}na_n(t-t_0)^{n-1}$ and
    $y''=\sum_{n=2}^{\infty}n(n-1)a_n(t-t_0)^{n-2}$.
  \item Multiply through by the coefficient functions and re-index so all
    series have the same power of $(t-t_0)$.
  \item Set all coefficients to zero to obtain a recurrence relation for
    the $a_n$.
  \item Solve the recurrence; identify the two linearly independent
    solutions from free parameters $a_0$ and $a_1$.
\end{enumerate}
\end{method}

\begin{example}[Airy equation]\label{ex:airy}
\index{Airy equation}
Solve $y'' - ty = 0$ near $t_0=0$.

Substituting $y=\sum a_n t^n$:
\[
  \sum_{n=2}^{\infty}n(n-1)a_n t^{n-2} - \sum_{n=0}^{\infty}a_n t^{n+1}=0.
\]
Shifting indices: let $m=n-2$ in the first sum and $m=n+1$ in the second:
\[
  \sum_{m=0}^{\infty}(m+2)(m+1)a_{m+2}\,t^m
  - \sum_{m=1}^{\infty}a_{m-1}\,t^m = 0.
\]
For $m=0$: $2a_2=0$, so $a_2=0$.  For $m\ge 1$:
\begin{equation}\label{eq:airy-recurrence}
  a_{m+2} = \frac{a_{m-1}}{(m+2)(m+1)}.
\end{equation}
The recurrence produces two independent solutions:
\begin{align*}
  y_1(t) &= a_0\!\left(1+\frac{t^3}{2\cdot 3}
    +\frac{t^6}{2\cdot 3\cdot 5\cdot 6}+\cdots\right), \\
  y_2(t) &= a_1\!\left(t+\frac{t^4}{3\cdot 4}
    +\frac{t^7}{3\cdot 4\cdot 6\cdot 7}+\cdots\right).
\end{align*}
Both series converge for all $t\in\R$ (no finite singular points).
The functions $\operatorname{Ai}(t)$ and $\operatorname{Bi}(t)$ are
specific linear combinations of $y_1$ and $y_2$.
\end{example}

\begin{example}[Hermite equation]\label{ex:hermite}
\index{Hermite equation}
Solve $y''-2ty'+2\alpha y=0$ ($\alpha\in\R$) near $t_0=0$.

With $y=\sum a_n t^n$, we obtain the recurrence
\[
  a_{n+2} = \frac{2(n-\alpha)}{(n+1)(n+2)}\,a_n, \qquad n\ge 0.
\]
The even coefficients depend on $a_0$, the odd on $a_1$.
When $\alpha=m\in\N$, one of the two series terminates, producing
the \emph{Hermite polynomials}\index{Hermite polynomial} $H_m(t)$.
\end{example}

\subsection{Regular Singular Points and the Frobenius Method}
\label{subsec:frobenius}
\index{Frobenius method}\index{regular singular point}

\begin{definition}[Regular singular point]\label{def:regular-singular}
A singular point $t_0$ of $y''+P(t)y'+Q(t)y=0$ is a
\textbf{regular singular point} if $(t-t_0)P(t)$ and
$(t-t_0)^2Q(t)$ are analytic at $t_0$.  Otherwise it is an
\textbf{irregular singular point}\index{irregular singular point}.
\end{definition}

\begin{theorem}[Frobenius method]\label{thm:frobenius}
\index{Frobenius method!theorem}
If $t_0$ is a regular singular point, there exists at least one solution
of the form
\begin{equation}\label{eq:frobenius-form}
  y(t) = (t-t_0)^r\sum_{n=0}^{\infty}a_n(t-t_0)^n, \qquad a_0\neq 0,
\end{equation}
where $r$ satisfies the \textbf{indicial equation}\index{indicial equation}.
\end{theorem}

\begin{method}[Frobenius procedure]\label{meth:frobenius}
Write the ODE in the form
\[
  t^2y''+t\bigl[tp_0+tp_1t+\cdots\bigr]y'+\bigl[q_0+q_1t+\cdots\bigr]y=0
\]
(assuming $t_0=0$ for simplicity).
\begin{enumerate}[label=\textbf{Step \arabic*:}]
  \item Write the \textbf{indicial equation}:
    $r(r-1)+p_0\,r+q_0=0$.
  \item Find the two roots $r_1\ge r_2$.
  \item \textbf{Case 1} ($r_1-r_2\notin\Z$): two Frobenius series.
  \item \textbf{Case 2} ($r_1=r_2$): second solution involves a
    logarithmic term.
  \item \textbf{Case 3} ($r_1-r_2\in\Z_{>0}$): the larger root always
    gives a Frobenius series; the second solution may involve a logarithm.
\end{enumerate}
\end{method}

\begin{example}[Bessel equation of order zero]\label{ex:bessel-0}
\index{Bessel equation}
Consider $t^2y''+ty'+t^2y=0$ (regular singular point at $t_0=0$).

The indicial equation is $r^2=0$, giving a double root $r=0$.
The Frobenius series is
\[
  y_1(t) = J_0(t) = \sum_{n=0}^{\infty}\frac{(-1)^n}{(n!)^2}
    \left(\frac{t}{2}\right)^{2n}
  = 1 - \frac{t^2}{4} + \frac{t^4}{64} - \cdots
\]
The second solution involves a logarithmic term:
\[
  y_2(t) = J_0(t)\ln t + \sum_{n=1}^{\infty}
    \frac{(-1)^{n+1}H_n}{(n!)^2}\left(\frac{t}{2}\right)^{2n},
\]
where $H_n=1+\frac{1}{2}+\cdots+\frac{1}{n}$ is the $n$th harmonic number.
\end{example}

\begin{example}[Euler--Cauchy equation]\label{ex:euler-cauchy-frobenius}
\index{Euler--Cauchy equation!Frobenius}
The equation $t^2y''+\alpha ty'+\beta y=0$ has $t_0=0$ as a regular
singular point.  The indicial equation $r(r-1)+\alpha r+\beta=0$
recovers the standard Euler--Cauchy approach:
\begin{itemize}
  \item Distinct real roots $r_1\neq r_2$:
    $y=c_1\abs{t}^{r_1}+c_2\abs{t}^{r_2}$.
  \item Double root $r_1=r_2=r$: $y=\abs{t}^r(c_1+c_2\ln\abs{t})$.
  \item Complex roots $r=\alpha'\pm i\beta'$:
    $y=\abs{t}^{\alpha'}[c_1\cos(\beta'\ln\abs{t})+c_2\sin(\beta'\ln\abs{t})]$.
\end{itemize}
\end{example}

%--------------------------------------------------------------------
\section{Reduction of Order}\label{sec:reduction-order}
\index{reduction of order}
%--------------------------------------------------------------------

\begin{theorem}[Reduction of order]\label{thm:reduction-order}
If $y_1$ is a nonzero solution of $y''+P(t)y'+Q(t)y=0$, then a
second linearly independent solution is
\begin{equation}\label{eq:reduction-order}
  y_2(t) = y_1(t)\int\frac{e^{-\int P(t)\,\dd t}}{y_1(t)^2}\,\dd t.
\end{equation}
\end{theorem}

\begin{proof}
Set $y_2=v(t)\,y_1(t)$.  Substituting into the equation:
\[
  v''y_1 + 2v'y_1' + vy_1'' + P(v'y_1+vy_1') + Qvy_1 = 0.
\]
Since $y_1''+Py_1'+Qy_1=0$, this simplifies to
\[
  v''y_1 + v'(2y_1'+Py_1) = 0.
\]
Setting $w=v'$: $w'y_1 + w(2y_1'+Py_1)=0$, i.e.,
$\frac{w'}{w} = -\frac{2y_1'}{y_1}-P$.
Integrating: $\ln\abs{w} = -2\ln\abs{y_1}-\int P\,\dd t$, hence
$w = \frac{e^{-\int P\,\dd t}}{y_1^2}$.
One more integration gives $v$, and thus $y_2=vy_1$.
\end{proof}

\begin{remark}\label{rem:reduction-wronskian}
Note the connection to Abel's theorem: the numerator
$e^{-\int P\,\dd t}$ in \eqref{eq:reduction-order} is proportional to
$W(t)/y_1^2$, consistent with $W=y_1y_2'-y_1'y_2$.
\end{remark}

\begin{example}[Reduction of order]\label{ex:reduction-order-1}
Given that $y_1=t$ solves $t^2y''-2y=0$ for $t>0$ (rewritten as
$y''-\frac{2}{t^2}y=0$, so $P=0$), find a second solution.
\[
  y_2 = t\int\frac{e^{0}}{t^2}\,\dd t = t\int t^{-2}\,\dd t
      = t\cdot(-t^{-1}) = -1.
\]
Hmm, a constant is trivial.  Let us redo more carefully: the equation is
$t^2y''-2y=0$, i.e., $y''-\frac{2}{t^2}y=0$.
We have $P(t)=0$, $Q(t)=-2/t^2$.  Formula:
\[
  y_2 = t\int\frac{1}{t^2}\,\dd t = t\cdot\left(-\frac{1}{t}\right)=-1.
\]
So $y_2=-1$, or equivalently $y_2=1$.  But $y_2=1$ does not solve the
equation: $0-2\neq 0$.  The issue is that we need to verify: the actual
equation $t^2y''-2y=0$ requires $P=0$ and $Q=-2/t^2$ only if we
divide by $t^2$.  Let us try $y_1=t^2$.

Check: $y_1=t^2$, $y_1'=2t$, $y_1''=2$.
Then $t^2\cdot 2-2\cdot t^2=0$.  So $y_1=t^2$ is a solution.
\[
  y_2 = t^2\int\frac{1}{t^4}\,\dd t = t^2\cdot\left(-\frac{1}{3t^3}\right)
      = -\frac{1}{3t}.
\]
So a second independent solution is $y_2=t^{-1}$.  General solution:
$y=c_1t^2+c_2t^{-1}$.
\end{example}

\begin{example}[From a known solution of a variable-coefficient equation]%
\label{ex:reduction-order-2}
Verify that $y_1=e^t$ solves $y''-\frac{2}{t}y'+\left(1+\frac{2}{t^2}\right)y=0$
(for $t>0$ -- note this is not a standard form and requires care) and
find a second solution.

We have $P(t)=-2/t$.  By the reduction of order formula:
\[
  y_2 = e^t\int\frac{\exp\bigl(\int\frac{2}{t}\,\dd t\bigr)}{e^{2t}}\,\dd t
      = e^t\int\frac{t^2}{e^{2t}}\,\dd t.
\]
The integral $\int t^2 e^{-2t}\,\dd t$ can be computed by integration by
parts (twice), yielding
\[
  \int t^2 e^{-2t}\,\dd t
  = -\frac{e^{-2t}}{2}\left(t^2+t+\frac{1}{2}\right).
\]
Hence $y_2 = -\frac{1}{2}\,e^{-t}\bigl(t^2+t+\frac{1}{2}\bigr)$,
or equivalently $y_2 = e^{-t}(2t^2+2t+1)$ (up to a constant multiple).
\end{example}

%--------------------------------------------------------------------
\section{Method Summary}\label{sec:method-summary}
%--------------------------------------------------------------------

\begin{table}[htbp]
\centering
\caption{Summary of explicit solution methods for linear ODEs.}
\label{tab:method-summary}
\index{method summary table}
\renewcommand{\arraystretch}{1.4}
\begin{tabular}{@{}>{\bfseries}p{3.8cm} p{4.5cm} p{4.5cm}@{}}
\toprule
Method & Applicable When & Key Formula/Idea \\
\midrule
Constant coeff.\newline (char.\ equation)
  & $a_ny^{(n)}+\cdots+a_0y=0$,\newline $a_i$ constant
  & Solve $a_nr^n+\cdots+a_0=0$ \\
Variation of\newline parameters
  & Any linear ODE\newline $Ly=g(t)$, given fund.\ system
  & $u_k'=W_k g/W$;\newline integrate \\
Green's function
  & Linear ODE with\newline homogeneous BCs or ICs
  & $y(t)=\int G(t,s)g(s)\,\dd s$ \\
Power series
  & Ordinary point;\newline analytic coefficients
  & Substitute $y=\sum a_n t^n$;\newline recurrence for $a_n$ \\
Frobenius method
  & Regular singular point
  & $y=t^r\sum a_n t^n$;\newline indicial equation \\
Reduction of order
  & One solution $y_1$ known
  & $y_2=y_1\int e^{-\int P}/y_1^2$ \\
Matrix exponential
  & Systems $X'=AX$,\newline $A$ constant
  & $X(t)=e^{tA}X_0$ \\
\bottomrule
\end{tabular}
\end{table}

%--------------------------------------------------------------------
\section{Additional Worked Examples}\label{sec:more-examples}
%--------------------------------------------------------------------

\begin{example}[Power series: Legendre equation]\label{ex:legendre}
\index{Legendre equation}\index{Legendre polynomial}
The Legendre equation $(1-t^2)y''-2ty'+\ell(\ell+1)y=0$ has $t_0=0$
as an ordinary point.  With $y=\sum a_n t^n$, the recurrence is
\[
  a_{n+2} = \frac{n(n+1)-\ell(\ell+1)}{(n+1)(n+2)}\,a_n.
\]
When $\ell\in\N$, one series terminates: the resulting polynomial (suitably
normalized) is the Legendre polynomial $P_\ell(t)$.
For $\ell=2$: $a_2=-3a_0$, $a_4=0$, so
$P_2(t)=\frac{1}{2}(3t^2-1)$ (with normalization $P_2(1)=1$).
\end{example}

\begin{example}[Frobenius: Bessel equation of order $\nu$]\label{ex:bessel-nu}
\index{Bessel equation!order $\nu$}
The equation $t^2y''+ty'+(t^2-\nu^2)y=0$ has indicial equation
$r^2-\nu^2=0$, so $r=\pm\nu$.  For $r=\nu$ with $\nu\ge 0$:
\[
  J_\nu(t) = \sum_{n=0}^{\infty}\frac{(-1)^n}{n!\,\Gamma(n+\nu+1)}
    \left(\frac{t}{2}\right)^{2n+\nu}.
\]
When $2\nu\notin\Z$, a second solution is $J_{-\nu}(t)$.
When $\nu\in\Z$, the second solution requires a logarithmic term
(the Bessel function of the second kind $Y_\nu$).
\end{example}

\begin{example}[Variation of parameters for a system]\label{ex:var-params-system}
Solve $X' = \begin{pmatrix}0&1\\-1&0\end{pmatrix}X
+ \begin{pmatrix}0\\\sec t\end{pmatrix}$, $X(0)=0$.

The homogeneous solution gives
$\Phi(t)=\begin{pmatrix}\cos t&\sin t\\-\sin t&\cos t\end{pmatrix}$,
$\Phi^{-1}(t)=\begin{pmatrix}\cos t&-\sin t\\\sin t&\cos t\end{pmatrix}$.
Then
\[
  \Phi^{-1}(t)\begin{pmatrix}0\\\sec t\end{pmatrix}
  = \begin{pmatrix}-\tan t\\1\end{pmatrix}.
\]
Integrating from $0$ to $t$:
$\int_0^t\begin{pmatrix}-\tan s\\1\end{pmatrix}\dd s
= \begin{pmatrix}\ln\cos t\\t\end{pmatrix}$.

So
\[
  X(t) = \begin{pmatrix}\cos t&\sin t\\-\sin t&\cos t\end{pmatrix}
  \begin{pmatrix}\ln\cos t\\t\end{pmatrix}
  = \begin{pmatrix}\cos t\ln\cos t+t\sin t\\
    -\sin t\ln\cos t+t\cos t\end{pmatrix}.
\]
\end{example}

\begin{example}[Combining reduction of order with variation of parameters]%
\label{ex:combined-methods}
Solve $y''-\frac{2}{t}y'+\frac{2}{t^2}y=t$ for $t>0$.

\textbf{Step 1.}  The homogeneous equation $t^2y''-2ty'+2y=0$ is
Euler--Cauchy with indicial equation $r(r-1)-2r+2=r^2-3r+2=(r-1)(r-2)=0$.
So $y_1=t$, $y_2=t^2$.

\textbf{Step 2.}  Wronskian: $W=ty_2'-y_2\cdot 1 = 2t^2-t^2=t^2$.

\textbf{Step 3.}  Variation of parameters for $y''-\frac{2}{t}y'+\frac{2}{t^2}y=t$.
\begin{align*}
  u_1' &= -\frac{y_2\cdot g}{W} = -\frac{t^2\cdot t}{t^2} = -t,
  & u_1 &= -\frac{t^2}{2}, \\
  u_2' &= \frac{y_1\cdot g}{W} = \frac{t\cdot t}{t^2} = 1,
  & u_2 &= t.
\end{align*}
Particular solution:
$y_p = u_1 y_1 + u_2 y_2 = -\frac{t^2}{2}\cdot t + t\cdot t^2
= \frac{t^3}{2}$.

General solution: $y = c_1 t + c_2 t^2 + \frac{t^3}{2}$.
\end{example}

%--------------------------------------------------------------------
\section{Exercises}\label{sec:methods-exercises}
%--------------------------------------------------------------------

\begin{exercise}\label{exo:meth-1}
Compute the Wronskian of $e^t,\,te^t,\,t^2e^t$ and verify Abel's
theorem for the ODE they satisfy.
\end{exercise}

\begin{exercise}\label{exo:meth-2}
Use variation of parameters to solve $y''+4y=\tan(2t)$.
\end{exercise}

\begin{exercise}\label{exo:meth-3}
Find the Green's function for $y''-y=g(t)$ with $y(0)=y'(0)=0$
and use it to express the solution as an integral.
\end{exercise}

\begin{exercise}\label{exo:meth-4}
Find the power series solution of $(1+t^2)y''+2ty'-2y=0$ about $t_0=0$.
Determine the radius of convergence.
\end{exercise}

\begin{exercise}\label{exo:meth-5}
Use the Frobenius method to find one solution of
$ty''+y'+ty=0$ (Bessel equation of order zero, written differently).
\end{exercise}

\begin{exercise}\label{exo:meth-6}
Given that $y_1=\sin t/t$ solves $t^2y''+ty'+(t^2-1)y=0$ for $t>0$
(Bessel of order $1$, related form), use reduction of order to find a
second linearly independent solution.
\end{exercise}

\begin{exercise}\label{exo:meth-7}
Solve $y''+\frac{1}{t}y'-\frac{1}{t^2}y=\frac{1}{t}$ for $t>0$,
given that $y_1=t$ is a homogeneous solution.
\end{exercise}

\begin{exercise}\label{exo:meth-8}
Find the first four nonzero terms of two linearly independent power
series solutions of $(1-t)y''+y=0$ about $t_0=0$.
\end{exercise}

\begin{exercise}\label{exo:meth-9}
Show that the Wronskian of any two solutions of
$y''+(\sin t)y'+(\cos t)y=0$ satisfies $W(t)=W(0)\,e^{\cos t-1}$.
\end{exercise}

\begin{exercise}\label{exo:meth-10}
Use the Duhamel formula to solve
$X'=\begin{pmatrix}1&0\\0&-1\end{pmatrix}X+\begin{pmatrix}1\\t\end{pmatrix}$,
$X(0)=\begin{pmatrix}0\\0\end{pmatrix}$.
\end{exercise}

\begin{exercise}\label{exo:meth-11}
Classify the singular points of the equation
$t^2(t-1)^3y''+3t(t-1)y'+5y=0$
as regular or irregular.
\end{exercise}

\begin{exercise}\label{exo:meth-12}
Find the Green's function for the boundary value problem
$y''+y=g(t)$, $y(0)=0$, $y(\pi)=0$.
When does this problem have a solution?
\end{exercise}

%--------------------------------------------------------------------
\section{Chapter Summary}\label{sec:methods-summary}
%--------------------------------------------------------------------

\begin{itemize}
  \item The \textbf{Wronskian} of $n$ solutions of a linear ODE measures
    their linear independence.  Abel's theorem provides an explicit formula:
    $W(t)=W(t_0)\exp(-\int_{t_0}^t p_{n-1})$.
  \item \textbf{Variation of parameters} produces a particular solution of
    any non-homogeneous linear ODE given a fundamental system.  The key
    system is $\sum u_k' y_k^{(j)}=0$ for $j<n-1$ and
    $\sum u_k' y_k^{(n-1)}=g$.
  \item The \textbf{Green's function} $G(t,s)$ is an integral kernel
    that encodes the response of the system to a point source.  It
    satisfies the homogeneous equation with specific jump conditions.
  \item \textbf{Power series} solutions work at ordinary points; the
    radius of convergence is at least the distance to the nearest
    singular point.
  \item The \textbf{Frobenius method} handles regular singular points
    with solutions of the form $t^r\sum a_n t^n$.  The indicial equation
    determines the exponent $r$.
  \item \textbf{Reduction of order}: knowing one solution $y_1$ of a
    second-order equation, the formula
    $y_2=y_1\int e^{-\int P}/y_1^2$ gives a second independent solution.
  \item All these methods are complementary; the method summary table
    (Table~\ref{tab:method-summary}) provides a quick reference for
    choosing the appropriate technique.
\end{itemize}

% === END OF CHUNK ch5-6 ===
% === START OF CHUNK ch7-8 ===

% ══════════════════════════════════════════════════════════════
%  CHAPTER 7 — LAPLACE TRANSFORM AND APPLICATIONS
% ══════════════════════════════════════════════════════════════
\chapter{Laplace Transform and Applications}
\label{ch:laplace}
\index{Laplace transform|(}

\section{Motivating Example}
\label{sec:laplace-motivation}

Consider an \(RLC\) series circuit\index{RLC circuit} with resistance \(R\), inductance \(L\),
and capacitance \(C\), driven by a voltage source \(E(t)\) that switches on at \(t=0\)
and off at \(t=T\):
\[
  E(t) = \begin{cases} E_0 & \text{if } 0 \le t < T, \\ 0 & \text{if } t \ge T. \end{cases}
\]
The charge \(q(t)\) on the capacitor satisfies
\begin{equation}\label{eq:rlc-ode}
  L\,q'' + R\,q' + \frac{1}{C}\,q = E(t), \qquad q(0)=0,\; q'(0)=0.
\end{equation}
The forcing function \(E(t)\) is discontinuous, so the usual power series or
undetermined-coefficients methods do not apply directly. The \emph{Laplace
transform}\index{Laplace transform!motivation} converts this problem into an
algebraic equation in a new variable \(s\), handles discontinuities
effortlessly, and incorporates initial conditions automatically.

\section{Definition and Existence}
\label{sec:laplace-def}

\begin{definition}[Laplace transform]
\label{def:laplace}
\index{Laplace transform!definition}
Let \(f\colon [0,\infty)\to\R\) be piecewise continuous. The
\textbf{Laplace transform} of~\(f\) is
\begin{equation}\label{eq:laplace-def}
  \mathcal{L}\{f\}(s) = F(s) = \int_0^\infty e^{-st} f(t)\,\dd t,
\end{equation}
whenever the improper integral converges.
\end{definition}

\begin{definition}[Exponential order]
\label{def:exp-order}
\index{exponential order}
A function \(f\) is of \textbf{exponential order}~\(\alpha\) if there exist
constants \(M>0\) and \(T\ge 0\) such that
\[
  \abs{f(t)} \le M\,e^{\alpha t} \quad \text{for all } t \ge T.
\]
\end{definition}

\begin{theorem}[Existence of the Laplace transform]
\label{thm:laplace-existence}
\index{Laplace transform!existence}
If \(f\) is piecewise continuous on \([0,\infty)\) and of exponential
order~\(\alpha\), then \(\mathcal{L}\{f\}(s)\) exists for all \(s>\alpha\),
and moreover \(\abs{F(s)} \le M/(s-\alpha)\) for \(s>\alpha\).
\end{theorem}

\begin{proof}
For \(s>\alpha\) and \(t\ge T\), we have
\[
  \abs{e^{-st}f(t)} \le M\,e^{-st}\,e^{\alpha t} = M\,e^{-(s-\alpha)t}.
\]
The integral \(\int_T^\infty M\,e^{-(s-\alpha)t}\,\dd t = M/(s-\alpha)\)
converges. Since \(f\) is piecewise continuous, \(\int_0^T e^{-st}f(t)\,\dd t\)
is finite, so the full integral converges absolutely and satisfies the stated
bound.
\end{proof}

\section{Table of Fundamental Transforms}
\label{sec:laplace-table-derivations}
\index{Laplace transform!table}

We derive the most important transforms from the definition.

\begin{example}[Transform of \(1\)]
\label{ex:laplace-1}
For \(f(t)=1\) and \(s>0\):
\[
  \mathcal{L}\{1\}(s) = \int_0^\infty e^{-st}\,\dd t
  = \left[-\frac{e^{-st}}{s}\right]_0^\infty = \frac{1}{s}.
\]
\end{example}

\begin{example}[Transform of \(e^{at}\)]
\label{ex:laplace-exp}
\index{Laplace transform!of exponential}
For \(s>a\):
\[
  \mathcal{L}\{e^{at}\}(s) = \int_0^\infty e^{-(s-a)t}\,\dd t
  = \frac{1}{s-a}.
\]
\end{example}

\begin{example}[Transform of \(t^n\)]
\label{ex:laplace-tn}
\index{Laplace transform!of power function}
We claim \(\mathcal{L}\{t^n\}(s)=n!/s^{n+1}\) for \(n\in\N\) and \(s>0\).
For \(n=0\) this is \(1/s\). Integration by parts gives
\[
  \mathcal{L}\{t^n\} = \left[-\frac{t^n e^{-st}}{s}\right]_0^\infty
  + \frac{n}{s}\int_0^\infty t^{n-1}e^{-st}\,\dd t
  = \frac{n}{s}\,\mathcal{L}\{t^{n-1}\},
\]
and induction yields \(\mathcal{L}\{t^n\}=n!/s^{n+1}\).
\end{example}

\begin{example}[Transform of \(\sin(bt)\) and \(\cos(bt)\)]
\label{ex:laplace-sincos}
\index{Laplace transform!of sine}\index{Laplace transform!of cosine}
Since \(e^{ibt}=\cos(bt)+i\sin(bt)\) and
\(\mathcal{L}\{e^{ibt}\}=1/(s-ib)\), we compute
\[
  \frac{1}{s-ib} = \frac{s+ib}{s^2+b^2}.
\]
Taking real and imaginary parts:
\[
  \mathcal{L}\{\cos(bt)\} = \frac{s}{s^2+b^2}, \qquad
  \mathcal{L}\{\sin(bt)\} = \frac{b}{s^2+b^2}.
\]
\end{example}

\section{Properties of the Laplace Transform}
\label{sec:laplace-properties}

\begin{theorem}[Linearity]
\label{thm:laplace-linear}
\index{Laplace transform!linearity}
For constants \(a,b\) and functions \(f,g\) with Laplace transforms:
\[
  \mathcal{L}\{af+bg\} = a\,\mathcal{L}\{f\} + b\,\mathcal{L}\{g\}.
\]
\end{theorem}

\begin{proof}
Immediate from linearity of the integral.
\end{proof}

\begin{theorem}[First shifting theorem (\(s\)-shift)]
\label{thm:s-shift}
\index{Laplace transform!s-shift@$s$-shift}
If \(\mathcal{L}\{f\}(s)=F(s)\) for \(s>\alpha\), then
\begin{equation}\label{eq:s-shift}
  \mathcal{L}\{e^{at}f(t)\}(s) = F(s-a) \quad \text{for } s>\alpha+a.
\end{equation}
\end{theorem}

\begin{proof}
\[
  \mathcal{L}\{e^{at}f(t)\}(s)
  = \int_0^\infty e^{-st}\,e^{at}\,f(t)\,\dd t
  = \int_0^\infty e^{-(s-a)t}f(t)\,\dd t
  = F(s-a). \qedhere
\]
\end{proof}

\begin{example}[Application of the \(s\)-shift]
\label{ex:s-shift-app}
\[
  \mathcal{L}\{e^{-3t}\sin(2t)\}
  = \frac{2}{(s+3)^2+4}.
\]
\end{example}

\begin{theorem}[Second shifting theorem (\(t\)-shift)]
\label{thm:t-shift}
\index{Laplace transform!t-shift@$t$-shift}\index{Heaviside step function}
Let \(u_c(t)\) denote the Heaviside step function\index{Heaviside step function}
(see Definition~\ref{def:heaviside}). If \(\mathcal{L}\{f\}=F(s)\), then
\begin{equation}\label{eq:t-shift}
  \mathcal{L}\{u_c(t)\,f(t-c)\}(s) = e^{-cs}\,F(s) \quad (c\ge 0).
\end{equation}
\end{theorem}

\begin{proof}
Substituting \(\tau=t-c\):
\[
  \int_0^\infty e^{-st}\,u_c(t)\,f(t-c)\,\dd t
  = \int_c^\infty e^{-st}\,f(t-c)\,\dd t
  = \int_0^\infty e^{-s(\tau+c)}f(\tau)\,\dd\tau
  = e^{-cs}\,F(s). \qedhere
\]
\end{proof}

\begin{theorem}[Laplace transform of derivatives]
\label{thm:laplace-deriv}
\index{Laplace transform!of derivatives}
Let \(f\) be continuous, \(f'\) piecewise continuous, both of exponential order.
Then
\begin{equation}\label{eq:laplace-deriv1}
  \mathcal{L}\{f'\}(s) = s\,F(s) - f(0).
\end{equation}
More generally, if \(f,f',\ldots,f^{(n-1)}\) are continuous and
\(f^{(n)}\) is piecewise continuous, all of exponential order, then
\begin{equation}\label{eq:laplace-deriv-n}
  \mathcal{L}\{f^{(n)}\}(s) = s^n F(s) - s^{n-1}f(0) - s^{n-2}f'(0)
  - \cdots - f^{(n-1)}(0).
\end{equation}
\end{theorem}

\begin{proof}
For the first derivative, integration by parts gives
\[
  \int_0^\infty e^{-st}f'(t)\,\dd t
  = \bigl[e^{-st}f(t)\bigr]_0^\infty + s\int_0^\infty e^{-st}f(t)\,\dd t
  = -f(0) + s\,F(s).
\]
The boundary term at \(\infty\) vanishes by the exponential order assumption.
The general formula follows by induction.
\end{proof}

\begin{theorem}[Multiplication by \(t^n\)]
\label{thm:laplace-mult-t}
\index{Laplace transform!multiplication by $t$}
If \(\mathcal{L}\{f\}=F(s)\), then
\begin{equation}\label{eq:laplace-mult-t}
  \mathcal{L}\{t^n f(t)\}(s) = (-1)^n F^{(n)}(s).
\end{equation}
\end{theorem}

\begin{proof}
Differentiating under the integral sign:
\[
  F'(s) = \frac{\dd}{\dd s}\int_0^\infty e^{-st}f(t)\,\dd t
  = \int_0^\infty (-t)\,e^{-st}f(t)\,\dd t
  = -\mathcal{L}\{tf(t)\}.
\]
The general formula follows by iterating.
\end{proof}

\begin{theorem}[Convolution theorem]
\label{thm:convolution}
\index{convolution}\index{Laplace transform!convolution}
Define the \textbf{convolution}\index{convolution!definition} of \(f\) and \(g\) by
\begin{equation}\label{eq:convolution-def}
  (f*g)(t) = \int_0^t f(\tau)\,g(t-\tau)\,\dd\tau.
\end{equation}
If \(\mathcal{L}\{f\}=F\) and \(\mathcal{L}\{g\}=G\), then
\begin{equation}\label{eq:convolution-thm}
  \mathcal{L}\{f*g\}(s) = F(s)\,G(s).
\end{equation}
\end{theorem}

\begin{proof}
We write
\[
  F(s)\,G(s)
  = \int_0^\infty e^{-s\sigma}f(\sigma)\,\dd\sigma
    \int_0^\infty e^{-s\tau}g(\tau)\,\dd\tau
  = \int_0^\infty\!\int_0^\infty e^{-s(\sigma+\tau)}f(\sigma)\,g(\tau)
    \,\dd\sigma\,\dd\tau.
\]
Substituting \(t=\sigma+\tau\) (so \(\sigma=t-\tau\)) and changing the order
of integration (justified by absolute convergence), the region
\(\sigma\ge 0,\;\tau\ge 0\) becomes \(t\ge 0,\;0\le\tau\le t\):
\[
  F(s)\,G(s) = \int_0^\infty e^{-st}\left(\int_0^t f(t-\tau)\,g(\tau)
  \,\dd\tau\right)\dd t = \mathcal{L}\{f*g\}(s). \qedhere
\]
\end{proof}

\begin{remark}
\label{rem:convolution-properties}
Convolution is commutative (\(f*g=g*f\)), associative, and distributes over
addition, but \(f*1\neq f\) in general (indeed \((f*1)(t)=\int_0^t f(\tau)\,\dd\tau\)).
\end{remark}

\section{Heaviside and Dirac Delta Functions}
\label{sec:heaviside-dirac}

\begin{definition}[Heaviside step function]
\label{def:heaviside}
\index{Heaviside step function!definition}
For \(c\ge 0\), the \textbf{Heaviside step function} is
\[
  u_c(t) = \begin{cases} 0 & \text{if } t<c, \\ 1 & \text{if } t\ge c. \end{cases}
\]
\end{definition}

\begin{proposition}[Transform of \(u_c\)]
\label{prop:laplace-heaviside}
For \(c\ge 0\) and \(s>0\):
\[
  \mathcal{L}\{u_c(t)\}(s) = \frac{e^{-cs}}{s}.
\]
\end{proposition}

\begin{proof}
\(\displaystyle\int_c^\infty e^{-st}\,\dd t = e^{-cs}/s.\)
\end{proof}

\begin{remark}[Piecewise functions via Heaviside]
\label{rem:piecewise-heaviside}
\index{Heaviside step function!piecewise functions}
Any piecewise-defined function can be written using Heaviside functions. For
example, the forcing \(E(t)\) from \eqref{eq:rlc-ode} is
\(E(t)=E_0\bigl(u_0(t)-u_T(t)\bigr)\), so
\[
  \mathcal{L}\{E\}(s) = E_0\left(\frac{1}{s}-\frac{e^{-Ts}}{s}\right)
  = \frac{E_0(1-e^{-Ts})}{s}.
\]
\end{remark}

\begin{definition}[Dirac delta function]
\label{def:dirac}
\index{Dirac delta function}
The \textbf{Dirac delta function} \(\delta(t-c)\) (for \(c\ge 0\)) is
informally a ``function'' satisfying \(\delta(t-c)=0\) for \(t\neq c\) and
\(\int_{-\infty}^{\infty}\delta(t-c)\,\varphi(t)\,\dd t = \varphi(c)\) for
any continuous~\(\varphi\). Rigorously, it is a distribution.
\end{definition}

\begin{proposition}[Transform of \(\delta(t-c)\)]
\label{prop:laplace-dirac}
\index{Dirac delta function!Laplace transform}
For \(c\ge 0\):
\[
  \mathcal{L}\{\delta(t-c)\}(s) = e^{-cs}.
\]
In particular, \(\mathcal{L}\{\delta(t)\}=1\).
\end{proposition}

\section{Inverse Laplace Transform and Partial Fractions}
\label{sec:inverse-laplace}
\index{Laplace transform!inverse}\index{partial fractions}

The inverse Laplace transform \(\mathcal{L}^{-1}\{F\}=f\) is determined
(for continuous functions) by uniqueness. In practice, we decompose rational
functions of~\(s\) using partial fractions and read off inverse transforms
from a table.

\begin{method}[Partial fraction decomposition]
\label{meth:partial-fractions}
\index{partial fractions!method}
Given \(F(s)=P(s)/Q(s)\) with \(\deg P<\deg Q\):
\begin{enumerate}[label=(\roman*)]
  \item Factor \(Q(s)\) into linear and irreducible quadratic factors.
  \item For each simple real root \(s=a\): term \(A/(s-a)\).
  \item For a repeated root \(s=a\) of multiplicity~\(m\):
    terms \(A_1/(s-a)+A_2/(s-a)^2+\cdots+A_m/(s-a)^m\).
  \item For an irreducible quadratic \((s-\alpha)^2+\beta^2\):
    term \((As+B)/\bigl((s-\alpha)^2+\beta^2\bigr)\).
  \item Solve for coefficients by equating numerators.
\end{enumerate}
\end{method}

\begin{example}[Inverse Laplace by partial fractions]
\label{ex:inverse-laplace-pf}
Find \(\mathcal{L}^{-1}\!\left\{\dfrac{3s+7}{(s+1)(s^2+4)}\right\}\).

\medskip\noindent\textbf{Solution.}
Write
\[
  \frac{3s+7}{(s+1)(s^2+4)} = \frac{A}{s+1} + \frac{Bs+C}{s^2+4}.
\]
Multiplying through: \(3s+7 = A(s^2+4)+(Bs+C)(s+1)\).
Setting \(s=-1\): \(4=5A\), so \(A=4/5\).
Comparing \(s^2\): \(0=A+B\), so \(B=-4/5\).
Comparing constant terms: \(7=4A+C\), so \(C=7-16/5=19/5\).
Thus
\[
  \mathcal{L}^{-1}\!\left\{\frac{3s+7}{(s+1)(s^2+4)}\right\}
  = \frac{4}{5}\,e^{-t} - \frac{4}{5}\cos(2t) + \frac{19}{10}\sin(2t).
\]
\end{example}

\section{Solving ODEs with the Laplace Transform}
\label{sec:laplace-solve-ode}
\index{Laplace transform!solving ODEs}

\begin{method}[Laplace transform method for IVPs]
\label{meth:laplace-ivp}
To solve an initial value problem:
\begin{enumerate}[label=\textbf{Step~\arabic*.}]
  \item Apply \(\mathcal{L}\) to both sides using linearity and the derivative
    rule (Theorem~\ref{thm:laplace-deriv}).
  \item Substitute initial conditions \(y(0),y'(0),\ldots\)
  \item Solve algebraically for \(Y(s)=\mathcal{L}\{y\}(s)\).
  \item Find \(y(t)=\mathcal{L}^{-1}\{Y\}\) using partial fractions and the
    table.
\end{enumerate}
\end{method}

\begin{example}[Second-order IVP]
\label{ex:laplace-2nd-order}
Solve \(y''-3y'+2y=e^{3t}\), \(y(0)=1\), \(y'(0)=0\).

\medskip\noindent\textbf{Solution.}
Applying \(\mathcal{L}\):
\[
  s^2 Y - s\cdot 1 - 0 - 3(sY - 1) + 2Y = \frac{1}{s-3}.
\]
Simplifying: \((s^2-3s+2)Y = s - 3 + \frac{1}{s-3}\), i.e.,
\[
  (s-1)(s-2)\,Y = \frac{(s-3)^2+1}{s-3} = \frac{s^2-6s+10}{s-3}.
\]
Partial fractions of \(Y(s)=\dfrac{s^2-6s+10}{(s-1)(s-2)(s-3)}\):
\[
  \frac{s^2-6s+10}{(s-1)(s-2)(s-3)}
  = \frac{A}{s-1}+\frac{B}{s-2}+\frac{C}{s-3}.
\]
Cover-up: \(A=\frac{1-6+10}{(-1)(-2)}=\frac{5}{2}\),
\(B=\frac{4-12+10}{(1)(-1)}=-2\),
\(C=\frac{9-18+10}{(2)(1)}=\frac{1}{2}\).
Therefore
\[
  y(t) = \frac{5}{2}\,e^t - 2\,e^{2t} + \frac{1}{2}\,e^{3t}.
\]
\end{example}

\begin{example}[ODE with discontinuous forcing]
\label{ex:laplace-discontinuous}
\index{Laplace transform!discontinuous forcing}
Solve \(y''+y=u_\pi(t)\), \(y(0)=0\), \(y'(0)=1\).

\medskip\noindent\textbf{Solution.}
Taking transforms:
\[
  s^2Y - 1 + Y = \frac{e^{-\pi s}}{s},
  \quad\text{so}\quad
  Y = \frac{1}{s^2+1} + \frac{e^{-\pi s}}{s(s^2+1)}.
\]
Partial fractions: \(\frac{1}{s(s^2+1)}=\frac{1}{s}-\frac{s}{s^2+1}\). By
the \(t\)-shift theorem:
\[
  y(t) = \sin t + u_\pi(t)\bigl[1-\cos(t-\pi)\bigr]
       = \sin t + u_\pi(t)\bigl[1+\cos t\bigr].
\]
\end{example}

\begin{example}[Impulse response via Dirac delta]
\label{ex:laplace-impulse}
\index{Dirac delta function!impulse response}
Solve \(y''+4y=3\,\delta(t-2)\), \(y(0)=0\), \(y'(0)=0\).

\medskip\noindent\textbf{Solution.}
\[
  s^2Y+4Y=3e^{-2s}, \quad Y=\frac{3e^{-2s}}{s^2+4}.
\]
Since \(\mathcal{L}^{-1}\{1/(s^2+4)\}=\frac{1}{2}\sin(2t)\), the
\(t\)-shift gives
\[
  y(t) = \frac{3}{2}\,u_2(t)\sin\bigl(2(t-2)\bigr).
\]
\end{example}

\section{Systems and Periodic Forcing}
\label{sec:laplace-systems}

\begin{example}[System of ODEs]
\label{ex:laplace-system}
\index{Laplace transform!systems}
Solve the system
\[
  x' = 2x + y, \quad y' = 3x + 4y, \qquad x(0)=1,\; y(0)=0.
\]
Taking Laplace transforms:
\begin{align*}
  sX - 1 &= 2X + Y, \\
  sY &= 3X + 4Y.
\end{align*}
From the second equation: \(Y=3X/(s-4)\). Substituting into the first:
\[
  (s-2)X - 1 = \frac{3X}{s-4},
  \quad\Longrightarrow\quad
  X\left[(s-2)-\frac{3}{s-4}\right]=1.
\]
Thus \(X=\frac{s-4}{s^2-6s+5}=\frac{s-4}{(s-1)(s-5)}\). Partial fractions:
\[
  X = \frac{3/4}{s-1}+\frac{1/4}{s-5},
  \qquad
  Y = \frac{3X}{s-4} = \frac{-3/4}{s-1}+\frac{3/4}{s-5}.
\]
Therefore \(x(t)=\frac{3}{4}e^t+\frac{1}{4}e^{5t}\) and
\(y(t)=-\frac{3}{4}e^t+\frac{3}{4}e^{5t}\).
\end{example}

\begin{theorem}[Transform of periodic functions]
\label{thm:laplace-periodic}
\index{Laplace transform!periodic functions}
If \(f\) has period~\(T>0\) (i.e., \(f(t+T)=f(t)\) for all \(t\ge 0\)), then
\begin{equation}\label{eq:laplace-periodic}
  \mathcal{L}\{f\}(s) = \frac{1}{1-e^{-sT}}\int_0^T e^{-st}f(t)\,\dd t.
\end{equation}
\end{theorem}

\begin{proof}
Split the integral into intervals of length~\(T\):
\[
  F(s)=\sum_{n=0}^\infty \int_{nT}^{(n+1)T} e^{-st}f(t)\,\dd t.
\]
In each term, substitute \(\tau=t-nT\), using \(f(\tau+nT)=f(\tau)\):
\[
  F(s)=\sum_{n=0}^\infty e^{-nsT}\int_0^T e^{-s\tau}f(\tau)\,\dd\tau
  = \frac{1}{1-e^{-sT}}\int_0^T e^{-s\tau}f(\tau)\,\dd\tau. \qedhere
\]
\end{proof}

\begin{example}[Square wave]
\label{ex:square-wave}
\index{square wave}
The square wave of period~\(2a\) is \(f(t)=1\) for \(0\le t<a\) and
\(f(t)=-1\) for \(a\le t<2a\). Then
\[
  \int_0^{2a} e^{-st}f(t)\,\dd t
  = \frac{1-e^{-as}}{s} - \frac{e^{-as}-e^{-2as}}{s}
  = \frac{(1-e^{-as})^2}{s}.
\]
By Theorem~\ref{thm:laplace-periodic}:
\[
  \mathcal{L}\{f\}(s) = \frac{(1-e^{-as})^2}{s(1-e^{-2as})}
  = \frac{1-e^{-as}}{s(1+e^{-as})}
  = \frac{1}{s}\tanh\!\left(\frac{as}{2}\right).
\]
\end{example}

\section{Complete Reference Table}
\label{sec:laplace-reference-table}
\index{Laplace transform!reference table}

\begin{center}
\renewcommand{\arraystretch}{1.4}
\begin{tabular}{@{}lll@{}}
\toprule
\(f(t)\) \((t\ge 0)\) & \(F(s)=\mathcal{L}\{f\}(s)\) & Domain \\
\midrule
\(1\) & \(1/s\) & \(s>0\) \\
\(t^n\) \((n\in\N)\) & \(n!/s^{n+1}\) & \(s>0\) \\
\(e^{at}\) & \(1/(s-a)\) & \(s>a\) \\
\(\sin(bt)\) & \(b/(s^2+b^2)\) & \(s>0\) \\
\(\cos(bt)\) & \(s/(s^2+b^2)\) & \(s>0\) \\
\(\sinh(bt)\) & \(b/(s^2-b^2)\) & \(s>\abs{b}\) \\
\(\cosh(bt)\) & \(s/(s^2-b^2)\) & \(s>\abs{b}\) \\
\(t^n e^{at}\) & \(n!/(s-a)^{n+1}\) & \(s>a\) \\
\(e^{at}\sin(bt)\) & \(b/\bigl((s-a)^2+b^2\bigr)\) & \(s>a\) \\
\(e^{at}\cos(bt)\) & \((s-a)/\bigl((s-a)^2+b^2\bigr)\) & \(s>a\) \\
\(t\sin(bt)\) & \(2bs/(s^2+b^2)^2\) & \(s>0\) \\
\(t\cos(bt)\) & \((s^2-b^2)/(s^2+b^2)^2\) & \(s>0\) \\
\(u_c(t)\) & \(e^{-cs}/s\) & \(s>0\) \\
\(\delta(t-c)\) & \(e^{-cs}\) & all \(s\) \\
\(u_c(t)\,f(t-c)\) & \(e^{-cs}F(s)\) & --- \\
\(e^{at}f(t)\) & \(F(s-a)\) & --- \\
\(t^n f(t)\) & \((-1)^n F^{(n)}(s)\) & --- \\
\((f*g)(t)\) & \(F(s)\,G(s)\) & --- \\
\(f'(t)\) & \(sF(s)-f(0)\) & --- \\
\(f''(t)\) & \(s^2F(s)-sf(0)-f'(0)\) & --- \\
\bottomrule
\end{tabular}
\end{center}

\section{Exercises}
\label{sec:laplace-exercises}

\begin{exercise}
\label{exo:laplace-compute}
Compute the Laplace transforms of: (a)~\(f(t)=3t^2-2e^{-t}+5\sin(4t)\);
(b)~\(g(t)=t\,e^{2t}\cos(3t)\); (c)~\(h(t)=u_2(t)\,(t-2)^3\,e^{-(t-2)}\).
\end{exercise}

\begin{exercise}
\label{exo:laplace-inverse}
Find the inverse Laplace transforms of:
(a)~\(\dfrac{2s+5}{s^2+6s+13}\); \quad
(b)~\(\dfrac{s^2}{(s+1)^3}\); \quad
(c)~\(\dfrac{e^{-3s}}{s^2(s+2)}\).
\end{exercise}

\begin{exercise}
\label{exo:laplace-ivp1}
Use the Laplace transform to solve
\(y''+2y'+5y=4e^{-t}\sin(2t)\), \(y(0)=1\), \(y'(0)=-1\).
\end{exercise}

\begin{exercise}
\label{exo:laplace-ivp2}
Solve \(y''-y=u_1(t)-u_3(t)\), \(y(0)=0\), \(y'(0)=0\).
Express the solution in piecewise form.
\end{exercise}

\begin{exercise}
\label{exo:laplace-convolution}
Use the convolution theorem to find
\(\mathcal{L}^{-1}\!\left\{\dfrac{1}{s^2(s^2+1)}\right\}\).
\end{exercise}

\begin{exercise}
\label{exo:laplace-system}
Use Laplace transforms to solve the system
\[
  x' = x - 2y, \quad y' = 5x - y, \qquad x(0)=2,\;y(0)=1.
\]
\end{exercise}

\begin{exercise}
\label{exo:laplace-periodic}
Find the Laplace transform of the sawtooth wave\index{sawtooth wave}
\(f(t)=t\) for \(0\le t<1\) with period~1.
\end{exercise}

\begin{exercise}
\label{exo:laplace-impulse}
A spring-mass system satisfies \(y''+\omega^2 y = \sum_{k=0}^{N}\delta(t-k\pi/\omega)\)
with \(y(0)=y'(0)=0\). Show that
\(y(t)=\frac{1}{\omega}\sum_{k=0}^{N} u_{k\pi/\omega}(t)\sin\bigl(\omega(t-k\pi/\omega)\bigr)\)
and discuss resonance when \(N\to\infty\).
\end{exercise}

\section{Chapter Summary}
\label{sec:laplace-summary}

\begin{itemize}
  \item The Laplace transform converts an ODE initial value problem into an
    algebraic equation in~\(s\), incorporating initial conditions automatically.
  \item Key properties---linearity, shifting, derivative rule, convolution---form
    a toolkit for transforming and inverting.
  \item The Heaviside function and Dirac delta allow systematic treatment of
    discontinuous and impulsive forcing.
  \item Partial fraction decomposition is the principal technique for computing
    inverse transforms.
  \item The method extends naturally to systems and to periodic forcing via
    the periodic transform formula.
\end{itemize}

\index{Laplace transform|)}


% ══════════════════════════════════════════════════════════════
%  CHAPTER 8 — STABILITY OF SOLUTIONS AND EQUILIBRIA
% ══════════════════════════════════════════════════════════════
\chapter{Stability of Solutions and Equilibria}
\label{ch:stability}
\index{stability|(}

\section{Motivating Example: The Pendulum}
\label{sec:stability-motivation}

Consider a rigid pendulum\index{pendulum} of length~\(\ell\) and mass~\(m\)
with angular displacement~\(\theta\) from the downward vertical, subject to
gravity and damping:
\begin{equation}\label{eq:pendulum}
  m\ell\,\theta'' + c\,\theta' + mg\sin\theta = 0,
\end{equation}
or equivalently the system
\[
  \theta' = \omega, \qquad
  \omega' = -\frac{g}{\ell}\sin\theta - \frac{c}{m\ell}\,\omega.
\]
This system has equilibria at \((\theta,\omega)=(n\pi,0)\) for \(n\in\Z\).
Physical intuition tells us the hanging position \(\theta=0\) is
\emph{stable} while the inverted position \(\theta=\pi\) is \emph{unstable}.
How can we make these notions precise, and prove them rigorously?

\begin{center}
\begin{tikzpicture}[>=Stealth,scale=1.0]
  % Pivot
  \fill (0,0) circle (2pt);
  \draw[thick] (-1,0) -- (1,0);
  \foreach \x in {-0.8,-0.5,-0.2,0.1,0.4,0.7}{
    \draw (\x,0) -- (\x+0.2,0.25);
  }
  % Pendulum
  \draw[thick] (0,0) -- (210:3) node[midway,left]{\(\ell\)};
  \fill (210:3) circle (5pt) node[below left]{\(m\)};
  % Angle
  \draw[dashed] (0,0) -- (0,-2.5);
  \draw[->] (0,-1.5) arc[start angle=-90,end angle=210,radius=1.5]
    node[midway,left]{\(\theta\)};
  % Gravity
  \draw[->,thick,blue] (210:3) -- ++(0,-1) node[right]{\(mg\)};
\end{tikzpicture}
\end{center}

\section{Autonomous Systems and Equilibrium Points}
\label{sec:equilibria}
\index{autonomous system}\index{equilibrium point}

Throughout this chapter we consider the autonomous system
\begin{equation}\label{eq:autonomous}
  \mathbf{x}' = \mathbf{f}(\mathbf{x}),
  \qquad \mathbf{x}\in\R^n,
\end{equation}
where \(\mathbf{f}\colon D\to\R^n\) is continuously differentiable on an
open set \(D\subseteq\R^n\).

\begin{definition}[Equilibrium point]
\label{def:equilibrium}
\index{equilibrium point!definition}
A point \(\mathbf{x}^*\in D\) is an \textbf{equilibrium point}
(or \textbf{critical point}, \textbf{fixed point}, \textbf{stationary point})
of~\eqref{eq:autonomous} if \(\mathbf{f}(\mathbf{x}^*)=\mathbf{0}\).
\end{definition}

\begin{remark}
\label{rem:constant-solution}
If \(\mathbf{x}^*\) is an equilibrium, then \(\mathbf{x}(t)\equiv\mathbf{x}^*\)
is a solution for all~\(t\). The question of stability concerns the behavior of
\emph{other} solutions that start near~\(\mathbf{x}^*\).
\end{remark}

\begin{definition}[Classification of equilibria for planar linear systems]
\label{def:equilibrium-classification}
\index{equilibrium point!classification}
For the linear system \(\mathbf{x}'=A\mathbf{x}\) with \(A\in\R^{2\times 2}\)
and eigenvalues \(\lambda_1,\lambda_2\), the equilibrium at the origin is
classified as:
\begin{enumerate}[label=(\roman*)]
  \item \textbf{Stable node}\index{node!stable}: \(\lambda_1,\lambda_2\) real,
    both negative.
  \item \textbf{Unstable node}\index{node!unstable}: \(\lambda_1,\lambda_2\)
    real, both positive.
  \item \textbf{Saddle point}\index{saddle point}: \(\lambda_1,\lambda_2\) real,
    opposite signs.
  \item \textbf{Stable spiral}\index{spiral!stable} (focus): complex eigenvalues
    with negative real part.
  \item \textbf{Unstable spiral}\index{spiral!unstable} (focus): complex
    eigenvalues with positive real part.
  \item \textbf{Center}\index{center}: purely imaginary eigenvalues.
  \item \textbf{Degenerate cases}: repeated eigenvalues (star node or
    improper node).
\end{enumerate}
\end{definition}

\section{Lyapunov Stability}
\label{sec:lyapunov-stability}
\index{Lyapunov stability|(}

\begin{definition}[Lyapunov stability]
\label{def:lyapunov-stability}
\index{stability!Lyapunov}
The equilibrium \(\mathbf{x}^*\) of~\eqref{eq:autonomous} is:
\begin{enumerate}[label=(\roman*)]
  \item \textbf{Stable (in the sense of Lyapunov)}\index{stability!Lyapunov}
    if for every \(\varepsilon>0\) there exists \(\delta>0\) such that
    \[
      \norm{\mathbf{x}(0)-\mathbf{x}^*}<\delta
      \implies \norm{\mathbf{x}(t)-\mathbf{x}^*}<\varepsilon
      \quad\text{for all } t\ge 0.
    \]
  \item \textbf{Asymptotically stable}\index{stability!asymptotic} if it is
    stable and moreover there exists \(\delta_0>0\) such that
    \[
      \norm{\mathbf{x}(0)-\mathbf{x}^*}<\delta_0
      \implies \lim_{t\to\infty}\mathbf{x}(t)=\mathbf{x}^*.
    \]
  \item \textbf{Unstable}\index{instability} if it is not stable.
\end{enumerate}
\end{definition}

\begin{remark}
\label{rem:stability-shift}
Without loss of generality, we may assume \(\mathbf{x}^*=\mathbf{0}\) by
the change of variables \(\mathbf{y}=\mathbf{x}-\mathbf{x}^*\). We adopt this
convention for the remainder of this chapter.
\end{remark}

\begin{example}[Linear stability]
\label{ex:linear-stability}
For the linear system \(\mathbf{x}'=A\mathbf{x}\), the origin is:
\begin{itemize}
  \item asymptotically stable if and only if all eigenvalues of~\(A\) have
    strictly negative real parts;
  \item stable if and only if all eigenvalues have non-positive real parts, and
    every eigenvalue with zero real part has equal algebraic and geometric
    multiplicity;
  \item unstable otherwise.
\end{itemize}
\end{example}

\section{Lyapunov's Direct Method}
\label{sec:lyapunov-direct}
\index{Lyapunov!direct method|(}

Lyapunov's direct (or second) method establishes stability without explicitly
solving the ODE. The idea is to find a generalized ``energy'' function that
decreases along trajectories.

\begin{definition}[Lyapunov function]
\label{def:lyapunov-function}
\index{Lyapunov function}
Let \(\mathbf{0}\) be an equilibrium of \(\mathbf{x}'=\mathbf{f}(\mathbf{x})\).
A continuously differentiable function \(V\colon D\to\R\) (with
\(\mathbf{0}\in D\)) is a \textbf{Lyapunov function} if:
\begin{enumerate}[label=(\roman*)]
  \item \(V(\mathbf{0})=0\);
  \item \(V(\mathbf{x})>0\) for all \(\mathbf{x}\in D\setminus\{\mathbf{0}\}\)
    (i.e., \(V\) is \emph{positive definite}\index{positive definite});
  \item \(\dot{V}(\mathbf{x}) \le 0\) for all \(\mathbf{x}\in D\), where the
    \textbf{orbital derivative}\index{orbital derivative} is
    \begin{equation}\label{eq:orbital-deriv}
      \dot{V}(\mathbf{x}) = \nabla V(\mathbf{x})\cdot\mathbf{f}(\mathbf{x})
      = \sum_{i=1}^n \frac{\partial V}{\partial x_i}\,f_i(\mathbf{x}).
    \end{equation}
\end{enumerate}
If moreover \(\dot{V}(\mathbf{x})<0\) for all \(\mathbf{x}\in D\setminus\{\mathbf{0}\}\),
we say \(V\) is a \textbf{strict Lyapunov function}\index{Lyapunov function!strict}.
\end{definition}

\begin{remark}
\label{rem:orbital-deriv}
The orbital derivative \(\dot{V}(\mathbf{x})\) gives the rate of change of~\(V\)
along solutions: if \(\mathbf{x}(t)\) is a solution, then
\(\frac{\dd}{\dd t}V(\mathbf{x}(t))=\dot{V}(\mathbf{x}(t))\) by the chain rule.
Hence \(\dot{V}\le 0\) means \(V\) is non-increasing along trajectories.
\end{remark}

\begin{theorem}[Lyapunov stability theorem]
\label{thm:lyapunov-stability}
\index{Lyapunov!stability theorem}
Let \(\mathbf{0}\) be an equilibrium of \(\mathbf{x}'=\mathbf{f}(\mathbf{x})\).
If there exists a Lyapunov function \(V\) on a neighborhood~\(D\) of the
origin (with \(\dot{V}\le 0\)), then the origin is \textbf{stable}.
\end{theorem}

\begin{proof}
Let \(\varepsilon>0\) be small enough that the closed ball
\(\overline{B}_\varepsilon=\{\mathbf{x}:\norm{\mathbf{x}}\le\varepsilon\}\subseteq D\).
Since \(V\) is continuous, positive definite, and \(V(\mathbf{0})=0\), define
\[
  m = \min_{\norm{\mathbf{x}}=\varepsilon} V(\mathbf{x}) > 0.
\]
By continuity of~\(V\) at the origin, there exists \(\delta\in(0,\varepsilon)\)
such that
\[
  \norm{\mathbf{x}}<\delta \implies V(\mathbf{x}) < m.
\]
Now suppose \(\norm{\mathbf{x}(0)}<\delta\). Then \(V(\mathbf{x}(0))<m\).
Since \(\dot{V}\le 0\) along trajectories, we have
\(V(\mathbf{x}(t))\le V(\mathbf{x}(0))<m\) for all \(t\ge 0\). But if
\(\mathbf{x}(t)\) ever reached the sphere \(\norm{\mathbf{x}}=\varepsilon\),
we would have \(V(\mathbf{x}(t))\ge m\), a contradiction. Therefore
\(\norm{\mathbf{x}(t)}<\varepsilon\) for all \(t\ge 0\).
\end{proof}

\begin{theorem}[Asymptotic stability theorem]
\label{thm:lyapunov-asymptotic}
\index{Lyapunov!asymptotic stability theorem}
Under the hypotheses of Theorem~\ref{thm:lyapunov-stability}, if moreover
\(\dot{V}(\mathbf{x})<0\) for all \(\mathbf{x}\in D\setminus\{\mathbf{0}\}\),
then the origin is \textbf{asymptotically stable}.
\end{theorem}

\begin{proof}
Stability is given by Theorem~\ref{thm:lyapunov-stability}. Choose \(\delta\)
as in that proof, and let \(\norm{\mathbf{x}(0)}<\delta\). The function
\(t\mapsto V(\mathbf{x}(t))\) is non-increasing and bounded below by~\(0\),
so it converges to some limit \(L\ge 0\). We claim \(L=0\).

Suppose \(L>0\). Then \(\mathbf{x}(t)\) remains in the compact annular region
\(K=\{\mathbf{x}: L\le V(\mathbf{x})\le V(\mathbf{x}(0))\}\subseteq
\overline{B}_\varepsilon\setminus\{\mathbf{0}\}\). On this compact set, the
continuous function \(\dot{V}\) achieves a maximum \(-\mu<0\). Therefore
\[
  V(\mathbf{x}(t)) \le V(\mathbf{x}(0)) - \mu t \to -\infty
  \quad\text{as } t\to\infty,
\]
contradicting \(V\ge 0\). Hence \(L=0\), and since \(V\) is positive definite,
\(\mathbf{x}(t)\to\mathbf{0}\).
\end{proof}

\begin{theorem}[Lyapunov instability theorem]
\label{thm:lyapunov-instability}
\index{Lyapunov!instability theorem}
Let \(\mathbf{0}\) be an equilibrium. Suppose there exists a continuously
differentiable function \(V\) on a neighborhood of the origin with
\(V(\mathbf{0})=0\) and \(\dot{V}(\mathbf{x})>0\) for \(\mathbf{x}\neq\mathbf{0}\).
If every neighborhood of the origin contains a point where \(V>0\), then the
origin is \textbf{unstable}.
\end{theorem}

\begin{example}[Damped nonlinear pendulum]
\label{ex:pendulum-lyapunov}
\index{pendulum!Lyapunov analysis}
For the pendulum \eqref{eq:pendulum} with \(c>0\), set
\(\omega_0^2=g/\ell\) and \(\gamma=c/(m\ell)\), and consider the energy
\begin{equation}\label{eq:pendulum-energy}
  V(\theta,\omega) = \frac{1}{2}\omega^2 + \omega_0^2(1-\cos\theta).
\end{equation}
Then \(V(0,0)=0\), \(V>0\) for \((\theta,\omega)\neq(0,0)\) with
\(\abs{\theta}<\pi\), and
\[
  \dot{V} = \omega\,\omega' + \omega_0^2\sin\theta\cdot\theta'
  = \omega(-\omega_0^2\sin\theta-\gamma\omega)+\omega_0^2\omega\sin\theta
  = -\gamma\,\omega^2 \le 0.
\]
By Theorem~\ref{thm:lyapunov-stability}, the origin is stable. Since
\(\dot{V}=0\) only when \(\omega=0\), we can in fact conclude asymptotic
stability via LaSalle's principle (Theorem~\ref{thm:lasalle}).
\end{example}

\subsection{LaSalle's Invariance Principle}
\label{subsec:lasalle}
\index{LaSalle's invariance principle|(}

When \(\dot{V}\le 0\) but not strictly negative everywhere, the asymptotic
stability theorem does not apply directly. LaSalle's principle fills this gap.

\begin{definition}[Invariant set]
\label{def:invariant-set}
\index{invariant set}
A set \(M\subseteq D\) is \textbf{positively invariant} under
\(\mathbf{x}'=\mathbf{f}(\mathbf{x})\) if every solution starting in~\(M\)
remains in~\(M\) for all \(t\ge 0\). It is \textbf{invariant} if solutions
remain in~\(M\) for all \(t\in\R\).
\end{definition}

\begin{theorem}[LaSalle's invariance principle]
\label{thm:lasalle}
\index{LaSalle's invariance principle}
Let \(\Omega\subseteq D\) be a compact positively invariant set for
\(\mathbf{x}'=\mathbf{f}(\mathbf{x})\). Let \(V\colon D\to\R\) be
continuously differentiable with \(\dot{V}(\mathbf{x})\le 0\) on~\(\Omega\).
Define
\[
  E = \{\mathbf{x}\in\Omega : \dot{V}(\mathbf{x})=0\},
\]
and let \(M\) be the largest invariant subset of~\(E\). Then every solution
starting in~\(\Omega\) converges to~\(M\) as \(t\to\infty\).
\end{theorem}

\begin{corollary}
\label{cor:lasalle-asymptotic}
If \(V\) is a Lyapunov function on a neighborhood of the origin,
\(\dot{V}\le 0\), and the only invariant subset of
\(\{x:\dot{V}(\mathbf{x})=0\}\) is \(\{\mathbf{0}\}\), then the origin is
asymptotically stable.
\end{corollary}

\begin{example}[Pendulum revisited]
\label{ex:pendulum-lasalle}
In Example~\ref{ex:pendulum-lyapunov}, \(\dot{V}=0\) iff \(\omega=0\). On the
set \(\{\omega=0\}\), the dynamics require \(\omega'=-\omega_0^2\sin\theta=0\),
hence \(\theta=0\). Thus \(M=\{(0,0)\}\), and LaSalle's principle gives
asymptotic stability of the hanging equilibrium.
\end{example}

\index{LaSalle's invariance principle|)}
\index{Lyapunov!direct method|)}

\section{Linearization and the Hartman--Grobman Theorem}
\label{sec:linearization}
\index{linearization|(}

\begin{definition}[Jacobian matrix]
\label{def:jacobian}
\index{Jacobian matrix}
The \textbf{Jacobian matrix}\index{Jacobian matrix} of
\(\mathbf{f}\colon\R^n\to\R^n\) at an equilibrium \(\mathbf{x}^*\) is the
\(n\times n\) matrix
\[
  J = D\mathbf{f}(\mathbf{x}^*) =
  \begin{pmatrix}
    \dfrac{\partial f_1}{\partial x_1} & \cdots & \dfrac{\partial f_1}{\partial x_n} \\[6pt]
    \vdots & \ddots & \vdots \\[3pt]
    \dfrac{\partial f_n}{\partial x_1} & \cdots & \dfrac{\partial f_n}{\partial x_n}
  \end{pmatrix}_{\mathbf{x}=\mathbf{x}^*}.
\]
The \textbf{linearization}\index{linearization!definition} of the system at
\(\mathbf{x}^*\) is \(\mathbf{y}'=J\,\mathbf{y}\), where
\(\mathbf{y}=\mathbf{x}-\mathbf{x}^*\).
\end{definition}

\begin{definition}[Hyperbolic equilibrium]
\label{def:hyperbolic}
\index{hyperbolic equilibrium}\index{equilibrium point!hyperbolic}
An equilibrium \(\mathbf{x}^*\) is \textbf{hyperbolic} if no eigenvalue of
\(J=D\mathbf{f}(\mathbf{x}^*)\) has zero real part.
\end{definition}

\begin{theorem}[Hartman--Grobman]
\label{thm:hartman-grobman}
\index{Hartman--Grobman theorem}
Let \(\mathbf{x}^*\) be a hyperbolic equilibrium of
\(\mathbf{x}'=\mathbf{f}(\mathbf{x})\) with Jacobian~\(J\). Then there exists
a homeomorphism of a neighborhood of~\(\mathbf{x}^*\) onto a neighborhood of
the origin that maps orbits of the nonlinear system to orbits of the linear
system \(\mathbf{y}'=J\,\mathbf{y}\), preserving the direction of time.
\end{theorem}

\begin{remark}
\label{rem:hartman-grobman}
The Hartman--Grobman theorem guarantees that near a hyperbolic equilibrium,
the nonlinear phase portrait is qualitatively the same as the linearized one.
In particular:
\begin{itemize}
  \item If all eigenvalues of~\(J\) have negative real parts, the equilibrium
    is asymptotically stable.
  \item If any eigenvalue has positive real part, the equilibrium is unstable.
\end{itemize}
For non-hyperbolic equilibria (eigenvalues on the imaginary axis), the
linearization is \emph{inconclusive}, and one must resort to Lyapunov methods
or center manifold theory.
\end{remark}

\begin{proposition}[Stability via linearization in \(\R^2\)]
\label{prop:linearization-2d}
\index{linearization!two-dimensional}
For a planar system \(\mathbf{x}'=\mathbf{f}(\mathbf{x})\) with
equilibrium at the origin and Jacobian~\(J\), let
\(\tau=\tr J\)\index{trace} and \(\Delta=\det J\)\index{determinant}. Then:
\begin{enumerate}[label=(\roman*)]
  \item If \(\Delta<0\): saddle (unstable).
  \item If \(\Delta>0\) and \(\tau<0\): asymptotically stable (node if
    \(\tau^2>4\Delta\), spiral if \(\tau^2<4\Delta\)).
  \item If \(\Delta>0\) and \(\tau>0\): unstable (node or spiral).
  \item If \(\Delta>0\) and \(\tau=0\): center for the linearization,
    inconclusive for the nonlinear system.
\end{enumerate}
\end{proposition}

\index{linearization|)}

\section{Phase Plane Analysis}
\label{sec:phase-plane}
\index{phase plane|(}

Phase plane analysis provides geometric insight into the behavior of
two-dimensional autonomous systems. We sketch trajectories in the
\((x_1,x_2)\)-plane, marking equilibria, determining their types, and
identifying features such as separatrices, limit cycles, and basins of
attraction.

\begin{method}[Phase plane analysis]
\label{meth:phase-plane}
\index{phase plane!method}
For \(\mathbf{x}'=\mathbf{f}(\mathbf{x})\) in \(\R^2\):
\begin{enumerate}[label=\textbf{Step~\arabic*.}]
  \item Find all equilibria by solving \(\mathbf{f}(\mathbf{x})=\mathbf{0}\).
  \item Compute the Jacobian at each equilibrium; determine eigenvalues.
  \item Classify each equilibrium (node, saddle, spiral, center).
  \item Sketch nullclines\index{nullcline}: curves where \(f_1=0\) or \(f_2=0\).
  \item Determine the direction field and sketch representative trajectories.
  \item Identify invariant sets, separatrices, limit cycles if present.
\end{enumerate}
\end{method}

\subsection{Phase Portraits}
\label{subsec:phase-portraits}

\begin{example}[Saddle point]
\label{ex:saddle-portrait}
\index{saddle point!phase portrait}
The system \(x'=x,\; y'=-y\) has eigenvalues \(\lambda_1=1,\;\lambda_2=-1\)
at the origin: a saddle. Solutions are \(x(t)=x_0e^t,\; y(t)=y_0e^{-t}\),
and trajectories satisfy \(xy=\text{const}\).

\begin{center}
\begin{tikzpicture}[>=Stealth,scale=0.9]
  \draw[->] (-2.8,0) -- (2.8,0) node[right]{\(x\)};
  \draw[->] (0,-2.3) -- (0,2.3) node[above]{\(y\)};
  % Stable manifold (y-axis)
  \draw[blue,thick,->] (0,2) -- (0,0.5);
  \draw[blue,thick,->] (0,-2) -- (0,-0.5);
  % Unstable manifold (x-axis)
  \draw[red,thick,->] (0.3,0) -- (2.2,0);
  \draw[red,thick,->] (-0.3,0) -- (-2.2,0);
  % Hyperbolas xy=c
  \foreach \c in {0.3,0.7,1.2}{
    \draw[gray,thick,domain=0.15:2.5,samples=40] plot (\x,{\c/\x});
    \draw[gray,thick,domain=-2.5:-0.15,samples=40] plot (\x,{\c/\x});
    \draw[gray,thick,domain=0.15:2.5,samples=40] plot (\x,{-\c/\x});
    \draw[gray,thick,domain=-2.5:-0.15,samples=40] plot (\x,{-\c/\x});
  }
  \fill (0,0) circle (2pt);
  \node[below left] at (0,0) {\(\mathbf{0}\)};
  \node[blue] at (0.6,1.9) {stable};
  \node[red] at (2.2,0.4) {unstable};
\end{tikzpicture}
\end{center}
\end{example}

\begin{example}[Stable spiral]
\label{ex:spiral-portrait}
\index{spiral!phase portrait}
The system \(x'=-x+2y,\; y'=-2x-y\) has Jacobian with
\(\tau=-2,\;\Delta=5\), and eigenvalues \(-1\pm 2i\): a stable spiral.

\begin{center}
\begin{tikzpicture}[>=Stealth,scale=0.9]
  \draw[->] (-2.5,0) -- (2.5,0) node[right]{\(x\)};
  \draw[->] (0,-2.5) -- (0,2.5) node[above]{\(y\)};
  % Spiral trajectory (approximate)
  \draw[blue,thick,domain=0:720,samples=50,->]
    plot ({2*exp(-\x/360)*cos(\x)},{2*exp(-\x/360)*sin(\x)});
  \draw[blue,thick,domain=0:720,samples=50,->]
    plot ({-2*exp(-\x/360)*cos(\x)},{-2*exp(-\x/360)*sin(\x)});
  \fill (0,0) circle (2pt);
  \node[below right] at (0.1,-0.2) {\(\mathbf{0}\)};
\end{tikzpicture}
\end{center}
\end{example}

\section{Examples: Nonlinear Systems}
\label{sec:nonlinear-examples}

\subsection{Nonlinear Pendulum}
\label{subsec:nonlinear-pendulum}
\index{pendulum!phase portrait}

The undamped pendulum \(\theta''+(g/\ell)\sin\theta=0\) has the phase portrait
below. The separatrices (heteroclinic orbits connecting saddle points at
\(\pm\pi\)) divide the phase plane into oscillatory regions (closed orbits) and
rotational regions (unbounded trajectories).

\begin{center}
\begin{tikzpicture}[>=Stealth]
  \begin{axis}[
    width=10cm, height=6cm,
    xlabel={\(\theta\)}, ylabel={\(\omega\)},
    xmin=-3.8, xmax=3.8, ymin=-3, ymax=3,
    axis lines=middle,
    xtick={-3.14159,0,3.14159},
    xticklabels={\(-\pi\),\(0\),\(\pi\)},
    ytick=\empty,
    samples=50,
    domain=-3.14:3.14,
  ]
    % Closed orbits (oscillation) using energy E = w^2/2 + (1-cos(theta)) = C
    \foreach \E in {0.3,0.7,1.2,1.6}{
      \addplot[blue,thick] ({x},{sqrt(max(2*(\E-(1-cos(deg(x)))),0))});
      \addplot[blue,thick] ({x},{-sqrt(max(2*(\E-(1-cos(deg(x)))),0))});
    }
    % Separatrix E=2
    \addplot[red,thick] ({x},{sqrt(max(2*(2-(1-cos(deg(x)))),0))});
    \addplot[red,thick] ({x},{-sqrt(max(2*(2-(1-cos(deg(x)))),0))});
    % Equilibria
    \addplot[only marks,mark=*,mark size=2.5pt] coordinates {(0,0)};
    \addplot[only marks,mark=o,mark size=2.5pt] coordinates
      {(-3.14159,0) (3.14159,0)};
    \node at (axis cs:0.5,-0.4) {center};
    \node at (axis cs:3.14,-0.5) {saddle};
  \end{axis}
\end{tikzpicture}
\end{center}

\subsection{Lotka--Volterra Predator-Prey Model}
\label{subsec:lotka-volterra}
\index{Lotka--Volterra model|(}

\begin{example}[Lotka--Volterra system]
\label{ex:lotka-volterra}
The classical predator-prey model is
\begin{equation}\label{eq:lotka-volterra}
  x' = \alpha x - \beta xy, \qquad y' = -\gamma y + \delta xy,
\end{equation}
where \(x\ge 0\) is the prey population, \(y\ge 0\) is the predator population,
and \(\alpha,\beta,\gamma,\delta>0\).

\medskip\noindent\textbf{Equilibria.}
Setting the right-hand side to zero gives two equilibria:
\[
  \mathbf{x}_1^* = (0,0), \qquad
  \mathbf{x}_2^* = \left(\frac{\gamma}{\delta},\frac{\alpha}{\beta}\right).
\]

\medskip\noindent\textbf{Linearization at \(\mathbf{x}_2^*\).}
The Jacobian is
\[
  J = \begin{pmatrix} \alpha-\beta y & -\beta x \\
                       \delta y & -\gamma+\delta x \end{pmatrix}.
\]
At \(\mathbf{x}_2^*=(\gamma/\delta,\alpha/\beta)\):
\[
  J = \begin{pmatrix} 0 & -\beta\gamma/\delta \\
      \alpha\delta/\beta & 0 \end{pmatrix},
  \qquad \lambda = \pm i\sqrt{\alpha\gamma}.
\]
The eigenvalues are purely imaginary, so the linearization predicts a center.
This is a non-hyperbolic case; the linearization is inconclusive.

\medskip\noindent\textbf{Lyapunov analysis.}
Define the conserved quantity (first integral)
\begin{equation}\label{eq:lv-conserved}
  V(x,y) = \delta x - \gamma\ln x + \beta y - \alpha\ln y.
\end{equation}
One verifies that \(\dot{V}=0\) along solutions, so trajectories lie on level
curves of~\(V\). Since \(V\) is strictly convex near \(\mathbf{x}_2^*\) with a
minimum there, the level curves are closed, confirming that \(\mathbf{x}_2^*\)
is a center (stable but not asymptotically stable).
\end{example}

\begin{center}
\begin{tikzpicture}[>=Stealth]
  \begin{axis}[
    width=8cm, height=7cm,
    xlabel={\(x\) (prey)}, ylabel={\(y\) (predator)},
    xmin=0, xmax=5, ymin=0, ymax=4,
    axis lines=left,
    samples=50,
  ]
    % Level curves of V = x - ln(x) + y - ln(y) (alpha=beta=gamma=delta=1)
    % Equilibrium at (1,1), V(1,1)=2
    % Approximate closed orbits around equilibrium (1,1)
    % Since contour plots are complex in basic pgfplots, draw parametric orbits
    % For illustration, draw approximate closed orbits
    \draw[blue,thick] (axis cs:1,1) circle[x radius=0.4cm, y radius=0.35cm];
    \draw[blue,thick,rotate around={-5:(axis cs:1,1)}]
      (axis cs:1,1) ellipse[x radius=0.85cm, y radius=0.7cm];
    \draw[blue,thick,rotate around={-10:(axis cs:1,1)}]
      (axis cs:1,1) ellipse[x radius=1.4cm, y radius=1.1cm];
    \draw[blue,thick,rotate around={-15:(axis cs:1,1)}]
      (axis cs:1,1) ellipse[x radius=2.0cm, y radius=1.5cm];
    \addplot[only marks,mark=*,mark size=2.5pt,red] coordinates {(1,1)};
    \node[red,above right] at (axis cs:1,1) {center};
  \end{axis}
\end{tikzpicture}
\end{center}

\index{Lotka--Volterra model|)}

\subsection{Van der Pol Oscillator}
\label{subsec:van-der-pol}
\index{van der Pol oscillator|(}

\begin{example}[Van der Pol oscillator]
\label{ex:van-der-pol}
The van der Pol equation\index{van der Pol oscillator!equation} is
\begin{equation}\label{eq:van-der-pol}
  x'' - \mu(1-x^2)x' + x = 0, \qquad \mu>0,
\end{equation}
or as a system:
\[
  x' = y, \qquad y' = \mu(1-x^2)y - x.
\]
The unique equilibrium is the origin.

\medskip\noindent\textbf{Linearization.}
The Jacobian at the origin is
\[
  J = \begin{pmatrix} 0 & 1 \\ -1 & \mu \end{pmatrix},
  \qquad \tau=\mu>0,\;\Delta=1>0.
\]
Since \(\tau>0\), the origin is an unstable spiral (for small \(\mu\)) or
unstable node (for \(\mu\ge 2\)).

\medskip\noindent\textbf{Limit cycle.}
Despite the instability of the equilibrium, all trajectories are bounded (one
can show this using Lyapunov-type arguments). By the Poincar\'e--Bendixson
theorem\index{Poincar\'e--Bendixson theorem}, there must exist a stable
limit cycle\index{limit cycle}. This is a distinctive feature of relaxation
oscillations: the system settles into a periodic orbit regardless of initial
conditions.

\begin{center}
\begin{tikzpicture}[>=Stealth]
  \begin{axis}[
    width=8cm, height=7cm,
    xlabel={\(x\)}, ylabel={\(y\)},
    xmin=-3.5, xmax=3.5, ymin=-5, ymax=5,
    axis lines=middle,
    samples=50,
  ]
    % Approximate limit cycle for mu=1 (roughly an oval)
    \addplot[red,very thick,domain=0:360,samples=50]
      ({2.0*cos(x)+0.15*cos(3*x)},{-2.0*sin(x)+1.2*sin(x)*cos(x)});
    % Spiral out from origin
    \draw[blue,thick,domain=0:500,samples=50,->]
      plot ({0.5*exp(\x/1200)*cos(\x)},{-0.5*exp(\x/1200)*sin(\x)});
    \addplot[only marks,mark=*,mark size=2pt] coordinates {(0,0)};
    \node[below right] at (axis cs:0.2,-0.3) {unstable};
    \node[red] at (axis cs:2.8,3.5) {limit cycle};
  \end{axis}
\end{tikzpicture}
\end{center}
\end{example}

\index{van der Pol oscillator|)}

\section{Lyapunov Functions: Construction Techniques}
\label{sec:lyapunov-construction}
\index{Lyapunov function!construction}

Finding Lyapunov functions is generally an art. We collect useful strategies.

\begin{method}[Energy-based Lyapunov functions]
\label{meth:energy-lyapunov}
\index{Lyapunov function!energy-based}
For mechanical systems \(q''+g(q)=0\) (possibly with damping), the total
energy
\[
  V(q,p) = \frac{1}{2}p^2 + \int_0^q g(s)\,\dd s
  \qquad (p=q')
\]
is a natural candidate.
\end{method}

\begin{method}[Quadratic Lyapunov functions]
\label{meth:quadratic-lyapunov}
\index{Lyapunov function!quadratic}
For a linear system \(\mathbf{x}'=A\mathbf{x}\) with \(A\) having all
eigenvalues with negative real parts, the quadratic form
\(V(\mathbf{x})=\mathbf{x}^T P\,\mathbf{x}\) is a Lyapunov function, where
\(P\) is the unique positive definite solution of the \textbf{Lyapunov
equation}\index{Lyapunov equation}
\begin{equation}\label{eq:lyapunov-equation}
  A^T P + P A = -Q
\end{equation}
for any chosen positive definite matrix~\(Q\) (commonly \(Q=I\)).
\end{method}

\begin{example}[Quadratic Lyapunov function]
\label{ex:quadratic-lyapunov}
Consider \(x'=-x+y,\;y'=-x-3y\). The matrix \(A=\bigl(\begin{smallmatrix}
-1&1\\-1&-3\end{smallmatrix}\bigr)\) has eigenvalues \(-2\pm i\) (both with
negative real part). Taking \(Q=I\), the Lyapunov equation
\(A^TP+PA=-I\) yields
\[
  P = \begin{pmatrix} 5/8 & -1/8 \\ -1/8 & 3/8 \end{pmatrix},
\]
and \(V(\mathbf{x})=\mathbf{x}^TP\mathbf{x}=\frac{1}{8}(5x^2-2xy+3y^2)\) is a
Lyapunov function with \(\dot{V}=-x^2-y^2\).
\end{example}

\begin{method}[Variable gradient method]
\label{meth:variable-gradient}
\index{Lyapunov function!variable gradient}
Seek \(V\) by postulating a gradient form \(\nabla V = M(\mathbf{x})\) with
undetermined coefficients, then imposing the integrability condition
\(\partial M_i/\partial x_j=\partial M_j/\partial x_i\) and the sign
condition \(\dot{V}=M\cdot\mathbf{f}\le 0\).
\end{method}

\section{Lyapunov Level Curves}
\label{sec:lyapunov-level-curves}

\begin{example}[Visualizing a Lyapunov function]
\label{ex:lyapunov-level-curves}
Consider the system \(x'=-x+y^2,\;y'=-y\). The origin is an equilibrium.
The function \(V(x,y)=x^2+y^2\) gives
\[
  \dot{V} = 2x(-x+y^2)+2y(-y) = -2x^2+2xy^2-2y^2.
\]
This is not sign-definite globally, but in a small neighborhood of the origin
(say \(\abs{y}<1\)), we have \(2xy^2\le 2\abs{x}y^2 \le x^2+y^4\le x^2+y^2\)
for \(\abs{y}\le 1\), so \(\dot{V}\le -x^2-y^2<0\) for small
\(\norm{(x,y)}\), confirming local asymptotic stability.

\begin{center}
\begin{tikzpicture}[>=Stealth]
  \begin{axis}[
    width=7cm, height=7cm,
    xlabel={\(x\)}, ylabel={\(y\)},
    xmin=-1.5, xmax=1.5, ymin=-1.5, ymax=1.5,
    axis lines=middle,
    view={0}{90},
  ]
    % Level curves of V = x^2 + y^2
    \foreach \r in {0.3,0.6,0.9,1.2}{
      \addplot[green!60!black,thick,domain=0:360,samples=50]
        ({\r*cos(x)},{\r*sin(x)});
    }
    % Direction field arrows (hand-placed for clarity)
    \foreach \xx in {-1.0,-0.5,0.5,1.0}{
      \foreach \yy in {-1.0,-0.5,0.5,1.0}{
        \edef\temp{\noexpand\draw[->,gray,thick] (axis cs:\xx,\yy) --
          (axis cs:{\xx+0.15*(-\xx+\yy*\yy)},{\yy+0.15*(-\yy)});}
        \temp
      }
    }
    \addplot[only marks,mark=*,mark size=2.5pt,red] coordinates {(0,0)};
    \node[green!50!black] at (axis cs:1.1,1.3) {\(V=c\)};
  \end{axis}
\end{tikzpicture}
\end{center}
\end{example}

\section{Summary of Stability Criteria}
\label{sec:stability-summary-table}

\begin{center}
\renewcommand{\arraystretch}{1.35}
\begin{tabular}{@{}p{3.8cm}p{5.5cm}p{4.5cm}@{}}
\toprule
\textbf{Method} & \textbf{Condition} & \textbf{Conclusion} \\
\midrule
Eigenvalues of~\(J\) &
All \(\operatorname{Re}(\lambda_i)<0\) &
Asymptotically stable \\
& Any \(\operatorname{Re}(\lambda_i)>0\) &
Unstable \\
& Some \(\operatorname{Re}(\lambda_i)=0\), rest \(<0\) &
Inconclusive (non-hyperbolic) \\
\midrule
Lyapunov function &
\(V>0\), \(\dot{V}\le 0\) &
Stable \\
& \(V>0\), \(\dot{V}<0\) &
Asymptotically stable \\
& LaSalle: \(\dot{V}\le 0\), largest invariant set in \(\{\dot{V}=0\}\) is \(\{\mathbf{0}\}\) &
Asymptotically stable \\
\midrule
Trace--determinant &
\(\Delta>0\), \(\tau<0\) &
Asymptotically stable (2D) \\
(\(\R^2\) only) &
\(\Delta<0\) &
Saddle (unstable) \\
\bottomrule
\end{tabular}
\end{center}

\section{Exercises}
\label{sec:stability-exercises}

\begin{exercise}
\label{exo:classify-equilibria}
Find and classify all equilibria of the system
\[
  x' = x(3-x-2y), \qquad y' = y(2-x-y).
\]
Determine the stability of each equilibrium using linearization.
\end{exercise}

\begin{exercise}
\label{exo:lyapunov-construct}
\index{Lyapunov function!construction exercise}
For the system \(x'=-x^3,\;y'=-y\), show that \(V(x,y)=x^4/4+y^2/2\) is a
strict Lyapunov function, and conclude that the origin is asymptotically stable.
\end{exercise}

\begin{exercise}
\label{exo:lasalle-exercise}
Consider \(x''+x'+x^3=0\). Write this as a first-order system, find a Lyapunov
function, and use LaSalle's invariance principle to prove asymptotic stability
of the origin.
\end{exercise}

\begin{exercise}
\label{exo:lotka-volterra-exercise}
For the Lotka--Volterra system \eqref{eq:lotka-volterra} with
\(\alpha=1,\beta=0.5,\gamma=0.75,\delta=0.25\):
\begin{enumerate}[label=(\alph*)]
  \item Find the non-trivial equilibrium.
  \item Verify that \(V(x,y)=0.25x-0.75\ln x+0.5y-\ln y\) is a conserved
    quantity.
  \item Sketch several level curves of~\(V\) in the first quadrant.
\end{enumerate}
\end{exercise}

\begin{exercise}
\label{exo:linearize-compete}
Two species compete according to
\[
  x' = x(1-x-\alpha y), \qquad y' = y(1-\beta x - y).
\]
Find all equilibria. For which values of \(\alpha,\beta>0\) is the coexistence
equilibrium stable?
\end{exercise}

\begin{exercise}
\label{exo:van-der-pol-exercise}
\index{van der Pol oscillator!exercise}
For the van der Pol oscillator \eqref{eq:van-der-pol} with \(\mu=0.5\):
\begin{enumerate}[label=(\alph*)]
  \item Show that the origin is an unstable spiral.
  \item Show that the set \(\{(x,y): x^2+y^2\le R^2\}\) is positively
    invariant for sufficiently large~\(R\).
  \item Conclude the existence of a limit cycle by the Poincar\'e--Bendixson
    theorem.
\end{enumerate}
\end{exercise}

\begin{exercise}
\label{exo:lyapunov-equation}
For the matrix \(A=\begin{pmatrix}-2&1\\0&-3\end{pmatrix}\), solve the
Lyapunov equation \(A^TP+PA=-I\) and verify that
\(V(\mathbf{x})=\mathbf{x}^TP\mathbf{x}\) is a Lyapunov function for
\(\mathbf{x}'=A\mathbf{x}\).
\end{exercise}

\begin{exercise}
\label{exo:global-stability}
Consider the system
\[
  x' = -x + 2\arctan(y), \qquad y' = -y + 2\arctan(x).
\]
\begin{enumerate}[label=(\alph*)]
  \item Show that the origin is the unique equilibrium.
  \item Show that the origin is locally asymptotically stable by linearization.
  \item Attempt to construct a Lyapunov function for global stability analysis.
  \item Does the system have other attractors? Justify your answer.
\end{enumerate}
\end{exercise}

\section{Chapter Summary}
\label{sec:stability-chapter-summary}

\begin{itemize}
  \item An equilibrium \(\mathbf{x}^*\) is Lyapunov stable if nearby solutions
    remain nearby, and asymptotically stable if they also converge to
    \(\mathbf{x}^*\).
  \item \textbf{Linearization}: near a hyperbolic equilibrium, the nonlinear
    system behaves qualitatively like its linearization (Hartman--Grobman
    theorem). Non-hyperbolic cases require further analysis.
  \item \textbf{Lyapunov's direct method}: a positive definite function with
    non-positive orbital derivative proves stability; strict negativity gives
    asymptotic stability.
  \item \textbf{LaSalle's invariance principle} extends asymptotic stability
    conclusions to cases where \(\dot{V}\le 0\) but is not strictly negative
    everywhere.
  \item The trace--determinant plane gives a complete classification of linear
    planar systems.
  \item Phase plane analysis---nullclines, direction fields, separatrices---provides
    geometric understanding of two-dimensional dynamics.
  \item Classical examples (pendulum, Lotka--Volterra, van der Pol) illustrate
    the interplay between linearization, Lyapunov methods, and geometric
    reasoning.
\end{itemize}

\index{Lyapunov stability|)}
\index{stability|)}

% === END OF CHUNK ch7-8 ===
% === START OF CHUNK ch9-10+end ===

%%%%%%%%%%%%%%%%%%%%%%%%%%%%%%%%%%%%%%%%%%%%%%%%%%%%%%%%%%%%%%%%%%%%%
%%  CHAPTER 9 — Introduction to Nonlinear Systems and Bifurcations %%
%%%%%%%%%%%%%%%%%%%%%%%%%%%%%%%%%%%%%%%%%%%%%%%%%%%%%%%%%%%%%%%%%%%%%
\chapter{Introduction to Nonlinear Systems and Bifurcations}
\label{ch:nonlinear-bifurcations}
\index{nonlinear systems}

\section{Motivation: Nonlinearity in Nature}
\label{sec:nonlinear-motivation}
\index{nonlinear systems!motivation}

Many phenomena in science and engineering are inherently nonlinear.
Population dynamics, chemical reactions, fluid mechanics, and climate
models all give rise to differential equations that cannot be reduced
to linear form. In this chapter we develop qualitative tools for
analysing autonomous planar systems
\begin{equation}\label{eq:planar-system}
  \dot{x} = f(x,y), \qquad \dot{y} = g(x,y),
\end{equation}
where $f,g\colon U\to\R$ are smooth functions defined on an open set
$U\subseteq\R^{2}$.

\begin{example}[Population growth with saturation]
\label{ex:logistic-motivation}
\index{logistic equation}
The logistic equation
\[
  \dot{N} = rN\!\left(1-\frac{N}{K}\right)
\]
models a population $N(t)$ with growth rate $r>0$ and carrying
capacity $K>0$. The nonlinear term $-rN^{2}/K$ prevents unbounded
exponential growth and produces a sigmoidal solution curve.
\end{example}

\begin{example}[Chemical kinetics]
\label{ex:chemical-motivation}
\index{chemical kinetics}
Consider two reacting species with concentrations $x$ and $y$ governed
by
\[
  \dot{x} = a - x - \frac{4xy}{1+x^{2}},
  \qquad
  \dot{y} = bx\!\left(1-\frac{y}{1+x^{2}}\right),
\]
where $a,b>0$ are parameters. The rational nonlinearities preclude
any closed-form solution; qualitative methods are essential.
\end{example}

%----------------------------------------------------------------------
\section{Phase Plane Analysis}
\label{sec:phase-plane}
\index{phase plane}

\subsection{Nullclines and Direction Fields}
\label{subsec:nullclines}
\index{nullclines}
\index{direction field}

\begin{definition}[Nullclines]
\label{def:nullclines}
For the system~\eqref{eq:planar-system}, the \emph{$x$-nullcline} is
the set $\{(x,y)\in U : f(x,y)=0\}$, and the \emph{$y$-nullcline}
is $\{(x,y)\in U : g(x,y)=0\}$. Equilibrium points lie at the
intersections of the two nullclines.
\end{definition}

On the $x$-nullcline the vector field is vertical ($\dot{x}=0$),
and on the $y$-nullcline it is horizontal ($\dot{y}=0$). The
nullclines therefore divide the phase plane into regions in which the
signs of $\dot{x}$ and $\dot{y}$ are constant.

\begin{example}[Nullclines of a competing-species model]
\label{ex:nullclines-competition}
\index{competition model}
Consider
\[
  \dot{x} = x(3-x-2y),\qquad \dot{y} = y(2-y-x).
\]
The $x$-nullcline consists of the lines $x=0$ and $x+2y=3$; the
$y$-nullcline consists of $y=0$ and $x+y=2$. The equilibrium points
are $(0,0)$, $(3,0)$, $(0,2)$, and $(1,1)$.
\end{example}

\begin{center}
\begin{tikzpicture}
  \begin{axis}[
    axis lines=middle,
    xlabel={$x$}, ylabel={$y$},
    xmin=-0.3, xmax=3.8, ymin=-0.3, ymax=2.8,
    width=9cm, height=7cm,
    samples=50,
    legend pos=north east
  ]
    % x-nullcline: x+2y=3
    \addplot[blue, thick, domain=0:3] {(3-x)/2};
    \addlegendentry{$x+2y=3$}
    % y-nullcline: x+y=2
    \addplot[red, thick, domain=0:2] {2-x};
    \addlegendentry{$x+y=2$}
    % equilibrium points
    \addplot[only marks, mark=*, mark size=2.5pt] coordinates {(0,0) (3,0) (0,2) (1,1)};
  \end{axis}
\end{tikzpicture}
\end{center}

\subsection{Equilibrium Classification in Nonlinear Systems}
\label{subsec:equilib-nonlinear}
\index{equilibrium!nonlinear classification}
\index{Hartman--Grobman theorem}

\begin{theorem}[Hartman--Grobman]
\label{thm:hartman-grobman}
Let $\mathbf{x}^{*}$ be a hyperbolic equilibrium of
$\dot{\mathbf{x}}=\mathbf{F}(\mathbf{x})$, i.e.\ the Jacobian
$D\mathbf{F}(\mathbf{x}^{*})$ has no eigenvalue with zero real part.
Then, near $\mathbf{x}^{*}$, the nonlinear flow is topologically
conjugate to the linearised flow
$\dot{\mathbf{u}}=D\mathbf{F}(\mathbf{x}^{*})\,\mathbf{u}$.
\end{theorem}

In practice, one computes the Jacobian
\[
  J = \begin{pmatrix}
    \partial f/\partial x & \partial f/\partial y \\
    \partial g/\partial x & \partial g/\partial y
  \end{pmatrix}\bigg|_{(x^{*},y^{*})}
\]
and classifies the equilibrium via the eigenvalues of $J$, exactly as
for linear systems, provided $\mathbf{x}^{*}$ is hyperbolic.

%----------------------------------------------------------------------
\section{The Poincar\'e--Bendixson Theorem}
\label{sec:poincare-bendixson}
\index{Poincar\'e--Bendixson theorem}

\begin{theorem}[Poincar\'e--Bendixson]
\label{thm:poincare-bendixson}
Let $\mathbf{F}\in C^{1}(\R^{2})$. Suppose that a forward orbit
$\gamma^{+}$ is contained in a compact set $K\subset\R^{2}$ that
contains no equilibrium points. Then $\gamma^{+}$ approaches a
periodic orbit (a \emph{limit cycle}\index{limit cycle}).
\end{theorem}

\begin{remark}
\label{rem:pb-dimension}
The Poincar\'e--Bendixson theorem is specific to $\R^{2}$.
In $\R^{3}$ and higher dimensions, bounded trajectories can exhibit
chaotic behaviour (cf.\ the Lorenz system).
\end{remark}

\begin{corollary}[Trapping region criterion]
\label{cor:trapping-region}
\index{trapping region}
If one can exhibit an annular region $A$ such that the vector field
points inward on both boundary components and $A$ contains no
equilibrium, then $A$ contains at least one limit cycle.
\end{corollary}

%----------------------------------------------------------------------
\section{Limit Cycles}
\label{sec:limit-cycles}
\index{limit cycle}

\begin{definition}[Limit cycle]
\label{def:limit-cycle}
A \emph{limit cycle} is an isolated periodic orbit $\Gamma$ of the
planar system~\eqref{eq:planar-system}. It is \emph{stable}
(attracting) if nearby orbits spiral towards $\Gamma$ as
$t\to+\infty$, \emph{unstable} if they spiral away, and
\emph{semi-stable} if the behaviour differs on each side.
\end{definition}

\subsection{The van der Pol Oscillator}
\label{subsec:van-der-pol}
\index{van der Pol oscillator}

The van der Pol equation
\begin{equation}\label{eq:van-der-pol}
  \ddot{x} + \mu(x^{2}-1)\dot{x} + x = 0, \qquad \mu>0,
\end{equation}
can be rewritten as the system
\[
  \dot{x} = y, \qquad \dot{y} = -x - \mu(x^{2}-1)y.
\]
The origin is the unique equilibrium. For any $\mu>0$, the Jacobian
there has eigenvalues with positive real part, so it is an unstable
spiral. By constructing a suitable trapping region (a large circle on
which the flow points inward), the Poincar\'e--Bendixson theorem
guarantees the existence of a stable limit cycle.

\begin{center}
\begin{tikzpicture}
  \begin{axis}[
    axis lines=middle,
    xlabel={$x$}, ylabel={$y$},
    xmin=-3.5, xmax=3.5, ymin=-5, ymax=5,
    width=9cm, height=7cm,
    title={Van der Pol limit cycle ($\mu=1$)},
    ticks=none
  ]
    % Approximate limit cycle for mu=1
    \addplot[blue, thick, domain=0:360, samples=50]
      ({2.0*cos(x)+0.15*cos(3*x)},{-2.0*sin(x)-0.9*sin(x)*abs(cos(x))});
    \addplot[only marks, mark=*, mark size=2pt] coordinates {(0,0)};
    \node at (axis cs:0.3,-0.4) {$O$};
  \end{axis}
\end{tikzpicture}
\end{center}

\subsection{Dulac's Criterion}
\label{subsec:dulac}
\index{Dulac's criterion}

Dulac's criterion provides a sufficient condition for the
\emph{non-existence} of periodic orbits.

\begin{theorem}[Dulac's criterion]
\label{thm:dulac}
Let $D\subseteq\R^{2}$ be a simply connected open set and
$B\colon D\to\R$ a $C^{1}$ function such that
$\nabla\cdot(B\mathbf{F})$ does not change sign in $D$ and is not
identically zero on any open subset. Then the
system~\eqref{eq:planar-system} has no periodic orbit lying entirely
in $D$.
\end{theorem}

\begin{proof}
Suppose $\Gamma$ is a periodic orbit enclosing a region $R\subset D$.
By Green's theorem,
\[
  \oint_{\Gamma} \bigl(-Bg\,\dd x + Bf\,\dd y\bigr)
  = \iint_{R} \left(\frac{\partial(Bf)}{\partial x}
    + \frac{\partial(Bg)}{\partial y}\right)\dd x\,\dd y.
\]
The left side equals
$\oint_{\Gamma}B(f\,\dd y - g\,\dd x)$; along $\Gamma$ we have
$\dd y/\dd t = g$ and $\dd x/\dd t = f$, so the integrand becomes
$B(f g - g f)=0$, giving $0$ on the left. But the right side has
a fixed sign, a contradiction.
\end{proof}

\begin{example}[Applying Dulac's criterion]
\label{ex:dulac-application}
For the system $\dot{x}=y$, $\dot{y}=-x-y^3$, choose $B=1$. Then
\[
  \frac{\partial f}{\partial x}+\frac{\partial g}{\partial y}
  = 0 + (-3y^{2}) = -3y^{2} \le 0,
\]
and $-3y^{2}=0$ only on the $x$-axis (a set of measure zero).
By Dulac's criterion there is no periodic orbit in $\R^{2}$.
\end{example}

%----------------------------------------------------------------------
\section{Bifurcation Theory}
\label{sec:bifurcations}
\index{bifurcation}

A \emph{bifurcation}\index{bifurcation!definition} occurs when a
small change in a parameter $\mu$ causes a qualitative change in the
phase portrait of $\dot{\mathbf{x}}=\mathbf{F}(\mathbf{x};\mu)$.

\subsection{Saddle-Node Bifurcation}
\label{subsec:saddle-node}
\index{bifurcation!saddle-node}

\begin{definition}[Saddle-node bifurcation]
\label{def:saddle-node}
The prototypical saddle-node (or fold) bifurcation is given by
\[
  \dot{x} = \mu + x^{2}.
\]
For $\mu<0$ there are two equilibria $x^{*}=\pm\sqrt{-\mu}$ (one
stable, one unstable); for $\mu>0$ there are none.
\end{definition}

\begin{center}
\begin{tikzpicture}
  \begin{axis}[
    axis lines=middle,
    xlabel={$\mu$}, ylabel={$x^*$},
    xmin=-2.5, xmax=1.5, ymin=-2, ymax=2,
    width=8cm, height=6cm,
    title={Saddle-node bifurcation},
    samples=50
  ]
    % Stable branch (lower)
    \addplot[blue, thick, domain=-2.2:0] {-sqrt(-x)};
    % Unstable branch (upper)
    \addplot[blue, thick, dashed, domain=-2.2:0] {sqrt(-x)};
    \addplot[only marks, mark=*, mark size=2.5pt] coordinates {(0,0)};
  \end{axis}
\end{tikzpicture}
\end{center}

\subsection{Transcritical Bifurcation}
\label{subsec:transcritical}
\index{bifurcation!transcritical}

\begin{definition}[Transcritical bifurcation]
\label{def:transcritical}
The normal form is
\[
  \dot{x} = \mu x - x^{2}.
\]
There are always two equilibria, $x^{*}=0$ and $x^{*}=\mu$, which
exchange stability as $\mu$ passes through~$0$.
\end{definition}

\begin{center}
\begin{tikzpicture}
  \begin{axis}[
    axis lines=middle,
    xlabel={$\mu$}, ylabel={$x^*$},
    xmin=-2, xmax=2, ymin=-2, ymax=2,
    width=8cm, height=6cm,
    title={Transcritical bifurcation},
    samples=50
  ]
    % x*=0: stable for mu<0, unstable for mu>0
    \addplot[blue, thick, domain=-2:0] {0};
    \addplot[blue, thick, dashed, domain=0:2] {0};
    % x*=mu: unstable for mu<0, stable for mu>0
    \addplot[red, thick, dashed, domain=-2:0] {x};
    \addplot[red, thick, domain=0:2] {x};
    \addplot[only marks, mark=*, mark size=2.5pt] coordinates {(0,0)};
  \end{axis}
\end{tikzpicture}
\end{center}

\subsection{Pitchfork Bifurcation}
\label{subsec:pitchfork}
\index{bifurcation!pitchfork}

\begin{definition}[Pitchfork bifurcation]
\label{def:pitchfork}
The \emph{supercritical} pitchfork has normal form
\[
  \dot{x} = \mu x - x^{3}.
\]
For $\mu\le 0$, $x^{*}=0$ is the unique (stable) equilibrium. For
$\mu>0$, the origin becomes unstable and two stable branches
$x^{*}=\pm\sqrt{\mu}$ appear.
\end{definition}

\begin{center}
\begin{tikzpicture}
  \begin{axis}[
    axis lines=middle,
    xlabel={$\mu$}, ylabel={$x^*$},
    xmin=-1.5, xmax=2.5, ymin=-2, ymax=2,
    width=8cm, height=6cm,
    title={Supercritical pitchfork},
    samples=50
  ]
    % x*=0: stable for mu<0, unstable for mu>0
    \addplot[blue, thick, domain=-1.5:0] {0};
    \addplot[blue, thick, dashed, domain=0:2.5] {0};
    % x*=+sqrt(mu)
    \addplot[blue, thick, domain=0.01:2.5] {sqrt(x)};
    % x*=-sqrt(mu)
    \addplot[blue, thick, domain=0.01:2.5] {-sqrt(x)};
    \addplot[only marks, mark=*, mark size=2.5pt] coordinates {(0,0)};
  \end{axis}
\end{tikzpicture}
\end{center}

\subsection{Hopf Bifurcation}
\label{subsec:hopf}
\index{bifurcation!Hopf}

\begin{theorem}[Hopf bifurcation]
\label{thm:hopf}
Consider a planar system
$\dot{\mathbf{x}}=\mathbf{F}(\mathbf{x};\mu)$ with an equilibrium
$\mathbf{x}^{*}(\mu)$ whose Jacobian has eigenvalues
$\lambda(\mu)=\alpha(\mu)\pm i\beta(\mu)$. Suppose that at
$\mu=\mu_{0}$:
\begin{enumerate}[label=(\roman*)]
  \item $\alpha(\mu_{0})=0$ and $\beta(\mu_{0})\neq 0$
    (a pair of purely imaginary eigenvalues);
  \item $\alpha'(\mu_{0})\neq 0$
    (the eigenvalues cross the imaginary axis with non-zero speed).
\end{enumerate}
Then a family of periodic orbits is born from the equilibrium at
$\mu=\mu_{0}$. In the \emph{supercritical}\index{Hopf bifurcation!supercritical}
case a stable limit cycle appears for
$\mu>\mu_{0}$; in the \emph{subcritical}\index{Hopf bifurcation!subcritical}
case an unstable limit cycle exists for $\mu<\mu_{0}$.
\end{theorem}

\begin{example}[Supercritical Hopf bifurcation]
\label{ex:hopf-super}
The system in polar coordinates
\[
  \dot{r} = \mu r - r^{3}, \qquad \dot{\theta} = 1,
\]
has equilibrium $r=0$ (the origin). For $\mu>0$, the stable limit
cycle $r=\sqrt{\mu}$ appears, illustrating a supercritical Hopf
bifurcation at $\mu=0$.
\end{example}

%----------------------------------------------------------------------
\section{Index Theory}
\label{sec:index-theory}
\index{index theory}

\begin{definition}[Index of a curve]
\label{def:index-curve}
Let $\mathcal{C}$ be a simple closed curve in the phase plane that
does not pass through any equilibrium. The \emph{index} of
$\mathcal{C}$ (with respect to the vector field $\mathbf{F}$) is the
number of complete counter-clockwise rotations made by the vector
$\mathbf{F}$ as one traverses $\mathcal{C}$ once counter-clockwise:
\[
  I_{\mathcal{C}}
  = \frac{1}{2\pi}\oint_{\mathcal{C}}\dd\!\left(\arctan\frac{g}{f}\right).
\]
\end{definition}

\begin{proposition}[Properties of the index]
\label{prop:index-properties}
\begin{enumerate}[label=(\roman*)]
  \item The index of a node, spiral, or centre is $+1$.
  \item The index of a saddle is $-1$.
  \item If $\mathcal{C}$ encloses no equilibrium, then
    $I_{\mathcal{C}}=0$.
  \item If $\mathcal{C}$ encloses equilibria $\mathbf{x}_{1}^{*},
    \dots,\mathbf{x}_{k}^{*}$, then
    $I_{\mathcal{C}}=\sum_{j=1}^{k}I_{j}$.
  \item A limit cycle must have index $+1$, so it must enclose
    equilibria whose indices sum to $+1$.
\end{enumerate}
\end{proposition}

\begin{corollary}
\label{cor:limit-cycle-equilibrium}
Every limit cycle must enclose at least one equilibrium point.
\end{corollary}

%----------------------------------------------------------------------
\section{Worked Examples}
\label{sec:nonlinear-examples}

\subsection{The Lotka--Volterra Predator--Prey Model}
\label{subsec:lotka-volterra-analysis}
\index{Lotka--Volterra equations}

\begin{example}[Lotka--Volterra system]
\label{ex:lotka-volterra-full}
The classical predator--prey model reads
\begin{equation}\label{eq:lotka-volterra}
  \dot{x} = \alpha x - \beta xy, \qquad
  \dot{y} = \delta xy - \gamma y,
\end{equation}
with $\alpha,\beta,\gamma,\delta>0$. Here $x$ denotes prey and $y$
denotes predators.

\emph{Equilibria.} Setting the right-hand sides to zero gives
$(0,0)$ and $(x^{*},y^{*})=(\gamma/\delta,\,\alpha/\beta)$.

\emph{Jacobian.} We compute
\[
  J = \begin{pmatrix}
    \alpha-\beta y & -\beta x \\
    \delta y & \delta x - \gamma
  \end{pmatrix}.
\]
At $(0,0)$: eigenvalues $\alpha>0$ and $-\gamma<0$ --- a saddle.

At $(\gamma/\delta,\alpha/\beta)$: the trace is $0$ and the
determinant is $\alpha\gamma>0$, so the eigenvalues are purely
imaginary. The equilibrium is a centre for the linearised system, and
the Lotka--Volterra system possesses the conserved quantity
\[
  H(x,y) = \delta x - \gamma\ln x + \beta y - \alpha\ln y,
\]
confirming that every orbit in the positive quadrant is a closed curve
(a true nonlinear centre).
\end{example}

\begin{center}
\begin{tikzpicture}
  \begin{axis}[
    axis lines=middle,
    xlabel={Prey $x$}, ylabel={Predator $y$},
    xmin=0, xmax=5, ymin=0, ymax=4,
    width=9cm, height=7cm,
    title={Lotka--Volterra phase portrait},
    ticks=none
  ]
    % Level curves of H = delta*x - gamma*ln(x) + beta*y - alpha*ln(y)
    % Using alpha=beta=delta=gamma=1 => H = x - ln x + y - ln y
    % Approximate concentric closed curves around (1,1)
    \addplot[blue, thick, domain=0:360, samples=50]
      ({1+0.4*cos(x)},{1+0.4*sin(x)});
    \addplot[blue, domain=0:360, samples=50]
      ({1+0.8*cos(x)-0.06*cos(2*x)},{1+0.8*sin(x)-0.06*sin(2*x)});
    \addplot[blue, domain=0:360, samples=50]
      ({1+1.3*cos(x)-0.15*cos(2*x)},{1+1.2*sin(x)-0.15*sin(2*x)});
    \addplot[only marks, mark=*, mark size=2.5pt] coordinates {(1,1)};
    \node at (axis cs:1.3,0.6) {$(x^{*},y^{*})$};
  \end{axis}
\end{tikzpicture}
\end{center}

\subsection{The Van der Pol Oscillator Revisited}
\label{subsec:vdp-worked}

\begin{example}[Li\'enard form and relaxation oscillations]
\label{ex:vdp-lienard}
\index{Li\'enard equation}
Setting $y = \dot{x}+\mu(x^{3}/3-x)$ transforms the van der Pol
equation~\eqref{eq:van-der-pol} into the Li\'enard form
\[
  \dot{x} = y - \mu\!\left(\frac{x^{3}}{3}-x\right), \qquad
  \dot{y} = -x.
\]
For large $\mu$, the limit cycle approaches the cubic nullcline and
the solution exhibits \emph{relaxation oscillations}\index{relaxation oscillations}:
slow drift along the nullcline branches alternating with fast jumps
between them.
\end{example}

\subsection{The SIR Epidemic Model}
\label{subsec:sir-analysis}
\index{SIR model}

\begin{example}[SIR dynamics]
\label{ex:sir-nonlinear}
The Kermack--McKendrick SIR model is
\[
  \dot{S} = -\beta SI,\qquad
  \dot{I} = \beta SI - \gamma I,\qquad
  \dot{R} = \gamma I,
\]
with $S+I+R=N$ constant. Since $R=N-S-I$, we reduce to the planar
system
\[
  \dot{S} = -\beta SI,\qquad
  \dot{I} = (\beta S - \gamma) I.
\]
The \emph{basic reproduction number}\index{basic reproduction number}
$\mathcal{R}_{0}=\beta S(0)/\gamma$ determines whether an epidemic
occurs: if $\mathcal{R}_{0}>1$, then $I$ initially increases.

\emph{Phase portrait.} All trajectories in the $(S,I)$-plane with
$S,I>0$ satisfy
\[
  \frac{\dd I}{\dd S} = -1 + \frac{\gamma}{\beta S},
\]
which integrates to $I = N - S + (\gamma/\beta)\ln(S/S_{0})$. The
disease-free line $I=0$ is an invariant set, and every orbit
approaches a point $(S_{\infty},0)$ with $S_{\infty}>0$.
\end{example}

%----------------------------------------------------------------------
\section{Exercises}
\label{sec:ch9-exercises}

\begin{exercise}
\label{exer:nullclines-practice}
For $\dot{x}=x(2-x-y)$, $\dot{y}=y(3-2y-x)$, find all equilibria
and sketch the nullclines and direction field in the first quadrant.
Classify each equilibrium using linearisation.
\end{exercise}

\begin{exercise}
\label{exer:dulac-exercise}
Show that the system $\dot{x}=2xy+x$, $\dot{y}=x^{2}-y^{2}$ has no
periodic orbit in the half-plane $\{y>0\}$. \textit{Hint:} try
$B(x,y)=1/y^{2}$.
\end{exercise}

\begin{exercise}
\label{exer:pb-exercise}
Consider $\dot{r}=r(1-r^{2})+\mu r$, $\dot{\theta}=1$ in polar
coordinates. For which values of $\mu$ does a limit cycle exist?
What is its radius?
\end{exercise}

\begin{exercise}
\label{exer:hopf-exercise}
For the system
\[
  \dot{x} = \mu x - y - x(x^{2}+y^{2}),\quad
  \dot{y} = x + \mu y - y(x^{2}+y^{2}),
\]
show that a Hopf bifurcation occurs at $\mu=0$. Is it supercritical
or subcritical? Find the radius of the limit cycle for $\mu>0$.
\end{exercise}

\begin{exercise}
\label{exer:bifurcation-diagram}
Sketch the bifurcation diagram for $\dot{x}=\mu-x^{2}+x^{3}$.
Identify the type of each bifurcation.
\end{exercise}

\begin{exercise}
\label{exer:lotka-volterra-conserved}
Verify that $H(x,y)=\delta x-\gamma\ln x+\beta y-\alpha\ln y$ is a
conserved quantity for the Lotka--Volterra
system~\eqref{eq:lotka-volterra}. Use this to prove that non-trivial
orbits in the positive quadrant are closed curves.
\end{exercise}

\begin{exercise}
\label{exer:index-exercise}
A planar system has exactly three equilibria: a stable node, an
unstable spiral, and a saddle. Can it have a limit cycle? If so,
which equilibria must the cycle enclose? Justify using index theory.
\end{exercise}

\begin{exercise}
\label{exer:sir-R0}
In the SIR model of Example~\ref{ex:sir-nonlinear}, prove that the
maximum of $I(t)$ is attained when $S=\gamma/\beta$. Express the
peak infection level $I_{\max}$ in terms of $S(0)$, $I(0)$,
$\gamma$, and $\beta$.
\end{exercise}

%----------------------------------------------------------------------
\section*{Chapter Summary}
\addcontentsline{toc}{section}{Chapter Summary}

\begin{itemize}
  \item Nullclines and direction fields provide geometric information
    about the phase portrait of a planar system.
  \item The Hartman--Grobman theorem justifies linearisation near
    hyperbolic equilibria.
  \item The Poincar\'e--Bendixson theorem guarantees the existence of
    limit cycles in bounded, equilibrium-free planar regions.
  \item Dulac's criterion rules out the existence of periodic orbits
    when a suitable multiplier function can be found.
  \item Bifurcations (saddle-node, transcritical, pitchfork, Hopf)
    describe qualitative changes as parameters vary.
  \item Index theory constrains which equilibria a limit cycle can
    enclose.
\end{itemize}


%%%%%%%%%%%%%%%%%%%%%%%%%%%%%%%%%%%%%%%%%%%%%%%%%%%%%%%%%%%%%%%%%%%%%
%%  CHAPTER 10 — Applications                                      %%
%%%%%%%%%%%%%%%%%%%%%%%%%%%%%%%%%%%%%%%%%%%%%%%%%%%%%%%%%%%%%%%%%%%%%
\chapter{Applications --- Mechanics, Circuits, Biology, Demography}
\label{ch:applications}
\index{applications of ODEs}

\section{Mechanical Oscillations}
\label{sec:mechanics}
\index{mechanical oscillations}

\subsection{The Simple Pendulum}
\label{subsec:pendulum}
\index{pendulum}

A rigid pendulum of length $\ell$ and mass $m$ in a gravitational
field $g$ satisfies
\begin{equation}\label{eq:pendulum}
  \ddot{\theta} + \frac{g}{\ell}\sin\theta = 0.
\end{equation}
Writing $\omega=\dot{\theta}$, the phase portrait in the
$(\theta,\omega)$-plane consists of:
\begin{itemize}
  \item \emph{Centres} at $\theta=2k\pi$ (stable equilibria:
    hanging vertically) surrounded by closed orbits (oscillations).
  \item \emph{Saddle points} at $\theta=(2k+1)\pi$ (inverted
    equilibrium), connected by heteroclinic orbits (separatrices).
  \item \emph{Rotational orbits} outside the separatrices
    (the pendulum goes over the top).
\end{itemize}

The energy integral is
\[
  E = \tfrac{1}{2}\omega^{2} - \frac{g}{\ell}\cos\theta.
\]

\begin{center}
\begin{tikzpicture}
  \begin{axis}[
    axis lines=middle,
    xlabel={$\theta$}, ylabel={$\omega$},
    xmin=-7, xmax=7, ymin=-3.5, ymax=3.5,
    width=10cm, height=6.5cm,
    xtick={-6.283,-3.1416,0,3.1416,6.283},
    xticklabels={$-2\pi$,$-\pi$,$0$,$\pi$,$2\pi$},
    ytick=\empty,
    title={Pendulum phase portrait},
    samples=50
  ]
    % Separatrices: E = g/l, take g/l=1 => omega^2/2 = 1+cos(theta)
    \addplot[red, thick, domain=-3.1416:3.1416]
      {sqrt(2*(1+cos(deg(x))))};
    \addplot[red, thick, domain=-3.1416:3.1416]
      {-sqrt(2*(1+cos(deg(x))))};
    % Oscillation orbits (E < g/l): omega^2/2 = 0.5+cos(theta) => E=0.5
    \addplot[blue, domain=-2.1:2.1] {sqrt(2*(0.5+cos(deg(x))))};
    \addplot[blue, domain=-2.1:2.1] {-sqrt(2*(0.5+cos(deg(x))))};
    % Smaller orbit
    \addplot[blue, domain=-1.2:1.2] {sqrt(2*(0.15+cos(deg(x))-1))};
    \addplot[blue, domain=-1.2:1.2] {-sqrt(2*(0.15+cos(deg(x))-1))};
    % Saddle points
    \addplot[only marks, mark=x, mark size=3pt] coordinates
      {(-3.1416,0) (3.1416,0)};
    % Centre
    \addplot[only marks, mark=*, mark size=2pt] coordinates {(0,0)};
  \end{axis}
\end{tikzpicture}
\end{center}

\begin{example}[Small-angle approximation]
\label{ex:pendulum-linear}
For small angles, $\sin\theta\approx\theta$, and
equation~\eqref{eq:pendulum} reduces to the simple harmonic oscillator
$\ddot{\theta}+\omega_{0}^{2}\theta=0$ with
$\omega_{0}=\sqrt{g/\ell}$ and period $T=2\pi/\omega_{0}$.
\end{example}

\subsection{Coupled Oscillators}
\label{subsec:coupled-oscillators}
\index{coupled oscillators}

\begin{example}[Two masses connected by springs]
\label{ex:coupled-springs}
Consider two unit masses connected in series by springs of stiffness
$k$, attached to fixed walls. The equations of motion are
\[
  \ddot{x}_{1} = -2kx_{1}+kx_{2}, \qquad
  \ddot{x}_{2} = kx_{1}-2kx_{2}.
\]
Introducing $u=x_{1}+x_{2}$ and $v=x_{1}-x_{2}$ decouples the
system:
\[
  \ddot{u}=-ku, \qquad \ddot{v}=-3kv.
\]
The normal-mode frequencies are $\omega_{1}=\sqrt{k}$ (in-phase mode)
and $\omega_{2}=\sqrt{3k}$ (out-of-phase mode).
\end{example}

\subsection{Resonance}
\label{subsec:resonance}
\index{resonance}

\begin{example}[Forced, damped harmonic oscillator]
\label{ex:resonance}
The equation
\[
  m\ddot{x} + c\dot{x} + kx = F_{0}\cos(\omega t)
\]
has a particular solution with amplitude
\[
  A(\omega) = \frac{F_{0}}{\sqrt{(k-m\omega^{2})^{2}+c^{2}\omega^{2}}}.
\]
The amplitude is maximal at the resonance frequency
\[
  \omega_{r} = \sqrt{\frac{k}{m}-\frac{c^{2}}{2m^{2}}},
\]
provided $c<\sqrt{2km}$. As $c\to 0$, $\omega_{r}\to\omega_{0}$ and
the peak amplitude $A(\omega_{r})\to\infty$.
\end{example}

\begin{center}
\begin{tikzpicture}
  \begin{axis}[
    axis lines=left,
    xlabel={$\omega$}, ylabel={$A(\omega)$},
    xmin=0, xmax=3, ymin=0, ymax=5,
    width=9cm, height=6cm,
    title={Resonance curves},
    samples=50,
    legend pos=north east
  ]
    % k=1, m=1, F0=1; varying c
    \addplot[blue, thick, domain=0.05:3]
      {1/sqrt((1-x^2)^2+0.04*x^2)};
    \addlegendentry{$c=0.2$}
    \addplot[red, thick, domain=0.05:3]
      {1/sqrt((1-x^2)^2+0.25*x^2)};
    \addlegendentry{$c=0.5$}
    \addplot[green!60!black, thick, domain=0.05:3]
      {1/sqrt((1-x^2)^2+x^2)};
    \addlegendentry{$c=1.0$}
  \end{axis}
\end{tikzpicture}
\end{center}

%----------------------------------------------------------------------
\section{Electrical Circuits}
\label{sec:circuits}
\index{electrical circuits}

By Kirchhoff's laws\index{Kirchhoff's laws}, the voltages and
currents in a circuit satisfy ODEs whose form depends on the circuit
elements.

\subsection{RC Circuit}
\label{subsec:RC}
\index{RC circuit}

\begin{example}[RC circuit with DC source]
\label{ex:RC-circuit}
A resistor $R$ in series with a capacitor $C$ driven by a constant
voltage $V_{0}$ satisfies
\[
  RC\,\dot{V}_{C} + V_{C} = V_{0},
\]
where $V_{C}(t)$ is the voltage across the capacitor. The solution
with $V_{C}(0)=0$ is
\[
  V_{C}(t) = V_{0}\!\left(1-e^{-t/(RC)}\right).
\]
The time constant is $\tau=RC$.
\end{example}

\begin{center}
\begin{tikzpicture}
  % Circuit diagram
  \draw (0,0) -- (1,0);
  \draw (1,0) -- (1,0.3) -- (1.4,0.3) -- (1.4,-0.3) -- (0.6,-0.3) -- (0.6,0.3) -- (1,0.3);
  \node at (1,0.6) {$R$};
  \draw (1.4,0) -- (2.5,0);
  % Capacitor
  \draw (2.5,0.3) -- (2.5,-0.3);
  \draw (2.8,0.3) -- (2.8,-0.3);
  \node at (2.65,0.6) {$C$};
  \draw (2.8,0) -- (3.8,0) -- (3.8,-1.5) -- (0,-1.5) -- (0,0);
  % Voltage source
  \draw (0,-0.5) circle (0.3);
  \node at (-0.55,-0.5) {$V_0$};
  \node at (0,-0.5) {$\sim$};
\end{tikzpicture}
\end{center}

\subsection{RL Circuit}
\label{subsec:RL}
\index{RL circuit}

\begin{example}[RL circuit]
\label{ex:RL-circuit}
A resistor $R$ in series with an inductor $L$ driven by $V_{0}$:
\[
  L\,\dot{I} + RI = V_{0}.
\]
With $I(0)=0$:
\[
  I(t) = \frac{V_{0}}{R}\!\left(1-e^{-Rt/L}\right).
\]
The time constant is $\tau=L/R$.
\end{example}

\subsection{RLC Circuit}
\label{subsec:RLC}
\index{RLC circuit}

\begin{example}[Series RLC circuit]
\label{ex:RLC-circuit}
The charge $q(t)$ on the capacitor in a series RLC circuit satisfies
\begin{equation}\label{eq:RLC}
  L\ddot{q} + R\dot{q} + \frac{q}{C} = V(t).
\end{equation}
This is a second-order linear ODE with constant coefficients. The
characteristic equation is $L\lambda^{2}+R\lambda+1/C=0$, giving
\[
  \lambda = \frac{-R\pm\sqrt{R^{2}-4L/C}}{2L}.
\]
Three regimes arise:
\begin{itemize}
  \item \textbf{Overdamped} ($R^{2}>4L/C$): two negative real roots.
  \item \textbf{Critically damped} ($R^{2}=4L/C$): a repeated real
    root.
  \item \textbf{Underdamped} ($R^{2}<4L/C$): complex conjugate roots
    leading to damped oscillations.
\end{itemize}
\end{example}

\begin{center}
\begin{tikzpicture}
  \begin{axis}[
    axis lines=left,
    xlabel={$t$}, ylabel={$q(t)$},
    xmin=0, xmax=8, ymin=-0.3, ymax=1.3,
    width=9cm, height=5.5cm,
    title={RLC circuit responses},
    samples=50,
    legend pos=north east
  ]
    % Overdamped: lambda = -0.5, -2
    \addplot[blue, thick, domain=0:8]
      {1 - (4/3)*exp(-0.5*x) + (1/3)*exp(-2*x)};
    \addlegendentry{Overdamped}
    % Critically damped: lambda = -1 (repeated)
    \addplot[red, thick, domain=0:8]
      {1 - (1+x)*exp(-x)};
    \addlegendentry{Critically damped}
    % Underdamped: alpha=0.3, omega=1.5
    \addplot[green!60!black, thick, domain=0:8]
      {1 - exp(-0.3*x)*(cos(deg(1.5*x))+0.2*sin(deg(1.5*x)))};
    \addlegendentry{Underdamped}
  \end{axis}
\end{tikzpicture}
\end{center}

%----------------------------------------------------------------------
\section{Population Dynamics}
\label{sec:population-dynamics}
\index{population dynamics}

\subsection{Malthusian Growth}
\label{subsec:malthus}
\index{Malthusian growth}

\begin{example}[Malthus model]
\label{ex:malthus}
The simplest population model is
\[
  \dot{N} = rN, \qquad N(0)=N_{0},
\]
with solution $N(t)=N_{0}e^{rt}$. If $r>0$ the population grows
exponentially; if $r<0$ it decays. The model is unrealistic for large
populations because it ignores resource limitations.
\end{example}

\subsection{Logistic Growth}
\label{subsec:logistic-growth}
\index{logistic equation!solution}

\begin{example}[Logistic equation --- complete solution]
\label{ex:logistic-full}
The logistic equation $\dot{N}=rN(1-N/K)$ is separable. The solution
is
\[
  N(t) = \frac{K}{1+\left(\frac{K}{N_{0}}-1\right)e^{-rt}}.
\]
As $t\to\infty$, $N(t)\to K$ regardless of $N_{0}\in(0,\infty)$.
The inflection point (maximum growth rate) occurs at $N=K/2$.
\end{example}

\begin{center}
\begin{tikzpicture}
  \begin{axis}[
    axis lines=left,
    xlabel={$t$}, ylabel={$N(t)$},
    xmin=0, xmax=10, ymin=0, ymax=1.2,
    width=9cm, height=5.5cm,
    title={Logistic growth ($r=1$, $K=1$)},
    samples=50,
    legend pos=north east
  ]
    \addplot[blue, thick, domain=0:10] {1/(1+(1/0.1-1)*exp(-x))};
    \addlegendentry{$N_0=0.1$}
    \addplot[red, thick, domain=0:10] {1/(1+(1/0.5-1)*exp(-x))};
    \addlegendentry{$N_0=0.5$}
    \addplot[green!60!black, thick, domain=0:10] {1/(1+(1/1.5-1)*exp(-x))};
    \addlegendentry{$N_0=1.5$}
    \addplot[dashed, gray] coordinates {(0,1)(10,1)};
    \node at (axis cs:8.5,1.07) {$K$};
  \end{axis}
\end{tikzpicture}
\end{center}

\subsection{Lotka--Volterra Predator--Prey}
\label{subsec:lotka-volterra-app}
\index{Lotka--Volterra equations!application}

\begin{example}[Oscillatory coexistence]
\label{ex:lv-oscillations}
With parameters $\alpha=\gamma=1$ and $\beta=\delta=0.5$, the
Lotka--Volterra system~\eqref{eq:lotka-volterra} has co-existence
equilibrium $(x^{*},y^{*})=(2,2)$ and conserved quantity
$H=0.5x-\ln x+0.5y-\ln y$. All orbits with $x,y>0$ are closed, and
both populations oscillate periodically with the predator lagging the
prey by a quarter period.
\end{example}

\subsection{The SIR Model}
\label{subsec:sir-app}
\index{SIR model!application}

\begin{example}[Epidemic threshold]
\label{ex:sir-threshold}
Consider a population of $N=1000$ with $\beta=0.3$ and
$\gamma=0.1$. If $I(0)=1$ and $S(0)=999$, then
$\mathcal{R}_{0}=\beta S(0)/\gamma = 2997$, so an epidemic occurs.
The fraction of the population that ultimately remains susceptible,
$s_{\infty}=S_{\infty}/N$, satisfies the transcendental equation
\[
  \ln s_{\infty} = \mathcal{R}_{0}(s_{\infty}-1).
\]
\end{example}

%----------------------------------------------------------------------
\section{Chemical Kinetics}
\label{sec:chemical-kinetics}
\index{chemical kinetics}

\subsection{First-Order Reactions}
\label{subsec:first-order-reaction}
\index{first-order reaction}

\begin{example}[Radioactive decay]
\label{ex:radioactive-decay}
A radioactive substance decays at a rate proportional to the amount
present:
\[
  \dot{A} = -\lambda A, \qquad A(0)=A_{0},
\]
with solution $A(t)=A_{0}e^{-\lambda t}$. The half-life is
$t_{1/2}=\ln 2/\lambda$.
\end{example}

\subsection{Second-Order and Reversible Reactions}
\label{subsec:second-order-reaction}
\index{second-order reaction}

\begin{example}[Second-order reaction $A+B\to C$]
\label{ex:second-order-reaction}
If the initial concentrations are $[A]_{0}=a$ and $[B]_{0}=b$ with
$a\neq b$, the rate law $\dot{x}=k(a-x)(b-x)$ (where $x$ is the
concentration of product) gives
\[
  x(t) = ab\,\frac{e^{k(a-b)t}-1}{ae^{k(a-b)t}-b}.
\]
\end{example}

\subsection{Michaelis--Menten Kinetics}
\label{subsec:michaelis-menten}
\index{Michaelis--Menten kinetics}

\begin{example}[Enzyme kinetics]
\label{ex:michaelis-menten}
In the Michaelis--Menten mechanism, substrate $S$ binds an enzyme $E$
to form a complex $C$, which then yields product $P$:
\[
  E + S \xrightleftharpoons[k_{-1}]{k_{1}} C \xrightarrow{k_{2}} E + P.
\]
Applying the law of mass action with $e_{0}=E+C$ constant yields
\begin{align*}
  \dot{s} &= -k_{1}e_{0}s + (k_{1}s+k_{-1})c, \\
  \dot{c} &= k_{1}e_{0}s - (k_{1}s+k_{-1}+k_{2})c.
\end{align*}
Under the \emph{quasi-steady-state approximation}\index{quasi-steady-state approximation}
($\dot{c}\approx 0$), the reaction velocity is
\[
  v = k_{2}c \approx \frac{V_{\max}s}{K_{m}+s},
\]
where $V_{\max}=k_{2}e_{0}$ and $K_{m}=(k_{-1}+k_{2})/k_{1}$.
\end{example}

\begin{center}
\begin{tikzpicture}
  \begin{axis}[
    axis lines=left,
    xlabel={Substrate $[S]$}, ylabel={Velocity $v$},
    xmin=0, xmax=5, ymin=0, ymax=1.2,
    width=8cm, height=5.5cm,
    title={Michaelis--Menten curve},
    samples=50
  ]
    % Vmax=1, Km=0.5
    \addplot[blue, thick, domain=0:5] {x/(0.5+x)};
    \addplot[dashed, gray] coordinates {(0,1)(5,1)};
    \node at (axis cs:4.2,1.08) {$V_{\max}$};
    \addplot[dashed, gray] coordinates {(0.5,0)(0.5,0.5)};
    \addplot[dashed, gray] coordinates {(0,0.5)(0.5,0.5)};
    \node at (axis cs:0.5,-0.08) {$K_m$};
    \node at (axis cs:-0.35,0.5) {$\frac{V_{\max}}{2}$};
  \end{axis}
\end{tikzpicture}
\end{center}

%----------------------------------------------------------------------
\section{Other Applications}
\label{sec:other-applications}

\subsection{Pursuit Curves}
\label{subsec:pursuit}
\index{pursuit curve}

\begin{example}[Linear pursuit]
\label{ex:pursuit-curve}
A pursuer at position $(x,y)$ always moves directly towards a prey
travelling along the $y$-axis at constant speed $v$. If the pursuer
has speed $w$, the pursuit curve $y=y(x)$ satisfies
\[
  xy'' = \frac{v}{w}\sqrt{1+(y')^{2}}.
\]
Setting $p=y'$ reduces the order:
\[
  xp' = \frac{v}{w}\sqrt{1+p^{2}},
\]
which is separable. When $v=w$, the pursuit curve is a parabola.
When $w>v$, the pursuer eventually catches the prey; when $w<v$, the
pursuer approaches asymptotically but never reaches it.
\end{example}

\subsection{The Catenary}
\label{subsec:catenary}
\index{catenary}

\begin{example}[Hanging chain]
\label{ex:catenary}
A uniform chain of linear density $\rho$ hanging under gravity
satisfies
\[
  y'' = \frac{\rho g}{T_{0}}\sqrt{1+(y')^{2}},
\]
where $T_{0}$ is the horizontal tension. Setting $a=T_{0}/(\rho g)$
and $p=y'$, we get the separable ODE $p'=\frac{1}{a}\sqrt{1+p^{2}}$,
whose solution is $p=\sinh(x/a)$, hence
\[
  y(x) = a\cosh\!\left(\frac{x}{a}\right) + C.
\]
\end{example}

\begin{center}
\begin{tikzpicture}
  \begin{axis}[
    axis lines=middle,
    xlabel={$x$}, ylabel={$y$},
    xmin=-3, xmax=3, ymin=0.5, ymax=4,
    width=8cm, height=5cm,
    title={Catenary $y=\cosh(x)$},
    samples=50
  ]
    \addplot[blue, thick, domain=-2.5:2.5] {cosh(x)};
  \end{axis}
\end{tikzpicture}
\end{center}

%----------------------------------------------------------------------
\section{Modelling Methodology}
\label{sec:modelling-methodology}
\index{mathematical modelling}

The examples in this chapter illustrate a common workflow:

\begin{method}[Mathematical modelling with ODEs]
\label{meth:ode-modelling}
\begin{enumerate}[label=\textbf{Step \arabic*.}]
  \item \textbf{Identify variables and parameters.} Determine the
    dependent variables (state), independent variable (usually time),
    and physical parameters.
  \item \textbf{Write balance/conservation laws.} Apply Newton's law,
    Kirchhoff's laws, mass action, or conservation of mass/energy to
    derive the ODE(s).
  \item \textbf{Non-dimensionalise.} Introduce dimensionless
    variables to reduce the number of parameters and reveal the
    essential scales.
  \item \textbf{Analyse.} Solve explicitly if possible; otherwise,
    find equilibria, linearise, study stability, identify
    bifurcations, and sketch phase portraits.
  \item \textbf{Interpret.} Translate mathematical results back into
    the language of the application.
\end{enumerate}
\end{method}

\begin{example}[Non-dimensionalisation of the pendulum]
\label{ex:nondim-pendulum}
\index{non-dimensionalisation}
Let $\tau=t\sqrt{g/\ell}$. Then
equation~\eqref{eq:pendulum} becomes
\[
  \frac{\dd^{2}\theta}{\dd\tau^{2}} + \sin\theta = 0,
\]
which is parameter-free. All simple pendulums of any length share the
same dimensionless dynamics.
\end{example}

%----------------------------------------------------------------------
\section{Extended Worked Examples}
\label{sec:extended-examples}

\begin{example}[Damped pendulum with torque]
\label{ex:damped-pendulum}
\index{pendulum!damped}
The equation
\[
  \ddot{\theta}+b\dot{\theta}+\sin\theta=\Gamma
\]
models a pendulum with damping $b>0$ and constant applied torque
$\Gamma$. Setting $\omega=\dot{\theta}$:
\[
  \dot{\theta}=\omega,\qquad
  \dot{\omega}=-b\omega-\sin\theta+\Gamma.
\]
\emph{Equilibria:} $\omega=0$ and $\sin\theta=\Gamma$, which exist
only for $\abs{\Gamma}\le 1$. If $\Gamma<1$, there is a stable node
at $\theta_{s}=\arcsin\Gamma$ and a saddle at
$\theta_{u}=\pi-\arcsin\Gamma$. At $\Gamma=1$, these coalesce (a
saddle-node bifurcation), and for $\Gamma>1$ no equilibrium exists;
the pendulum rotates continuously.
\end{example}

\begin{example}[Coupled RLC and mechanical analogy]
\label{ex:rlc-mechanical}
\index{RLC circuit!mechanical analogy}
The series RLC equation~\eqref{eq:RLC} is mathematically identical to
$m\ddot{x}+c\dot{x}+kx=F(t)$ under the correspondence
\[
  q\leftrightarrow x,\quad L\leftrightarrow m,\quad
  R\leftrightarrow c,\quad 1/C\leftrightarrow k,\quad
  V(t)\leftrightarrow F(t).
\]
This analogy allows one to transfer intuition freely between
mechanical and electrical systems.
\end{example}

\begin{example}[Lotka--Volterra with harvesting]
\label{ex:lv-harvesting}
\index{Lotka--Volterra equations!harvesting}
Adding a constant harvesting rate $h\ge 0$ to the prey equation:
\[
  \dot{x}=\alpha x-\beta xy - h,\qquad
  \dot{y}=\delta xy-\gamma y.
\]
The co-existence equilibrium shifts to
$y^{*}=(\alpha-h\delta/\gamma)/\beta$ (for sufficiently small $h$).
When $h$ exceeds a critical threshold, the prey equilibrium disappears
via a saddle-node bifurcation, leading to extinction of both species.
\end{example}

\begin{example}[Logistic growth with periodic harvesting]
\label{ex:logistic-periodic}
\index{logistic equation!periodic harvesting}
Consider
\[
  \dot{N} = rN\!\left(1-\frac{N}{K}\right) - h(1+\varepsilon\sin\omega t),
\]
where $h>0$ is the mean harvest rate and $\varepsilon\in[0,1)$
modulates seasonality. This is a non-autonomous ODE. For small
$\varepsilon$, perturbation methods show that the average population
oscillates near the equilibrium of the time-averaged equation,
$N^{*}=K(1-h/(rK/4))/2$ (assuming the average equation has two
positive roots).
\end{example}

%----------------------------------------------------------------------
\section{Exercises}
\label{sec:ch10-exercises}

\begin{exercise}
\label{exer:pendulum-energy}
For the simple pendulum~\eqref{eq:pendulum}, show that the energy
$E=\frac{1}{2}\dot{\theta}^{2}-(g/\ell)\cos\theta$ is constant along
solutions. Classify the orbits in the phase plane according to the
value of $E$.
\end{exercise}

\begin{exercise}
\label{exer:rlc-forced}
Find the steady-state (particular) solution of the RLC
equation~\eqref{eq:RLC} with $V(t)=V_{0}\cos(\omega t)$. For what
$\omega$ is the amplitude of the charge $q$ maximised?
\end{exercise}

\begin{exercise}
\label{exer:logistic-solve}
Solve the logistic equation $\dot{N}=2N(1-N/500)$ with $N(0)=50$ by
separation of variables. Find the time at which $N=250$.
\end{exercise}

\begin{exercise}
\label{exer:lv-linearise}
Consider the Lotka--Volterra system with $\alpha=2$, $\beta=1$,
$\gamma=1$, $\delta=1$. Find the co-existence equilibrium and
linearise. Determine the period of oscillation predicted by the
linearisation.
\end{exercise}

\begin{exercise}
\label{exer:sir-peak}
In the SIR model, show that the peak number of infected individuals is
\[
  I_{\max} = I_{0}+S_{0}-\frac{\gamma}{\beta}
    -\frac{\gamma}{\beta}\ln\!\left(\frac{\beta S_{0}}{\gamma}\right).
\]
\end{exercise}

\begin{exercise}
\label{exer:michaelis-menten-solve}
Derive the Michaelis--Menten approximation
$v=V_{\max}s/(K_{m}+s)$ from the full enzyme system by applying the
quasi-steady-state assumption $\dot{c}\approx 0$.
\end{exercise}

\begin{exercise}
\label{exer:pursuit-solve}
Solve the pursuit curve ODE $xp'=(v/w)\sqrt{1+p^{2}}$ with
$p(x_{0})=0$ in the case $v/w=1$. Show that the solution is a
parabola.
\end{exercise}

\begin{exercise}
\label{exer:catenary-BC}
A chain of length $L$ hangs between two supports at equal height
separated by a horizontal distance $d<L$. Show that the parameter
$a$ in $y=a\cosh(x/a)$ satisfies $\sinh(d/(2a))=L/(2a)$ (after
centering the catenary). Prove that this equation has a unique
solution $a>0$ when $d<L$.
\end{exercise}

\begin{exercise}
\label{exer:nondim-exercise}
Non-dimensionalise the damped pendulum equation
$\ddot{\theta}+b\dot{\theta}+\omega_{0}^{2}\sin\theta=0$ using the
time scale $1/\omega_{0}$. Identify the single dimensionless
parameter and discuss the phase portrait for small and large values
of this parameter.
\end{exercise}

\begin{exercise}
\label{exer:rc-ac}
An RC circuit is driven by an alternating source
$V(t)=V_{0}\sin(\omega t)$. Find the steady-state solution and
compute the phase lag between the driving voltage and the capacitor
voltage. For what value of $\omega$ is the lag $\pi/4$?
\end{exercise}

\begin{exercise}[Modelling project]
\label{exer:modelling-project}
Choose one of the following scenarios and set up, analyse, and
interpret a model using ODEs:
\begin{enumerate}[label=(\alph*)]
  \item The spread of a rumour in a population of $N$ people, where
    the rate of spread is proportional to the product of those who
    know it and those who do not.
  \item A parachutist falling under gravity with air resistance
    proportional to $v^{2}$, opening a parachute at a specified
    altitude.
  \item The temperature of a building with heating, insulation, and
    an outdoor temperature that varies sinusoidally over 24 hours.
\end{enumerate}
\end{exercise}

%----------------------------------------------------------------------
\section*{Chapter Summary}
\addcontentsline{toc}{section}{Chapter Summary}

\begin{itemize}
  \item The simple pendulum, coupled oscillators, and resonance
    illustrate how mechanical systems naturally give rise to
    second-order ODEs.
  \item RC, RL, and RLC circuits provide first- and second-order
    linear ODEs; the mechanical--electrical analogy connects the two
    domains.
  \item Population models (Malthus, logistic, Lotka--Volterra, SIR)
    demonstrate both linear and nonlinear dynamics, including
    conservation laws and threshold phenomena.
  \item Michaelis--Menten kinetics illustrates quasi-steady-state
    reduction of a system.
  \item Pursuit curves and the catenary give elegant geometric
    applications of separable and second-order ODEs.
  \item A systematic modelling methodology ---~identifying variables,
    writing balance laws, non-dimensionalising, analysing, and
    interpreting ---~ties together all applications.
\end{itemize}


%%%%%%%%%%%%%%%%%%%%%%%%%%%%%%%%%%%%%%%%%%%%%%%%%%%%%%%%%%%%%%%%%%%%%
%%  END MATTER                                                     %%
%%%%%%%%%%%%%%%%%%%%%%%%%%%%%%%%%%%%%%%%%%%%%%%%%%%%%%%%%%%%%%%%%%%%%

\chapter*{Quick Reference Card}
\addcontentsline{toc}{chapter}{Quick Reference Card}
\label{ch:quick-reference}
\index{quick reference}

\section*{Standard ODE Types and Solution Methods}

\renewcommand{\arraystretch}{1.4}

\begin{center}
\begin{tabular}{@{}p{4.5cm}p{4cm}p{5.5cm}@{}}
\toprule
\textbf{ODE Type} & \textbf{Form} & \textbf{Method} \\
\midrule
Separable &
  $y'=f(x)g(y)$ &
  $\displaystyle\int\frac{\dd y}{g(y)}=\int f(x)\,\dd x$ \\[6pt]
Linear first-order &
  $y'+P(x)y=Q(x)$ &
  Integrating factor $\mu=e^{\int P\,\dd x}$ \\[6pt]
Exact &
  $M\,\dd x+N\,\dd y=0$, $M_{y}=N_{x}$ &
  Find $F$: $F_{x}=M$, $F_{y}=N$ \\[6pt]
Bernoulli &
  $y'+Py=Qy^{n}$ &
  Substitution $v=y^{1-n}$ \\[6pt]
Homogeneous &
  $y'=\phi(y/x)$ &
  Substitution $v=y/x$ \\[6pt]
Riccati &
  $y'=P+Qy+Ry^{2}$ &
  One particular solution $\to$ Bernoulli \\[6pt]
Second-order const.\ coeff. &
  $ay''+by'+cy=f(x)$ &
  Characteristic equation + undetermined coefficients or variation of parameters \\[6pt]
Euler--Cauchy &
  $x^{2}y''+bxy'+cy=0$ &
  Try $y=x^{r}$ \\[6pt]
System $\dot{\mathbf{x}}=A\mathbf{x}$ &
  Constant matrix $A$ &
  $\mathbf{x}(t)=e^{At}\mathbf{x}_{0}$ via eigenvalues \\
\bottomrule
\end{tabular}
\end{center}

\section*{Key Theorems}

\begin{center}
\begin{tabular}{@{}p{4.5cm}p{9.5cm}@{}}
\toprule
\textbf{Theorem} & \textbf{Statement (informal)} \\
\midrule
Picard--Lindel\"of &
  If $f$ is Lipschitz in $y$, the IVP $y'=f(x,y)$, $y(x_{0})=y_{0}$
  has a unique local solution. \\[4pt]
Peano &
  If $f$ is continuous, the IVP has at least one solution (uniqueness
  may fail). \\[4pt]
Gr\"onwall's inequality &
  Controls the growth of solutions; key tool for continuous dependence
  and uniqueness proofs. \\[4pt]
Hartman--Grobman &
  Near a hyperbolic equilibrium, the nonlinear flow is topologically
  conjugate to the linearised flow. \\[4pt]
Poincar\'e--Bendixson &
  A bounded orbit of a planar system that avoids equilibria must
  approach a periodic orbit. \\[4pt]
Liouville's formula &
  $W(t)=W(t_{0})\exp\!\bigl(\int_{t_0}^{t}\tr A(s)\,\dd s\bigr)$. \\
\bottomrule
\end{tabular}
\end{center}

\section*{Table of Laplace Transforms}
\index{Laplace transform!table}

\begin{center}
\begin{tabular}{@{}lcl@{}}
\toprule
$f(t)$ \quad ($t\ge 0$) & \qquad & $\mathcal{L}\{f\}(s)=F(s)$ \\
\midrule
$1$ && $\dfrac{1}{s}$, \; $s>0$ \\[8pt]
$t^{n}$ \; ($n\in\N$) && $\dfrac{n!}{s^{n+1}}$ \\[8pt]
$e^{at}$ && $\dfrac{1}{s-a}$, \; $s>a$ \\[8pt]
$\sin(bt)$ && $\dfrac{b}{s^{2}+b^{2}}$ \\[8pt]
$\cos(bt)$ && $\dfrac{s}{s^{2}+b^{2}}$ \\[8pt]
$e^{at}\sin(bt)$ && $\dfrac{b}{(s-a)^{2}+b^{2}}$ \\[8pt]
$e^{at}\cos(bt)$ && $\dfrac{s-a}{(s-a)^{2}+b^{2}}$ \\[8pt]
$t^{n}e^{at}$ && $\dfrac{n!}{(s-a)^{n+1}}$ \\[8pt]
$u_{c}(t)=\begin{cases}0&t<c\\1&t\ge c\end{cases}$ &&
  $\dfrac{e^{-cs}}{s}$ \\[8pt]
$u_{c}(t)\,f(t-c)$ && $e^{-cs}F(s)$ \\[8pt]
$\delta(t-c)$ && $e^{-cs}$ \\[8pt]
$f'(t)$ && $sF(s)-f(0)$ \\[8pt]
$f''(t)$ && $s^{2}F(s)-sf(0)-f'(0)$ \\[8pt]
$(f*g)(t)$ && $F(s)\cdot G(s)$ \\
\bottomrule
\end{tabular}
\end{center}

\section*{Stability Summary for $\dot{\mathbf{x}}=A\mathbf{x}$ in $\R^{2}$}
\index{stability!summary}

\begin{center}
\begin{tabular}{@{}lll@{}}
\toprule
\textbf{Eigenvalues} & \textbf{Type} & \textbf{Stability} \\
\midrule
$\lambda_{1}<\lambda_{2}<0$ & Stable node & Asymp.\ stable \\
$0<\lambda_{1}<\lambda_{2}$ & Unstable node & Unstable \\
$\lambda_{1}<0<\lambda_{2}$ & Saddle & Unstable \\
$\alpha\pm i\beta$, $\alpha<0$ & Stable spiral & Asymp.\ stable \\
$\alpha\pm i\beta$, $\alpha>0$ & Unstable spiral & Unstable \\
$\pm i\beta$ & Centre & Stable (not asymp.) \\
$\lambda<0$ repeated, $2$ eigenvectors & Stable star node &
  Asymp.\ stable \\
$\lambda<0$ repeated, $1$ eigenvector & Stable improper node &
  Asymp.\ stable \\
\bottomrule
\end{tabular}
\end{center}

%----------------------------------------------------------------------
% Bibliography
%----------------------------------------------------------------------
\begin{thebibliography}{99}

\bibitem{tenenbaum}
M.~Tenenbaum and H.~Pollard,
\emph{Ordinary Differential Equations},
Dover, 1985.
\index{Tenenbaum}

\bibitem{arnold}
V.\,I.~Arnold,
\emph{Ordinary Differential Equations},
Springer, 1992.
\index{Arnold}

\bibitem{chicone}
C.~Chicone,
\emph{Ordinary Differential Equations with Applications},
Springer, 2nd edition, 2006.
\index{Chicone}

\bibitem{hirsch}
M.\,W.~Hirsch, S.~Smale, and R.\,L.~Devaney,
\emph{Differential Equations, Dynamical Systems, and an Introduction
to Chaos},
Academic Press, 3rd edition, 2013.
\index{Hirsch}\index{Smale}\index{Devaney}

\bibitem{verhulst}
F.~Verhulst,
\emph{Nonlinear Differential Equations and Dynamical Systems},
Springer, 2nd edition, 1996.
\index{Verhulst}

\bibitem{braun}
M.~Braun,
\emph{Differential Equations and Their Applications},
Springer, 4th edition, 1993.
\index{Braun}

\bibitem{perko}
L.~Perko,
\emph{Differential Equations and Dynamical Systems},
Springer, 3rd edition, 2001.
\index{Perko}

\end{thebibliography}

\printindex
\end{document}
